\documentclass[a4paper,11pt]{article}
\usepackage[top=23mm,bottom=23mm,left=25mm,right=25mm]{geometry}
\usepackage{lmodern}
\usepackage{mathtools,amssymb,amsthm,mathrsfs}
\usepackage{booktabs}
\usepackage{xcolor}
\usepackage[unicode]{hyperref}

\newcommand{\paperlanguage}[2]{#2}
\newcommand{\upstream}{https://github.com/lsernau/grh-lean/blob/72fead102f663d5c7a407964a13dc0cf9d977207}
\newcommand{\leansrc}[2]{\href{\upstream/grh/src/#1.lean}{#2}}
\newcommand{\leandecl}[3]{\href{\upstream/grh/src/#1.lean\#L#2}{\texttt{#3}}}

\newcounter{leanentry}
\renewcommand{\theleanentry}{L\arabic{leanentry}}
\newcommand{\leanref}[1]{\unskip\nobreak\hspace{0.2em}%
  \textsuperscript{\normalfont\hyperref[lean:#1]{[\ref*{lean:#1}]}}}
\newcommand{\leanentry}[4]{%
  \item[]\refstepcounter{leanentry}\label{lean:#1}%
  \textbf{[\theleanentry]}\enspace\paperlanguage{#2}{#3}\par\nobreak
  {\small #4\par}}

\definecolor{supplementblue}{HTML}{174A7A}
\newenvironment{editorialnote}
  {\begin{quote}\color{supplementblue}\small\textbf{Note}\quad\ignorespaces}
  {\end{quote}}
\newcommand{\RePart}{\operatorname{Re}}
\newcommand{\ImPart}{\operatorname{Im}}
\newcommand{\res}{\operatorname*{Res}}

\theoremstyle{plain}
\newtheorem{theorem}{Theorem}[section]
\newtheorem{proposition}[theorem]{Proposition}
\newtheorem{lemma}[theorem]{Lemma}
\newtheorem{corollary}[theorem]{Corollary}
\theoremstyle{definition}
\newtheorem{definition}[theorem]{Definition}
\numberwithin{equation}{section}
\allowdisplaybreaks[1]
\setlength{\emergencystretch}{1.5em}

\hypersetup{
  pdftitle={The generalized Riemann hypothesis for Dirichlet L-functions},
  pdfsubject={An exposition of the Lean proof of Alpöge and Claude}
}
\title{The generalized Riemann hypothesis for Dirichlet \(L\)-functions\\[0.45em]
  \large An exposition of the formal proof via theta lifts}
\author{}
\date{\small Expository draft \quad 26 November 2026}

\begin{document}
\maketitle

\begin{abstract}
We give an account of the formal proof by Levent Alpöge and Claude
of the generalized Riemann hypothesis for Dirichlet \(L\)-functions,
following the readable edition of \emph{grh-lean}. The proof obtains a
one-sided bound for the remainder in a sum over prime powers in an
arithmetic progression with a quadratic character weight. Iterated
theta lifts between unitary groups furnish a comparison between the
squared covolume of the lifted forms and the index required for
integrality. Divisibility of theta derivatives and a second moment for
Poincaré series give the upper bound. Twists chosen separately at the
finite places, ratios of Gram determinants, and the product formula
give the lower bound. To pass from the one-sided estimate to
nonvanishing in a half-plane, a zero-density estimate for quadratic
twists prevents cancellation of residues in the difference of
logarithmic derivatives, and Landau's lemma then excludes the poles.
\end{abstract}

\section{The main theorem}

\begin{theorem}[Alpöge's formal proof]\label{thm:grh}
Let \(q\ge1\) be an integer and let \(\chi\) be a Dirichlet character
modulo \(q\).
\begin{enumerate}
\item If \(\RePart s>1/2\), then \(L(s,\chi)\ne0\), with the exception
      of \(s=1\) when \(\chi\) is principal.
\item If \(0<\RePart s<1\) and \(L(s,\chi)=0\), then \(\RePart s=1/2\).
\item All zeros of the Riemann zeta function other than the trivial
      zeros \(-2,-4,\ldots\) lie on \(\RePart s=1/2\).
\end{enumerate}
\end{theorem}

\noindent
Here \(L(s,\chi)\) is the meromorphic continuation of
\(\sum_{n\ge1}\chi(n)n^{-s}\) from \(\RePart s>1\).
The character \(\chi\) may be imprimitive.
The \(L\)-function of the character modulo \(1\) is \(\zeta(s)\).
Of the five formal statements, two express (3).
The other three express (1), (2), and the specialization of (2)
to primitive characters.
See the \leansrc{Statements}{formal statements} and the
\leansrc{Main}{main declarations}.

We first deduce (1) from a one-sided bound for a weighted prime-power sum
(Section~\ref{sec:analytic}). We then place weights at the infinite places
of an auxiliary number field to express the remainder in terms of
divisibility of theta derivatives (Section~\ref{sec:data}).
Section~\ref{sec:lattices} combines the infinite and finite places for
a single Gram determinant. Sections~\ref{sec:upper-bound}--\ref{sec:determinant-estimates}
establish the comparison's upper bound, dimension estimate, and bound
for the determinant ratios under twisting. Section~\ref{sec:conclusion}
deduces the theorem.

The proof follows Alpöge's announcement \cite{alpoege} and the
\texttt{grh/} project, which the source repository calls its readable
edition \cite{grh-source}. References beginning with L point to the
corresponding declarations in the appendix. The proofs of the estimates
for theta valuations, the Poincaré second moment, and the choice of
twists and projection denominators appear in their respective sections.
For auxiliary results on the Weil representation and local orbit sums,
we give the forms used and cite the corresponding sources. Blue notes are editorial
explanations. The overview that follows describes the main ideas and
points to related accounts.

\section*{Overview of the proof}
\addcontentsline{toc}{section}{Overview of the proof}
\label{sec:overview}

To estimate a prime-power sum of length \(x\), the proof composes
\(m\asymp\sqrt{x}\) theta lifts. The prime-sum remainder then appears with
a factor of \(m\) in a bound for a Gram determinant. Upper and lower bounds
for the same determinant cancel the common main term and leave a one-sided
bound for the remainder.

\subsection*{Why iterate the lift?}

In Stepanov's method for curves over finite fields, an auxiliary function
vanishes to high order at many points, and the sum of the orders of
vanishing is compared with its degree. In his announcement, Alpöge
describes iterated theta lifts over number fields as performing the role
of this multiplicity. In his account of the original proof, Zhang
also emphasizes that the leading finite-place valuation gains grow with
the number \(m\) of stages. \cite[pp.~1, 4]{alpoege} \cite[Section~4]{zhang}

In the readable edition, one estimates divisibility at finite places
after differentiating a theta kernel. For example, if \(p\) is inert
in \(K=\mathbb Q(\sqrt{-d})\), then \(1-\chi_K(p^j)\) is \(2\) for odd
\(j\) and \(0\) for even \(j\). The contribution of \(p^j\) to the local
valuation includes the cutoff weight. Summing the local gains with weight
\(2\log p\) produces the required weighted prime-power sums. A seesaw
identity connects the chain whose divisibility can be calculated to the
chain whose Gram determinant is estimated.
See Section~\ref{sec:upper-bound}.

\subsection*{The main term cancels in the determinant comparison}

Let \(c\) bound the logarithm of the squared covolume of \(n\) lifted
forms, and let \(\delta\) bound the logarithmic index required to achieve
integrality at the finite places. The index occurs as \(2\delta\)
because one change-of-coordinates matrix enters for each variable of
the Hermitian pairing. Let \(h\) be the target dimension, \(H_0\) the
common main term, and \(\mathcal R\) the total prime-sum remainder.
Write \(f\) for the number of infinite places and \(b\) for the weight
at the selected places, and set \(a=xb\) and \(E_j=fa(\log b)^j\).
The bounds take the form
\[
 h(H_0-3E_4)
 \le c+2\delta
 \le n(H_0-m\mathcal R+E_4)
     +20h\log(1+2b^2)E_3 .
\]

For the upper bound, divisibility enters the normalization of the
Poincaré map, after which one estimates the squared norms at infinity.
A second moment for translates of the map yields
\(n\ge(1-4/b^2)h\) linearly independent images of short forms.
For the lower bound, one selects a twist at each finite place and
compares the untwisted Gram determinant with one whose integrality
is controlled at that place. The height of the determinant ratio uses
Hecke \(L\)-values and bounds for denominators of projections to the
spaces of lifts. The product formula then gives the same \(H_0\).
See Sections~\ref{sec:lattices}--\ref{sec:determinant-estimates}.

Using \(n\le h\) and \(H_0\ge0\) cancels \(H_0\), and \(n\ge h/2\)
gives \(m\mathcal R\le fa(\log b)^5\). An identity using two weights
at the infinite places extracts the remainder \(R_{q_0}(x)\) in
the residue class \(1\) modulo \(q_0\). The fixed-modulus prime number
theorem controls the terms at the larger cutoff \(x_1\). The leading
term of the remaining bound is \(\varphi(q_0)(\log b)^6/m\).
Since \(m\ge\sqrt{x}\) and \(\log b\) can be chosen below an
appropriately small positive power of \(x\), this gives the one-sided
bound \(R_{q_0}(x)\le x^{-1/2+\varepsilon}\).
See Sections~\ref{sec:data} and \ref{sec:conclusion}.

\begin{editorialnote}
The working guide to the zeta case published by Rubinstein explains
the square-root scale in the original, longer proof. After
normalization by the coefficient of the prime sum, errors on scales
\(x/m\) and \(m\) balance when \(m\asymp\sqrt{x}\). The comparison and
the bound above are those of the readable edition.
\cite[Section~I.8]{rubinstein}
\end{editorialnote}

\subsection*{Why a one-sided bound excludes zeros}

The remainder \(R_{q_0}(x)\) comes from summing over
\(n\equiv1\pmod{q_0}\) with weight
\(\Lambda(n)\bigl(1-(n/d)\bigr)(1/n-1/x)\) and subtracting
\(\log x/\varphi(q_0)+c_{q_0}\). Zeros of the original
\(L\)-functions and zeros of their quadratic twists enter the
transform of this remainder with opposite residue signs.
The choice of \(d\), made with the height of the candidate zero
in mind, excludes twisted zeros in the relevant region. The
remaining residues on the real axis are nonnegative. The
one-sided bound supplies a nonnegative density for Landau's lemma,
which rules out the pole associated to the candidate zero.
The functional equation carries the result to the other side
of the critical line. See Section~\ref{sec:analytic}.

For local divisibility and the slope comparison in the original proof,
see the accounts by Slijepčević and Zhang. The analytic reduction described
by Slijepčević uses bounds on intervals \([Y,Y^{3/2}]\) with arbitrarily
large \(Y\), as in the original proof.
\cite[Section~1.3]{slijepcevic} \cite[Sections~1 and~4]{zhang}

\section{Prime-power sums and zeros}\label{sec:analytic}

\subsection{The Laplace transform of the remainder}

Let \(d\) be an odd prime and \(q\) a positive integer coprime to \(d\).
Write \(\chi_d(n)=(n/d)\) for the Legendre symbol, \(\Lambda\) for the von Mangoldt function,
and \(\varphi\) for Euler's totient function. Put
\[
  \lambda_q(n)=\mathbf 1_{n\equiv1\bmod q}\Lambda(n)(1-\chi_d(n)).
\]
For the moment, \(d\) is fixed and is suppressed from the notation \(\lambda_q,R_q,G_q\).
If \(\psi\) is principal, let \(L^\circ(s,\psi)=(s-1)L(s,\psi)\), with the value at
\(s=1\) given by holomorphic continuation. For nonprincipal \(\psi\), put
\(L^\circ(s,\psi)=L(s,\psi)\). Both definitions give entire functions.
We regard \(\psi\chi_d\) as a character modulo \(qd\). Since \(d\) is an odd prime
and \((d,q)=1\), \(\psi\chi_d\) is nonprincipal. Set
\begin{align}
  G_q(s)
   &=\frac1{\varphi(q)}
     \left[
       \sum_{\psi\bmod q}
       \left(
        \frac{(L^\circ)'(s,\psi\chi_d)}{L^\circ(s,\psi\chi_d)}
        -\frac{(L^\circ)'(s,\psi)}{L^\circ(s,\psi)}
       \right)-1
     \right], \label{eq:log-derivative}\\
  R_q(y)
   &= \sum_{n\le y}\lambda_q(n)\left(\frac1n-\frac1y\right)
      -\frac{\log y}{\varphi(q)}-c_q,\qquad c_q=\RePart G_q(1).
  \label{eq:prime-remainder}
\end{align}
Logarithmic derivatives are taken with respect to \(s\). For a principal character,
\(L^\circ\) is nonzero at \(s=1\) after removal of the pole. For a nonprincipal
character, \(L(1,\psi)\ne0\). Thus \(G_q\) is holomorphic at \(1\).

By character orthogonality, \(\lambda_q(n)\) is the average of
\(\Lambda(n)(\psi(n)-(\psi\chi_d)(n))\) over all characters modulo \(q\).
Logarithmic differentiation of the Euler products for \(\RePart s>1\), with the term
\(1/(s-1)\) from removing the pole of the principal \(L\)-function accounted for, gives
\begin{equation}\label{eq:lambda-dirichlet}
  F_q(s):=\sum_{n\ge1}\frac{\lambda_q(n)}{n^s}
        =\frac{s}{\varphi(q)(s-1)}+G_q(s).
\end{equation}
The right-hand side is real for real \(s>1\), so \(G_q(1)=c_q\).

\begin{lemma}\label{lem:prime-transform}
For \(\RePart s>2\) one has
\begin{equation}\label{eq:prime-transform}
  \mathcal T_q(s):=\int_0^\infty R_q(e^v)e^{(1-s)v}\,dv
     =\frac{G_q(s)/s-G_q(1)}{s-1}.
\end{equation}
The meromorphic continuation given by the right-hand side is holomorphic at \(s=1\).
At a point \(\rho\ne1\) with \(\RePart\rho>0\), it has at most a simple pole, with residue
\[
  \res_{s=\rho}\mathcal T_q(s)=
    \frac{\nu(\rho)-\nu_d(\rho)}{\varphi(q)\rho(1-\rho)}.
\]
Here \(\nu\) and \(\nu_d\) are the sums over all characters modulo \(q\) of the
orders of vanishing at \(\rho\) of \(L^\circ(s,\psi)\) and \(L^\circ(s,\psi\chi_d)\),
respectively.
\end{lemma}

\begin{proof}
Put \(A_r(y)=y^{-r}\sum_{n\le y}\lambda_q(n)n^{r-1}\).
For \(r=0,1\), integrating the \(n\)-th summand over \(v\ge\log n\) gives
\[
 \int_0^\infty A_r(e^v)e^{(1-s)v}\,dv=\frac{F_q(s)}{s+r-1}.
\]
The termwise integration is absolutely convergent for \(\RePart s>2\). Consequently
\[
 \mathcal T_q(s)
 = F_q(s)\left(\frac1{s-1}-\frac1s\right)
    -\frac1{\varphi(q)(s-1)^2}-\frac{G_q(1)}{s-1}.
\]
Substitution of \eqref{eq:lambda-dirichlet} gives \eqref{eq:prime-transform}.
The numerator vanishes at \(s=1\). At \(s=\rho\ne1\), apply to
\eqref{eq:log-derivative} the fact that the residue of the logarithmic derivative
of an entire function equals its order of vanishing.
\leanlinks{\leanref{integral_primeRemainder}}
\end{proof}

\subsection{Selecting the quadratic twist}

\begin{lemma}\label{lem:select-quadratic}
Let \(q\) be a positive multiple of \(8\), with \(d_0\equiv7\pmod8\) and \((d_0,q)=1\).
For every \(\beta>1/2\) and \(H\ge0\) there is a prime \(d\equiv d_0\pmod q\)
for which \(L^\circ(s,\psi\chi_d)\) is nonzero throughout
\(\RePart s\ge\beta,\ |\ImPart s|\le H\), for every character \(\psi\) modulo \(q\).
\end{lemma}

\begin{proof}
For \(\beta\ge1\), the case \(\beta=3/4\) suffices.
Assume \(1/2<\beta<1\) and write \(\beta=1/2+9\eta\), with \(\eta>0\).
The formal proof gives the bound \(O_{q,\eta,H}(X^{1-\eta/2})\) for the number
of primes \(d\le X\) for which at least one twisted function has a zero
in the specified region, by the following argument.

Let \(\mu\) be the Möbius function. For \(j\in\{0,1\}\) put
\(w_j(u)=u^j e^{-\pi u^2}\) and write
\(\Gamma_j(s)=\int_0^\infty w_j(u)u^{s-1}\,du\) for its Mellin transform.
For a sequence of coefficients \(c=(c(n))_{n\ge1}\), set
\[
 \theta_c(u)=u^{1/2}\sum_{n\ge1}c(n)w_j(nu),\qquad
 J(c,T)=\int_{1/T}^\infty
        \bigl(|\theta_c(u)|+|\theta_{\mu c}(u)|\bigr)\,\frac{du}{u}.
\]
Let \(\chi\) be a nonprincipal primitive character of conductor \(Q\le T^2\)
and parity \(j\). The Mellin transform of \(\theta_\chi\) at \(s-1/2\) equals
\(\Gamma_j(s)L(s,\chi)\). For \(\RePart s>1\), the corresponding integral for
\(\theta_{\mu\chi}\) gives \(\Gamma_j(s)/L(s,\chi)\). Let \(M_T(s)\) denote
that integral restricted to \(u>1/T\). On \(\RePart s=1/2\) we have
\(|M_T(s)|\le J(\chi,T)\).
The functional equation for the theta function gives
\(\int_0^\infty |\theta_\chi(u)|\,du/u\) as twice
\(\int_{Q^{-1/2}}^\infty|\theta_\chi(u)|\,du/u\).
Since \(Q^{-1/2}\ge T^{-1}\), it follows that
\(|\Gamma_j(s)L(s,\chi)|\le2J(\chi,T)\) on this line.

Suppose \(J(\chi,T)\le T^{2\eta}\) and put
\(\Psi(s)=\Gamma_j(s)L(s,\chi)M_T(s)-\Gamma_j(s)^2\).
On the left boundary \(\RePart s=1/2\) we have \(|\Psi(s)|\ll T^{4\eta}\).
On the right boundary \(\RePart s=2\), the bound \(3e^{-u}\) for the
Mellin integrand of a coefficient sequence bounded by \(1\) in modulus gives
\[
 |\Gamma_j(s)L(s,\chi)|\le3,\qquad
 \left|\frac{\Gamma_j(s)}{L(s,\chi)}-M_T(s)\right|\le 3/T.
\]
Thus \(|\Psi(s)|\le9/T\) there. The Mellin representation of \(\Gamma_j L\) and
the truncated integral defining \(M_T\) give the holomorphy and boundedness
required by the three-lines theorem on the strip. For
\(\RePart s\in[1/2+9\eta,1]\), the interpolation weight of the right boundary is
\(t=(\RePart s-1/2)/(3/2)\ge6\eta\).
Hence \(|\Psi(s)|\ll T^{4\eta(1-t)-t}\), and the exponent is at most \(-2\eta\).
The Euler product rules out zeros in \(\RePart s>1\).
If \(L(s,\chi)=0\), then \(\Psi(s)=-\Gamma_j(s)^2\).
The modulus of \(\Gamma_j\) has a positive lower bound on
\(1/2\le\RePart s\le1,\ |\ImPart s|\le H\). Therefore a zero in the required
region forces \(J(\chi,T)>T^{2\eta}\) for all sufficiently large \(T\).
\leanlinks{\leanref{eventually_rpow_lt_thetaTailIntegral}}

Next let \(n^\circ\) be the odd part of \(n\), and for \(|\alpha_n|\le1\) put
\(c_v(n)=\alpha_n(-v/n^\circ)\), using the Jacobi symbol.
When \(u\ge T^{-1}\), expansion of \(\sum_{1\le v\le T^2}|\theta_{c_v}(u)|^2\)
gives the character sum
\[
  \sum_{1\le v\le T^2}\left(\frac{-v}{(nn')^\circ}\right)
\]
for the pair \(n,n'\). If the odd part of \(nn'\) is a square, the absolute
value of the sum is at most \(T^2\). With \(A=1/2+\eta\) in
\(|w_j(x)|\ll_A x^{-A}\), the contribution of the pair is bounded by a constant
times \(T^2u^{-2\eta}(nn')^{-1/2-\eta}\).
For such a pair there is a \(k\in\{1,2\}\) for which \(knn'\) is a square, so
we may write \(kn=ab^2,\ n'=ac^2\). The sum \(\sum(nn')^{-1/2-\eta}\) over these
pairs is therefore bounded by a constant times the convergent triple sum
\(\sum_{a,b,c\ge1}(abc)^{-1-2\eta}\).

If the odd part is not a square, the Pólya--Vinogradov inequality for the Jacobi
character gives \(O_\eta((nn')^{1/2+\eta/2})\) for the character sum, after
absorbing the logarithmic factor. Taking \(A=3/2+\eta\) and using \(u^{-2}\le T^2\)
bounds the contribution of a pair by a constant times
\(T^2u^{-2\eta}(nn')^{-1-\eta/2}\). This double sum also converges. Combining the
two bounds gives
\[
 \sum_{v\le T^2}|\theta_{c_v}(u)|^2\ll_\eta T^2u^{-2\eta}.
\]
Cauchy--Schwarz followed by integration against \(du/u\) yields
\(\sum_{v\le T^2}J(c_v,T)\ll_\eta T^{2+\eta}\).
For the term involving \(\theta_{\mu c_v}\), the coefficients are
\(\alpha_n\mu(n)\).
\leanlinks{\leanref{sum_thetaTailIntegral_le}}

For \(d\equiv7\pmod8\), quadratic reciprocity and \((2/d)=1\) give
\((n/d)=(-d/n^\circ)\). For each \(\psi\bmod q\), let \(\psi^\ast\) be the
primitive character inducing \(\psi\), with conductor \(Q_0\mid q\).
Then \(\psi^\ast\chi_d\) is primitive and nonprincipal, with conductor \(Q_0d\).
The finitely many Euler factors introduced by induction are nonzero in
\(\RePart s>0\). Set \(T^2=qX\) and \(\alpha_n=\psi^\ast(n)\).
Counting each parity separately, the number of primes producing a zero
in the region is \(O(T^{2+\eta}/T^{2\eta})=O(X^{1-\eta/2})\).
There are finitely many characters modulo \(q\), so the bound has the same
exponent after summing over them.

Were every prime in the residue class to produce such a zero, this counting
estimate would imply convergence of \(\sum_{p\equiv d_0\bmod q}(\log p)/p\).
That series diverges in a reduced residue class.
\leanlinks{\leanref{zero_density_quadTwist}}
\end{proof}

\subsection{Zero exclusion from a one-sided bound}

\begin{lemma}[Landau's lemma in the required form]\label{lem:landau}
Let \(U(v)\) be a measurable real function satisfying \(|U(v)|\le C e^v\)
for \(v\ge0\). Suppose its transform \(T(s)=\int_0^\infty U(v)e^{(1-s)v}\,dv\)
has a meromorphic continuation to \(\RePart s>0\) with at most simple poles.
Suppose also that, for some \(\beta>0\), the inequality \(U(v)\le e^{(\beta-1)v}\)
holds for all sufficiently large \(v\), and the real part of the residue of
\(T\) at each real \(\sigma>\beta\) is nonnegative. Then \(T\) is holomorphic
on \(\RePart s>\beta\).
\end{lemma}

\begin{proof}
Choose \(K\) so that \(K e^{\beta v}-U(v)e^v\ge0\) for all \(v\ge0\).
The Laplace transform \(A(s)\) of this nonnegative density satisfies
\(A(s)=K/(s-\beta)-T(s)\) sufficiently far to the right.
Let \(\sigma_c\) be its abscissa of convergence and suppose \(\sigma_c>\beta\).
If \(A\) were holomorphic at \(\sigma_c\), then for \(\sigma_1>\sigma_c\)
sufficiently close to \(\sigma_c\) the Taylor series about \(\sigma_1\) would
converge at some real \(\sigma_2<\sigma_c\). Differentiation under the integral
in the half-plane of convergence identifies its \(j\)-th term with the integral
of the density against \(e^{-\sigma_1 v}((\sigma_1-\sigma_2)v)^j/j!\).
All terms are nonnegative, so monotone convergence makes the sum equal to the
Laplace integral at \(\sigma_2\), and that integral is finite. This contradicts
the definition of the abscissa of convergence.

Thus \(A\) has a pole at \(\sigma_c\). Its residue is the negative of the residue
of \(T\) there, so its real part is nonpositive. At a simple pole, however, the
residue is the limit
\(\lim_{\sigma\downarrow\sigma_c}(\sigma-\sigma_c)A(\sigma)\).
This is a nonnegative real number and is nonzero at a simple pole, hence is
strictly positive. The contradiction gives \(\sigma_c\le\beta\).
The integral, and therefore \(K/(s-\beta)-A(s)\), is holomorphic on
\(\RePart s>\beta\).
\leanlinks{\leanref{IsPolar.landau}}
\end{proof}

\begin{proposition}\label{prop:one-sided-criterion}
Let \(q,d_0\) be as in Lemma \ref{lem:select-quadratic} and let \(\beta>1/2\).
Suppose that for every prime \(d\equiv d_0\pmod q\) one has
\(R_q(y)\le y^{\beta-1}\) for all sufficiently large real \(y\).
Then \(L(\rho,\psi)\ne0\) for every \(\psi\bmod q\) and every \(\rho\ne1\)
with \(\RePart\rho>\beta\).
\end{proposition}

\begin{proof}
Fix \(\rho\) and choose \(d\) by Lemma \ref{lem:select-quadratic}, with
\(H=|\ImPart\rho|\). The residue in Lemma \ref{lem:prime-transform} is then
\(\nu(s)/\bigl(\varphi(q)s(1-s)\bigr)\) for
\(\RePart s>\beta,\ |\ImPart s|\le H,\ s\ne1\).
In particular, it is nonnegative on the real interval \(\beta<s<1\).
For real \(s\ge1\), the entire functions \(L^\circ(s,\psi)\) are nonzero and the
residue vanishes. The function \(U(v)=R_q(e^v)\) is \(O(e^v)\), and the
hypothesis gives \(U(v)\le e^{(\beta-1)v}\) for all sufficiently large \(v\).
By Lemma \ref{lem:landau} the residue at \(\rho\) is zero.
The denominator is nonzero, and \(\nu(\rho)\) is a sum of nonnegative orders
of vanishing, so every summand is zero. Since \(\rho\ne1\), nonvanishing of
\(L^\circ\) implies nonvanishing of \(L\).
\leanlinks{\leanref{LFunction_ne_zero_of_primeRemainder_le}}
\end{proof}

\subsection{The prime number theorem for a fixed modulus}

\begin{lemma}[Uniform logarithmic derivative bound]\label{lem:uniform-log-derivative}
Let \(\chi\) be a Dirichlet character modulo \(Q\), with \(\mathcal L\ge10^{10}\)
and \(Q\le e^{\mathcal L}\). Write \(\delta_\chi=1\) for the principal character and
\(\delta_\chi=0\) otherwise, and put \(f_\chi(s)=L^\circ(s,\chi)/s^{\delta_\chi}\).
If \(\chi^2\) is principal, assume that \(|f_\chi(z)|\ge1/\mathcal L\) for
\(|z-1|\le1/\mathcal L\). Then for \(\RePart s\ge1-\mathcal L^{-25}\) and
\(|s|\le e^{\mathcal L}\) one has
\begin{equation}\label{eq:uniform-log-derivative}
  L^\circ(s,\chi)\ne0,\qquad
  \left|\frac{(L^\circ)'(s,\chi)}{L^\circ(s,\chi)}\right|\le2\mathcal L^{25}.
\end{equation}
\end{lemma}

\begin{proof}
Put \(\kappa_\chi=Q^{-1}\sum_{v\bmod Q}\chi(v)\) and
\(r_\chi(t)=\sum_{1\le n\le t}\chi(n)-\kappa_\chi t\).
Subtracting the contribution of each complete period leaves fewer than \(Q\)
terms, each of absolute value at most \(2\), and a contribution of at most \(1\)
from \(t-\lfloor t\rfloor\). Thus \(|r_\chi(t)|\le2Q\).
Partial summation and analytic continuation give
\[
 f_\chi(s)=\kappa_\chi+(s-\delta_\chi)
            \int_1^\infty r_\chi(t)t^{-s-1}\,dt
 \qquad(\RePart s>0).
\]
Indeed, in \(\RePart s>1\) the \(L\)-function is \(s\) times the integral
plus \(s\kappa_\chi/(s-1)\). For a nonprincipal character \(\kappa_\chi=0\).
For the principal character the factor \((s-1)/s\) removes the pole.
Since \(|s-\delta_\chi|\le|s|\) in \(\RePart s\ge1/2\), the integral gives
\(|f_\chi(s)|\le1+4Q|s|\) there.

Fix \(s\) with \(\RePart s\ge1-1/\mathcal L\) and \(|s|\le3e^{\mathcal L}\).
On the vertical strip of width \(1/3\) with right edge
\(\RePart s+2/\mathcal L\), consider \(g(z)=f_\chi(z)/(z-s+2)\).
Throughout the strip \(|z-s+2|\ge1\) and \(|z|\le4e^{\mathcal L}|z-s+2|\),
so the partial summation bound gives \(|g(z)|\le e^{3\mathcal L}\).
The right edge is in \(\RePart z\ge1+1/\mathcal L\), where bounding the
Dirichlet series by \(\zeta(\RePart z)\) gives \(|g(z)|\le1+\mathcal L\).
The weight of the left edge in the three-lines theorem is \(6/\mathcal L\).
Consequently
\[
 |f_\chi(s)|=2|g(s)|\le2e^{18}(1+\mathcal L)\le\mathcal L^2.
\]
For \(\RePart z\ge 1-1/(2\mathcal L)\) and \(|z|\le2e^\mathcal L\), applying
this bound on the circle of radius \(1/(2\mathcal L)\) gives
\(|f_\chi'(z)|\le2\mathcal L^3\) by Cauchy's estimate.
\leanlinks{\leansrc{Lib/Analytic/PrimeNumberTheorem/LFunctionGrowth}{Partial summation, three lines, and the derivative bound}}

Now set \(\eta=\mathcal L^{-25}\) and take \(v=\sigma+it\) with
\(\sigma\ge1+\eta\) and \(|v|\le3e^\mathcal L/2\).
In the logarithms of the Euler products, the coefficient at each prime
power with \(p\nmid Q\) is multiplied by
\(3+4\cos\theta+\cos(2\theta)=2(1+\cos\theta)^2\ge0\), where \(\theta\) is
the argument of \(\chi(p)^j p^{-ijt}\). Hence
\[
 1\le |L(\sigma,1_Q)|^3 |L(v,\chi)|^4 |L(\sigma+2it,\chi^2)| .
\]
The first \(L\)-value is bounded by \(2/\eta\) using the zeta function in
the region of absolute convergence. If \(\chi^2\ne1_Q\) or
\(|v-1|\ge1/\mathcal L\), the other two \(L\)-values are bounded by
\(2\mathcal L|f_\chi(v)|\) and \(2\mathcal L^3\), respectively.
For a principal character the factor that is restored has modulus
\(|v|/|v-1|\le1+\mathcal L\).
The same estimate applies at the second point since
\(|\sigma+2it-1|\ge|v-1|\) and \(|\sigma+2it|\le3e^\mathcal L\).
Thus \(\eta^3\le256\mathcal L^7|f_\chi(v)|^4\). From
\(256\cdot12^4\le\mathcal L^6\) we obtain \(|f_\chi(v)|\ge12/\mathcal L^{22}\).
When \(\chi^2=1_Q\) and \(|v-1|<1/\mathcal L\), the hypothesis yields
the same bound since \(1/\mathcal L\ge12/\mathcal L^{22}\).

For \(s\) in the region of the lemma, the lower bound applies at \(v=s+2\eta\).
The derivative is bounded by \(2\mathcal L^3\) along the segment \([s,v]\).
It follows that
\[
 |f_\chi(v)-f_\chi(s)|\le4/\mathcal L^{22},
 \qquad |f_\chi(s)|\ge8/\mathcal L^{22}.
\]
Therefore \(|f_\chi'(s)/f_\chi(s)|\le\mathcal L^{25}/4\).
On passing back to \(L^\circ(s,\chi)=s^{\delta_\chi}f_\chi(s)\), the added
term in the logarithmic derivative is \(\delta_\chi/s\), of modulus at
most \(2\). This proves \eqref{eq:uniform-log-derivative}.
\leanlinks{\leansrc{Lib/Analytic/PrimeNumberTheorem/ZeroFreeRegion}{Proof of the zero-free region}}
\end{proof}

For a fixed modulus \(Q\), one has \(f_\chi(1)\ne0\) for every character.
By continuity the hypothesis near \(1\) holds for all sufficiently large
\(\mathcal L\). For the principal character \(f_\chi(1)\) is the residue at
the pole, and for a nonprincipal character the assertion is Dirichlet's
nonvanishing at \(1\). See also \cite[Chapter 11]{mv} for the fixed-modulus
zero-free region.
\leanlinks{\leanref{eventually_norm_logDeriv_entireL_le}}

\begin{lemma}\label{lem:fixed-modulus-pnt}
Fix an odd prime \(d\) and a positive integer \(q\) with \((d,q)=1\).
For every integer \(B\ge0\) and all sufficiently large \(y\), one has
\[
 |R_q(y)|\le(\log y)^{-B}.
\]
If \(z\ge y\ge1\), then \(|R_q(z)-R_q(y)|\le13(z-y)/y\).
\end{lemma}

\begin{proof}
Put \(\mathcal E_q(t)=\sum_{n\le t}\lambda_q(n)-t/\varphi(q)\).
The bounds \(\lambda_q\le2\Lambda\) and \(\sum_{n\le t}\Lambda(n)\le6t\)
give \(|\mathcal E_q(t)|\le13t\). Partial summation gives
\[
 R_q(y)=\int_1^y\mathcal E_q(t)\frac{dt}{t^2}-c_q.
\]
The bound for differences follows. For \(0\le u\le v\le u+1\), it gives
\(|R_q(e^v)-R_q(e^u)|\le13(e^{v-u}-1)\le26(v-u)\). Partitioning into intervals
of length at most \(1\) shows that \(U(v)=R_q(e^v)\) is Lipschitz on \(v\ge0\)
with constant \(26\).

Apply Lemma \ref{lem:uniform-log-derivative} simultaneously to the finite
families \(\psi,\psi\chi_d\) of fixed moduli \(q,qd\). Each difference in the
average in \eqref{eq:log-derivative} has modulus at most \(4\mathcal L^{25}\),
and \(1/|s|\le2\). Thus, for all sufficiently large \(\mathcal L\),
\(F(s)=G_q(s)/s\) is holomorphic on \(\RePart s\ge1-\mathcal L^{-26}\),
\(|s|\le e^\mathcal L\), with
\(|F(s)|\le2(4\mathcal L^{25}+1)\le\mathcal L^{26}\).
\leanlinks{\leanref{eventually_norm_logDerivMean_div_le}}

Take \(u\) large and put \(k=\lfloor u^{2B+4}\rfloor\), \(H=(k+1)/u\),
\(\mathcal L=(2B+5)\log u\), and \(\eta=\mathcal L^{-26}\).
Write
\[
 A=\frac{(-H)^k H}{k!}\mathcal T_q^{(k)}(1+H)
\]
for the normalized Taylor coefficient of
\(\mathcal T_q(s)=(F(s)-F(1))/(s-1)\) at \(1+H\).
Differentiation under the integral expresses \(A\) as the average of \(U(v)\)
against the density \(H^{k+1}v^k e^{-Hv}/k!\), which has mean \(u\) and variance
\(u^2/(k+1)\). For any \(\delta>0\),
\(|v-u|\le\delta+(v-u)^2/\delta\). Combining this inequality with the Lipschitz
bound and choosing \(\delta=u^{-B-1}\) gives \(|A-U(u)|\le52u^{-B-1}\).

The disk of center \(1+H\) and radius \(H+\eta\) lies in the holomorphy region
above. On its boundary \(|s-1|\ge\eta\), so
\(|\mathcal T_q(s)|\le 2\mathcal L^{26}/\eta\). Cauchy's estimate and
\(H/(H+\eta)\le e^{-\eta/(2H)}\) give
\[
 |A|\le4\mathcal L^{26}H\eta^{-1}
       \exp\!\left(-\frac{(k+1)\eta}{2H}\right)
      =4\mathcal L^{52}H e^{-u\eta/2}.
\]
For sufficiently large \(u\) we have \(\mathcal L^{52}\le u\),
\(4H\le u^{2B+5}\), and \(u^{3B+7}\le e^{\sqrt{u}/2}\), so
\(|A|\le u^{-B-1}\). Consequently \(|U(u)|\le53u^{-B-1}\le u^{-B}\).
The substitution \(u=\log y\) finishes the proof.
\leanlinks{\leanref{eventually_abs_le_inv_pow_of_transform}}
\end{proof}

\section{Auxiliary data and a prime-sum identity}\label{sec:data}

\subsection{Number fields, weights, and levels}

Let \(q_0\) be a positive multiple of \(18480\), and let \(d\) be a prime
congruent to \(q_0/2-1\) modulo \(q_0\). Set \(K=\mathbb Q(\sqrt{-d})\).
Then \(d\equiv7\pmod8\). The character \(\chi_d\) from the preceding section
is the quadratic character of \(K\). We also write it as \(\chi_K\).

The admissible data \(D\) include positive integers \(x,x_1,b\) and integers
\(m,r\) that determine the ranks. Set \(a=xb\), \(k=r+1-m\), and \(t=5r\).
The integers \(m,r\) are positive powers of \(2,3\), respectively, and satisfy
\begin{gather*}
 (\log b)^{100}\le x,\qquad
 x^{10^6}\le x_1,\qquad x_1^{10^6}\le b,\qquad
 (\log x_1)^6\le\log b,\\
 \sqrt{x}\le m\le 2\sqrt{x},\qquad x^5\le r\le 3x^5.
\end{gather*}
Admissibility also requires \(\log\log b\ge
10^{12}\operatorname{constantBound}(q_0,d)\), where the constant and the full
set of conditions are fixed in the formalization. \leanref{Parameters.constantBound}
Write \(2^s=m^{22}\), so that \(s\le32\log x\), and set
\[
  l_0=3^{\lfloor\log_3 s\rfloor+4},\qquad
  q_1=3l_0\,2^s q_0,\qquad q_2=(q_0d)^{s+6},\qquad
  E=K(\zeta_{q_1}).
\]
The integer \(b\) is a multiple of \(q_1x_1\).
Identify the set \(\Phi\) of infinite places of \(E\) with
\((\mathbb Z/q_1\mathbb Z)^\times\). Let \(f=|\Phi|\), and let \(\Phi_a\)
be the set of places congruent to \(1\) modulo \(q_0\).
Choose a subset \(\Phi_b\) of the places congruent to \(-1\) modulo \(q_0\),
and impose the following cardinality and weight conditions.
\[
  f=|\Phi_a|\varphi(q_0),\qquad
  |\Phi_a|\le|\Phi_b|\log b,\qquad
  \kappa(\tau)=
  \begin{cases}
    b=a/x,&\tau\in\Phi_b,\\
    a/x_1,&\tau\notin\Phi_b.
  \end{cases}
\]

The set \(\Phi_b\) satisfies a sparsity condition. For any prime \(\ell\le kt\)
and \(0<i<j<l_0\), there are no \(\sigma,\tau,\rho\in\Phi_b\) satisfying
\(\tau/\sigma\equiv\pm\ell^i\) and \(\rho/\sigma\equiv\pm\ell^j\) simultaneously.
The modulus is the part of \(q_1\) prime to \(\ell\), and the two signs are independent.
The formalization specifies the identification of the infinite places
and the sparsity condition.
\leansrc{New/Reduction/Data/FieldsAndPlaces}{Infinite places}. \leanref{Parameters.IsSparse}

The level primes \(p_0,p_1,N\) are all congruent to \(1\) modulo \(q_2^2\).
The distinct primes \(p_0,p_1\) lie in \([b^{10},b^{11}]\), and \(M=p_0p_1\).
Choose algebraic integers \(\pi_M\) in \(K\) and \(\pi_N\) in \(E\) such that
\(\pi_M\overline{\pi_M}=M\) and \(\pi_N\overline{\pi_N}=N\).
Let \(\bar\kappa=\exp(f^{-1}\sum_{\tau\in\Phi}\log\kappa(\tau))\) be the
geometric mean of the weights and
\(\mathrm{rd}_E=|\operatorname{disc} E|^{1/[E:\mathbb Q]}\) the root discriminant.
The range for \(N\) is
\begin{equation}\label{eq:level-window}
 \log N\in[\log N_-,\log N_-+\sqrt{\log b}],\qquad
 \log N_-=
 m\log\bar\kappa+(m-1)\log\frac{M\mathrm{rd}_E}{2\pi}
 -\frac{15}{2}\log b.
\end{equation}
With this choice of level, the term \(a\log N\) offsets the growth over \(m\)
steps in the archimedean estimate below.

\subsection{Existence of admissible data}

\begin{proposition}\label{prop:admissible-data}
Fix \(q_0,d\) as above and let \(0<\varepsilon<1/100\).
For every sufficiently large integer \(x\), admissible data exist.
One may take \(x_1=\lfloor e^{x^{\varepsilon/6}}\rfloor\), \(m\) to be the
least power of \(2\) at least \(\sqrt{x}\), and \(r\) to be the least power
of \(3\) at least \(x^5\).
The least multiple \(b\) of \(q_1x_1\) exceeding \(e^{x^\varepsilon}\) satisfies
\[
 x^{\varepsilon/6}-1\le\log x_1,\qquad
 x^\varepsilon\le\log b\le x^\varepsilon+1.
\]
\end{proposition}

\begin{proof}
Put \(W=x^{\varepsilon/6}\) and choose \(m,r\) as in the statement.
For \(x_1=\lfloor e^W\rfloor\) one has \(W-1\le\log x_1\le W\).
For all sufficiently large \(x\), \(q_1x_1\le e^{W^6}\), so that the
chosen multiple satisfies \(e^{W^6}<b\le2e^{W^6}\).
It follows that \(\log b\le W^6+1\).
Since \(s=O(\log x)\) and \(W\) is a positive power of \(x\), the required
logarithmic scale conditions hold.

For the sparse set, use Behrend's construction. In an \(n\)-dimensional
box with coordinates in \(\{0,\ldots,H-1\}\), at least \(H^n/(nH^2)\)
points have the same squared norm. Reading each point's coordinates as
digits in base \(2l_0H\) gives a set of exponents.
If \(a_1,a_2\) are positive integers with \(a_1+a_2\le2l_0\), then
\[
 (a_1+a_2)e_2\equiv a_1 e_1+a_2 e_3 \pmod {2^s}
\]
implies \(e_1=e_2=e_3\). Indeed, under \((2l_0H)^n\le2^s\), the integer
representatives on both sides are less than \(2^s\), and the equality
holds digit by digit without carries. The middle point is a strict
weighted average of the endpoints unless they coincide. Since all three
points lie on the same sphere, they coincide. \leanref{exists_three_term_free}

Let \(g=1+3l_0q_0\) and take places of the form \(-g^e\).
If \(2^v\) is the highest power of \(2\) dividing \(q_0\), then \(g\) has
order \(2^s\) modulo \(2^{s+v}\). Hence the selected exponents give
distinct places, all congruent to \(-1\) modulo \(q_0\).
For \(\ell=2\) their ratios are \(1\) modulo \(3l_0\), whereas \(4\) has
order \(l_0\) modulo \(3l_0\). This excludes the forbidden ratios.
For an odd prime \(\ell\), raising the two ratios to \(2j\) and \(2i\),
respectively, and comparing them modulo \(2^{s+v}\) gives
\(2j e_2\equiv2(j-i)e_1+2i e_3\pmod{2^s}\).
The property of the exponents gives \(e_1=e_2=e_3\).
Then \(\ell^{2i}\equiv1\pmod{2^{s+1}}\), which by the order bound in the
formalization would imply \(2^s\le(\ell+1)2i\).
The parameters satisfy \(2l_0(kt+1)<2^s\), yielding a contradiction.
Fix \(n\) in terms of \(\varepsilon\), and choose \(H\) to be a power of \(2\)
with \(H^n\le 2^s/W<(2H)^n\). Under the numerical conditions
\[
 W\le 2^s\le W^n,\qquad (2l_0)^n\le W,\qquad
 3l_0 n2^n\le W^3,
\]
one has \(H\le W\) and the encoded box fits within the modulus.
The number of selected places is at least \(3l_0 2^s/W^6\).
Since \(3l_0 2^s=|\Phi_a|\) and \(W^6\le\log b\), this gives the required
cardinality. \leanref{Parameters.exists_isSparse}

A Siegel--Walfisz type estimate supplies the level primes in windows
of width \(W\) for \(\log p\).
Since \(q_0,d\) are fixed and \(s=O(\log x)\), one can fix \(n\) with
\(q_2^2\le W^{n+7}\).
For \(W\le U\le W^{n+7}\) one can choose a prime in the residue class \(1\)
modulo \(q_2^2\) with \(\log p\) in \([U,U+W]\).
The values of \(U\) used here lie in \([\log b,24m\log b]\).
Let \(h_K\) be the class number of \(K\). One can place \(h_K+1\) disjoint
windows in \([10\log b,11\log b]\).
Each prime \(p\) selected from these windows splits in \(K\). If \(I_p,J_p\)
are the two primes above \(p\), their classes satisfy \([I_p][J_p]=1\).
Choose \(I_{p_0},I_{p_1}\) in the same class. The ideal \(I_{p_0}J_{p_1}\)
is principal, and one may take its generator to be \(\overline{\pi_M}\).
For \(N\), start the window at \(\log N_-\).
Since \(W\le\sqrt{\log b}\), the window lies in the range of
\eqref{eq:level-window}.
Since \(N\equiv1\pmod4\), one can write \(N=u^2+v^2\).
As \(E\) contains \(\sqrt{-1}\), one may take \(\pi_N=u+v\sqrt{-1}\).
The existence argument in the formalization supplies the windows and
verifies the numerical conditions used above. \leanref{eventually_exists_data}
\end{proof}

\subsection{Source and target spaces of the lifts}

The construction uses Hermitian spaces \(X_j,Y_j\) over \(E\).
Their dimensions are \(\dim_E X_j=k+j(t+1)\) and \(\dim_E Y_j=r+jt\),
and \(Y_m\subset X_m\) is a definite subspace of codimension \(1\).
Let \(S_0\) be the space of algebraic automorphic forms on \(X_0\) of the
specified level, with weight \((a,0,\ldots,0)\) at places in \(\Phi_a\)
and weight \(0\) at the remaining infinite places.
The level is principal congruence at the selected places above \(M\).
Modulo \(\pi_N\) the last row is congruent to \((0,\ldots,0,1)\).
The subspace \(S\subseteq S_0\) consists of forms whose average under the
central scalar action supported at \(w,\bar w\) vanishes for every
place \(w\) dividing \(\pi_N\).
\leansrc{New/UpperBound/Objects/Columns}{Source forms and their level}.
\leanref{ramifiedSourceForms}

The principal-level forms on \(Y_0\) carry an action of
\(\bar G=\prod_w\mathrm{GL}_r(\mathbb F_w)\), where \(w\) ranges over the
finite places above \(p_0,p_1\) selected by \(\pi_M\).
Let \(V\) be the image of the projection given by the central idempotent
\(e_0\), and set \(h=\dim_{\mathbb C}V\).
The idempotent \(e_0\) sums the projectors for the irreducible characters
with the following labels. At each \(w\) the label is a monic squarefree
polynomial of degree \(r\) with nonzero constant term, and none of its
roots in the residue field is a \(q_2\)-th root of unity.
Moreover, above each of \(p_0,p_1\) at least one label is a primitive
polynomial, with a root that generates the multiplicative group of the
degree-\(r\) extension.
\leansrc{New/UpperBound/Objects/LabelsAndIndexTypes}{The finite group and its central idempotent}.
\leansrc{New/UpperBound/Objects/Rows}{The target space and its archimedean type}.

Each twist \(\xi\) gives \(m\) successive theta lifts along
\(Y_0\to\cdots\to Y_m\). Their composite restricted to \(V\) is \(\Theta_\xi\).
The character in the \(j\)-th lift is the fixed character \(\mu_j^0\) of
the formalization multiplied by \(\xi^j\).
Write \(1\) for the trivial twist. Applying the normalization and projection of
Section~\ref{sec:upper-bound} to the Poincar\'e map from \(X_0\) to \(Y_0\) gives
\(\mathfrak p_g\colon S\to V\).
\leansrc{New/UpperBound/Objects/RowChain}{The chain of lifts}.
\leansrc{New/Reduction/IntegralTargetForms}{The integral structure for \(\Theta_\xi\)}.

\subsection{Prime sums in the local divisibility estimates}

The image of an integer in \(\Phi\) is its residue class if it is prime
to \(q_1\), and the identity element otherwise. With this convention we
write \(-\nu\tau\). Put \(u_+=\max(u,0)\), and for \(0\le y\le a\) set
\begin{align*}
 s_{\mathrm{fin}}(y,\tau)
   &=\sum_{1\le\nu\le a}\lambda_1(\nu)
      \left(\frac{y}{\nu}-\kappa(-\nu\tau)\right)_+,\\
 s_{\infty}(j)&=\log(j!)+j\log(M/\Omega),\qquad
 \Omega=
    \exp\!\left(\gamma-\RePart\frac{L'}{L}(1,\chi_d)\right)
        \frac{2\pi}{\mathrm{rd}_E},\\
 s_{\mathrm{main}}&=a\bigl(\log(a/\bar\kappa)+c_1\bigr),\qquad
 \mathcal R=\sum_{\tau\in\Phi_a}\bigl(s_{\mathrm{fin}}(a,\tau)-s_{\mathrm{main}}\bigr),\\
 A_0&=\sum_{\tau\in\Phi}
       \bigl(s_\infty(\kappa(\tau))-s_{\mathrm{fin}}(\kappa(\tau),\tau)\bigr),
 \qquad H_0=mkA_0,\qquad E_j=fa(\log b)^j .
\end{align*}
Here \(\gamma\) is Euler's constant and \(c_1=G_1(1)\) is the constant of
\eqref{eq:prime-remainder}.
We will use \(H_0\ge0\), which follows from
\(s_\infty(\kappa(\tau))-s_{\mathrm{fin}}(\kappa(\tau),\tau)
\ge\kappa(\tau)\log b\) at every \(\tau\).
This inequality uses the scale conditions and
\(s_{\mathrm{fin}}(y,\tau)\le4y\log x_1\). \leanref{primeSum_le}

\begin{proposition}\label{prop:total-remainder}
For admissible data,
\begin{equation}\label{eq:total-remainder}
 \mathcal R =
 a\left[
 |\Phi_b|\bigl(R_{q_0}(x)-R_{q_0}(x_1)\bigr)
 +|\Phi_a|R_1(x_1)
 \right].
\end{equation}
\end{proposition}

\begin{proof}
The number of \(\tau\in\Phi_a\) for which \(-\nu\tau\in\Phi_b\) is
\(|\Phi_b|\) when \(\nu\equiv1\pmod{q_0}\), and zero otherwise.
Indeed, if \(\nu\) is prime to \(q_1\), multiplication determines the
image residue class, and \(\Phi_b\) lies in the class \(-1\) modulo \(q_0\).
If \(\nu\) is not prime to \(q_1\), the image of \(-\nu\) is the identity.
Since \(q_1\) and \(q_0\) have the same prime divisors, in this case
\(\nu\) cannot be \(1\) modulo \(q_0\).

Thus, to evaluate \(\sum_{\tau\in\Phi_a}s_{\mathrm{fin}}(a,\tau)\), first
assign the weight \(a/x_1\) at all \(|\Phi_a|\) places. For each term with
\(\nu\equiv1\pmod{q_0}\), change that weight to \(a/x\) at \(|\Phi_b|\) places.
Writing the positive part \(\bigl(a/\nu-a/y\bigr)_+\) as a sum restricted
to \(\nu\le y\) gives
\[
 \frac1a\sum_{\tau\in\Phi_a}s_{\mathrm{fin}}(a,\tau)
  =
 |\Phi_a|B_1(x_1)+|\Phi_b|\bigl(B_{q_0}(x)-B_{q_0}(x_1)\bigr),
 \qquad B_q(y)=R_q(y)+\frac{\log y}{\varphi(q)}+c_q .
\]
The constant \(c_{q_0}\) cancels in the difference. The definition of
the weights also gives
\[
 |\Phi_a|\bigl(\log(a/\bar\kappa)+c_1\bigr)
 =|\Phi_a|(\log x_1+c_1)
   +\frac{|\Phi_b|}{\varphi(q_0)}(\log x-\log x_1).
\]
This is the main term of the preceding sum. Subtracting it yields
\eqref{eq:total-remainder}. \leanref{totalRemainder_eq}
\end{proof}

\section{Twisted lattices and the product formula}\label{sec:lattices}

\subsection{Integrality at finite places and size at infinity}

Let \(\iota\) denote a prime \(\ell\) together with a field isomorphism
\(\mathbb C\simeq\overline{\mathbb Q_\ell}\). Write
\(|z|_\iota=|\iota(z)|_\ell\), with the normalization \(|\ell|_\ell=\ell^{-1}\),
and let \(\bar\iota\) be the isomorphism obtained by precomposing with complex conjugation.
A twist \(\xi\) is a character of \(U(1)\) of weight \(0\) at the infinite places.
Its set of ramified places \(Q_\xi\) is disjoint from the places above \(MN\)
and from its conjugate. For each \(1\le j\le m\), the power \(\xi^j\) is ramified
at every place of \(Q_\xi\). We say that \(\xi\) is admissible at \(\ell\) if the
norm \(N_\xi\) of its modulus is coprime to \(\ell\).
\leansrc{New/Reduction/Data/LevelAndTwists}{Twists and admissibility}.

\begin{editorialnote}
The index \(\iota\) specifies both the prime and the field isomorphism.
The distinction between \(\iota\) and \(\bar\iota\) above the same prime \(\ell\)
accounts for the complex conjugation in one argument of the Hermitian inner product.
\end{editorialnote}

Let \(W_\xi\) be the space of forms into which \(\Theta_\xi\) maps, and let
\(W^{\mathrm{int}}_{\xi,\iota}\) be its module of forms integral at \(\iota\).
Integrality means that the \(\ell\)-adic Weyl gauge of every value of the form
is at most \(1\). For \(\xi\) admissible at \(\ell\), we have
\[
  \phi\in W^{\mathrm{int}}_{\xi,\bar\iota},\quad
  \psi\in W^{\mathrm{int}}_{\xi,\iota}
  \quad\Longrightarrow\quad
  |\langle\phi,\psi\rangle|_\iota\le1 .
\]
Indeed, the inner product is the sum of the Weyl pairings at class
representatives. At this level the stabilizers of the representatives are
trivial, so their orders contribute no denominators. At each point of
evaluation, write \(A\) for the local matrix, \(\rho_{\mathrm W}\) for the
Weyl representation, and \(G\) for the integral Gram matrix. After conjugating
the first argument, apply \(\rho_{\mathrm W}(A^{\mathsf T})\) to it and
\(\rho_{\mathrm W}(A^{-1})\) to the second argument. The gauge bounds imply that
both resulting vectors have integral coordinates. Since
\(\rho_{\mathrm W}(A^{\mathsf T})^{\mathsf T}G\rho_{\mathrm W}(A^{-1})=G\),
the pairing computed in these integral coordinates equals the
\(\iota\)-image of the original pairing. \leanref{gramForm_mem_embValRing}
For the trivial twist, take as coordinates the values at class
representatives in the standard Weyl basis. The inner product in these
coordinates has a positive-definite matrix with rational entries.
\leanref{Data.targetLattice}

\begin{definition}
A vector \(y\in V\) is integral at \(\iota\) if \(\Theta_\xi y\) lies in
\(W^{\mathrm{int}}_{\xi,\iota}\) for every \(\xi\) admissible at \(\ell\).
A finite family \(y=(y_1,\ldots,y_n)\) has logarithmic index at most \(\delta\)
if there are nonnegative reals \(\delta_\ell\) with the following properties.
\begin{enumerate}
\item For every finite set \(P\) of primes,
\(\sum_{\ell\in P}\delta_\ell\log\ell\le\delta\).
\item For each \(\iota\) above \(\ell\), there is an \(n\times n\) matrix
\(U_\iota\) with entries integral at \(\iota\) such that every component of
\(yU_\iota\) is integral at \(\iota\) and
\(|\det U_\iota|_\iota\ge\ell^{-\delta_\ell}\).
\end{enumerate}
The same \(U_\iota\) serves for every twist admissible at that prime.
\end{definition}

Apply \(\sigma\in\operatorname{Aut}(\mathbb C)\) to the coordinates of the
untwisted lift, and write \(\|\Theta_1 y_i\|_\sigma^2\) for the squared norm
computed using the rational inner product matrix. If the Gram determinant
of the lifted family \(y\), computed with these coordinates and this inner
product, is at most \(e^c\) for every \(\sigma\), we say that the logarithm of
the squared covolume is at most \(c\). In particular, if
\(\|\Theta_1 y_i\|_\sigma^2\le e^B\) for every \(i,\sigma\), then Hadamard's
inequality gives the bound \(c=nB\). \leanref{IndexLe}

\subsection{The product formula for heights}

\begin{definition}
A family \(u=(u_\iota)_\iota\) has height at most \(C\), written
\(\operatorname{ht}(u)\le C\), if there are a common multiplier
\(H\in\mathbb C^\times\), a real number \(\alpha\), and nonnegative reals
\(e_\ell\) such that for every \(\sigma,\iota\) and every finite set \(P\)
of primes,
\[
 |\sigma H|\le e^\alpha,\qquad
 |H u_\iota|_\iota\le\ell^{e_\ell},\qquad
 \alpha+\sum_{\ell\in P} e_\ell\log\ell\le C .
\]
\end{definition}

When multiplying two families, we can multiply their common multipliers
and add their \(\alpha\)'s and \(e_\ell\)'s. Thus height bounds add under
multiplication. Every upper bound for the height of the constant family
\(1\) is nonnegative. To see this, take a multiplier \(H\) for that family.
Its absolute values \(|H|_\iota\) are bounded uniformly in \(\iota\) above
\(\ell=2\), so \(H\) is algebraic. In fact, a transcendental \(H\) can be
mapped by a suitable isomorphism to any transcendental element of
\(\overline{\mathbb Q_2}\). The elements algebraic over \(\mathbb Q\)
are countable, and their complement is dense, so there are transcendental
elements of arbitrarily large absolute value. The extension of the
isomorphism is the one used in the algebraicity criterion.
\leanref{isAlgebraic_of_norm_complexToCp_le}

Set \(Q=\operatorname{Nm}_{\mathbb Q(H)/\mathbb Q}H\) and
\(n_H=[\mathbb Q(H):\mathbb Q]\). Then \(Q\ne0\). Writing \(Q\) as a product
of conjugates of \(H\) gives \(\log|Q|\le n_H\alpha\) at the infinite place
and \(\log|Q|_\ell\le n_H e_\ell\log\ell\) at the finite places. Take \(P\)
to contain the primes dividing the numerator or denominator of \(Q\).
The product formula over \(\mathbb Q\) gives
\[
 0=\log|Q|+\sum_{\ell\in P}\log|Q|_\ell
   \le n_H\left(\alpha+\sum_{\ell\in P} e_\ell\log\ell\right)
   \le n_H C .
\]
Thus \(C\ge0\). In particular, if the constant family \(z\ne0\) has height
at most \(C\) and \(|\sigma z|\le e^{C'}\) for every \(\sigma\), then
\(C+C'\ge0\). Indeed, the constant family \(z^{-1}\) has height at most
\(C'\) with multiplier \(H=z\), and the two families have product \(1\).
\leanref{HeightLe.nonneg_of_one}

\subsection{The Gram determinant of a basis}

Assume that the maps \(\Theta_\xi\) are injective and put
\(Z(\xi)=\det(\Theta_\xi^*\Theta_\xi)\). For each \(\iota\), choose a
twist \(\xi_\iota\) admissible at its prime. Let \(y=(y_1,\ldots,y_n)\)
be linearly independent, and complete it by \(z=(z_1,\ldots,z_{n'})\)
to a basis of \(V\). Write \(w=(y,z)\) and set
\[
 G_\xi(w)=\bigl(\langle\Theta_\xi w_i,\Theta_\xi w_j\rangle\bigr)_{ij},
 \qquad \Delta_\xi(w)=\det G_\xi(w).
\]

\begin{proposition}\label{prop:gram-product}
Suppose that the logarithm of the squared covolume of \(y\) is at most
\(c\) and that its logarithmic index is at most \(\delta\). For each twist,
suppose there is a scalar \(\gamma(\xi)\ne0\) such that
\(\gamma(\xi)\Theta_\xi z_j\) is integral at every place where the twist is
admissible. Suppose also that \(\|\Theta_1 z_j\|_\sigma^2\le e^{B_{\mathrm{Ar}}}\)
for all \(\sigma,j\). If the families \((|\gamma(\xi_\iota)|^{-2})_\iota\)
and \((Z(1)/Z(\xi_\iota))_\iota\) have heights at most
\(B_{\mathrm{fin}}\) and \(B_Z\), respectively, then
\begin{equation}\label{eq:gram-product}
  c+2\delta+n'(B_{\mathrm{Ar}}+B_{\mathrm{fin}})+B_Z\ge0.
\end{equation}
\end{proposition}

\begin{proof}
Since \(\bar\iota\) lies above the same \(\ell\) as \(\iota\), both index
matrices \(U_\iota,U_{\bar\iota}\) for \(y\) can be used with
\(\xi=\xi_\iota\). Put
\[
  P_\iota=\operatorname{diag}(U_\iota,\gamma(\xi)I_{n'}),\qquad
  P_{\bar\iota}=\operatorname{diag}(U_{\bar\iota},\gamma(\xi)I_{n'}).
\]
Every entry of \(\overline{P_{\bar\iota}}^{\,t}G_\xi(w)P_\iota\) pairs a form
integral at \(\bar\iota\) with one integral at \(\iota\), so the entry is
integral at \(\iota\). Taking determinants gives
\begin{align*}
 1&\ge|\det P_{\bar\iota}|_{\bar\iota}
             |\Delta_\xi(w)|_\iota\,|\det P_\iota|_\iota,\\
 \bigl|\,|\gamma(\xi)|^{2n'}\Delta_\xi(w)\,\bigr|_\iota
     &\le\ell^{2\delta_\ell}.
\end{align*}
The second inequality uses
\(\overline{\gamma(\xi)}\,\gamma(\xi)=|\gamma(\xi)|^2\) and the lower bounds
on the determinants of the two index matrices. With \(H=1\), the last
inequality gives height at most \(2\delta\) for the family
\((|\gamma(\xi_\iota)|^{2n'}\Delta_{\xi_\iota}(w))_\iota\).

Since \(w\) is a basis of \(V\), the Gram determinant obeys
\(\Delta_\xi(w)=Z(\xi)\det(\langle w_i,w_j\rangle)_{ij}\).
Therefore, for each \(\iota\),
\[
 \Delta_1(w)
 =\bigl(|\gamma(\xi_\iota)|^{2n'}\Delta_{\xi_\iota}(w)\bigr)
     \bigl(|\gamma(\xi_\iota)|^{-2}\bigr)^{n'}
     \frac{Z(1)}{Z(\xi_\iota)} .
\]
It follows that the constant family \(\Delta_1(w)\) has height at most
\(2\delta+n'B_{\mathrm{fin}}+B_Z\).

For the Archimedean bound, apply Fischer's inequality to the Gram
determinant of the full basis after applying \(\sigma\) to the coordinates.
It is at most the product of the determinants of the diagonal blocks for
\(y\) and \(z\). Hadamard's inequality on the auxiliary block for \(z\)
then bounds the determinant of the whole basis by \(e^{c+n'B_{\mathrm{Ar}}}\).
To bound the conjugate \(\sigma(\Delta_1(w))\) itself for an arbitrary
field automorphism \(\sigma\), we use exterior products to account for
complex conjugation. Put \(\sigma^\dagger(z)=\overline{\sigma(\bar z)}\).
The inner product matrix is rational, so \(\sigma(\Delta_1(w))\) is the
inner product of the two exterior products formed using \(\sigma^\dagger\)
and \(\sigma\), respectively. Each has squared norm at most
\(e^{c+n'B_{\mathrm{Ar}}}\). The Cauchy--Schwarz inequality gives
\(\left|\sigma(\Delta_1(w))\right|\le e^{c+n'B_{\mathrm{Ar}}}\).

Injectivity of \(\Theta_1\) gives \(\Delta_1(w)\ne0\). Combining its height
and Archimedean bounds by the product formula proves \eqref{eq:gram-product}.
\leanref{covolume_add_index_nonneg}
\end{proof}

\section{Upper bounds for the theta lifts}\label{sec:upper-bound}

The ramified source subspace \(S\) provides forms whose lifts have
controlled length and whose images in \(V\) have controlled index.
The index estimate applies to all twists allowed at each finite place.
The length estimate concerns the trivial twist.

\begin{theorem}[Upper bounds]\label{thm:upper-bound}
\begin{enumerate}
\item The forms \(\phi\in S\) that are integral at every finite place and satisfy
\[
  \|\Theta_1\mathfrak p_g\phi\|_\sigma^2
  \le \exp\!\bigl(H_0-m\mathcal R+E_4\bigr)
\]
for all \(g\in\overline G\) and \(\sigma\in\operatorname{Aut}(\mathbb C)\)
span \(S\).
\item Let \(\mathcal T\subseteq\overline G\) be finite.
Suppose that every member \(y_i\) of a family can be written as
\(y_i=\mathfrak p_{g_i}\phi_i\), with \(g_i\in\mathcal T\) and
\(\phi_i\in S\) integral at every finite place.
Then the logarithmic index of \((y_i)_i\) is at most \(2|\mathcal T|(\dim S)E_3\).
\end{enumerate}
\end{theorem}

\subsection{The seesaw identity and the projection}

The principal-level forms on \(Y_0\) carry an action of \(\overline G\),
the product of the general linear groups over the residue fields at the
selected places above \(M\). The central idempotent \(e_0\) of its group
algebra projects onto \(V\). Write \(\rho\) for the group action and
\(\mathcal P\) for the Poincar\'e series map. The lower chain of lifts
from principal level, along \(Y_0\to\cdots\to Y_m\), is denoted by \(L_\xi\).
Its restriction to \(V\) is \(\Theta_\xi\).

\emph{Unfolding the integrals of the lower chain.}
Let \(F=E^+\), and take adeles over \(F\) here.
Write \(G_j=U(Y_j)\) and \([G_j]=G_j(F)\backslash G_j(\mathbb A)\). Let \(J^\circ_j\) be the stabilizer of the standard
lattice at the finite places and \(J_{\xi,j}\subset J^\circ_j\) the
level for the twist \(\xi\).
The lower theta kernel \(K_{\xi,j}(g,h)\) is the theta value of the
Schwartz datum after the Weil action of
\((g,h)\in G_j(\mathbb A)\times G_{j+1}(\mathbb A)\).
The lower lift \(L_{\xi,j}\) integrates \(K_{\xi,j}(g,h)f(g)\)
over \([G_j]\) and multiplies by \(d_{\xi,j}=[J^\circ_j:J_{\xi,j}]\).
For a form \(z\) of conjugate type, put
\[
 L_{\xi,j}^{\mathrm t}z(g)=d_{\xi,j}
          \int_{[G_{j+1}]}K_{\xi,j}(g,h)z(h)\,dh .
\]
Exchanging the order of integration gives
\[
  \int_{[G_{j+1}]} z\,L_{\xi,j} f
      = \int_{[G_j]}(L_{\xi,j}^{\mathrm t}z)f .
\]
Apply this identity from \(j=m-1\) down to \(0\).
The pairing with the test form \(x_m\) becomes the integral of
\(\mathcal P\phi\) against the successive descents of \(\overline{x_m}\).
Rapid decrease gives the absolute convergence used in this step.
\leansrc{New/UpperBound/Seesaw/SeesawIdentity}{Transpose of the lift and absolute convergence}.

Let \(R_j\) be the Bessel subgroup for the grading of \(Y_j\).
The initial function \(\mathcal W_0\) on \([R_0]\) extends the
matrix coefficient of \(\phi\) on \(X_0\) by the prescribed
unipotent character. After \(j\) unfoldings, write \(\mathcal W_j\)
for the resulting function with the earlier outer variables fixed.
Fix the next outer variable \(g_j=(a_j,y_j)\), its infinite and
finite components. Put \(W_{j,g_j}(r)=\mathcal W_j(rg_j)\) for \(j>0\)
and \(W_{0,g_0}(r)=\mathcal W_0(r)\) at the initial step.
Integrate this function times the kernel over \(r\in[R_j]\).
Unfolding the theta sum leaves the constant term along maps that
raise the grading, giving the next function \(\mathcal W_{j+1}\).
Thus
\[
 \begin{aligned}
 \int_{[R_j]}W_{j,g_j}(r)(L_{\xi,j}^{\mathrm t}z)(rg_j)\,dr
   &=d_{\xi,j}\,\Xi_{j+1}(z),\\
 \Xi_{j+1}(z)
   &=\int_{R_{j+1}(F)\backslash G_{j+1}(\mathbb A)}
                   \mathcal W_{j+1}(h)z(h)\,dh .
 \end{aligned}
\]
\leansrc{New/UpperBound/Seesaw/UnfoldingStages}{One step of the Bessel-period unfolding}.

At the first step, decompose the support at principal level \(J_M\)
into the classes \(\Omega_0=J_M/J_{\xi,0}\).
Unfolding the Poincar\'e series gives the average over these classes,
so the remaining factor is
\[
 \frac{d_{\xi,0}}{[J_M:J_{\xi,0}]}=[J^\circ_0:J_M],
\]
independent of the twist.
For \(1\le j<m\), a finite set \(\Omega_j\) of classes modulo
\(J_{\xi,j}\) meets the \(R_j(\mathbb A_f)\)-orbits of the finite
components on which \(\mathcal W_j\) is nonzero.
The multiplicity \(m_j(c)\) of \(c\in\Omega_j\) counts the classes
of elements of \(R_j(\mathbb A_f)\), modulo its intersection with
the standard lattice stabilizer, that move \(c\) back into the set.
Weighting by \(1/m_j(c)\) gives orbital integral \(1\).
Since \(J^\circ_j\) has volume \(1\), each \(J_{\xi,j}\)-coset has
volume \(d_{\xi,j}^{-1}\), which cancels the \(d_{\xi,j}\) in the
period. The remaining infinite integrals use the measures of this
weighted decomposition, and the finite sum runs over
\(\prod_{j<m}\Omega_j\).
The first multiplicity is \(1\), and each other multiplicity
divides a power of \(\mathrm NQ(\xi)\), the norm of the twist modulus.
\leansrc{New/UpperBound/Seesaw/IntegralAsCosetSum}{Orbital weights and the volume of the level}.
\leansrc{New/UpperBound/Seesaw/Cut}{Finite classes and their multiplicities}.

\emph{Decomposing the kernels in pairs.}
Fix the outer variables and take the product of two consecutive
constant terms for \(Y_j,Y_{j+1},Y_{j+2}\), with \(j\) even.
Reverse the first kernel by taking the adjoint of its map, and join
it with the second to obtain a datum on
\(\operatorname{Hom}(Y_{j+1},Y_j^{-}\oplus Y_{j+2})\), where \(Y_j^{-}\)
has the negative Hermitian form. If \(T,T'\) are the two maps along
the graded strings, the translation of the joined datum is
\((T^\ast,T')\).
The two means over maps that raise the grading give the projection
of the joined datum to its degree-zero part.
On the unipotent subgroup of \(R_j\), the characters of
\(\mathcal W_j\) and the constant term cancel.
Averaging over the unipotent subgroup of \(R_{j+1}\) gives the mean
over the remaining isotropic coordinates in the hyperbolic planes.
The identity \(-TT^\ast+(T')^\ast T'=0\) makes the Heisenberg phase
for joining the translations equal to \(1\).
The remaining kernel, up to a measure constant, is a theta kernel
on \(\operatorname{Hom}(X_{j+1},X_j^{-}\oplus X_{j+2})\).
Integrate first over \([U(X_{j+1})]\) and then over \([U(X_j)]\)
against the preceding matrix coefficient.
For products of Fock polynomials the kernel factors into the product
of the two definite theta kernels, so the two integrals give the
matrix coefficient of their composite lift.
Density of Fock polynomials and continuity of the integrals give the
identity for the infinite component used here.
\leansrc{New/UpperBound/Seesaw/UnfoldingStages}{Product of constant terms and the definite double integral}.

The infinite and finite factors in this calculation transform separately.
The infinite weight of the twist is fixed. The infinite factor therefore
changes only the linear functional used to take the matrix coefficient,
independently of the twist, \(\phi\), and the finite outer variables.
The finite factor gives the evaluated data \(\mathcal E_i(\xi;\boldsymbol c)\)
for each pair at the tuple of classes \(\boldsymbol c\).
Induction by pairs gives the final function \(\mathcal W_m\) as a matrix
coefficient of the lift along \(X_0\to\cdots\to X_m\) with datum
\[
  \mathcal E(\xi;\boldsymbol c)
  =\bigotimes_{i=0}^{m/2-1}\mathcal E_i(\xi;\boldsymbol c).
\]

At the top, \(Y_m\) is definite and \(R_m=U(Y_m)\).
Its point \(h\) embeds into \(X_m\) as \(\operatorname{diag}(h,1)\).
In integrating the matrix coefficient against \(\overline{x_m}\),
orthogonality at infinity gives the finite class mean after applying
the integral branching map. The inner product with the restriction to
\(Y_m\) is that same mean divided by the dimension of the Weyl module
on \(Y_m\) and by \(\prod_\tau\tau(\overline{\pi_N})^{u_\tau}\), where
\(u_\tau\) is the branching exponent.
Thus the top period is the pairing of the restricted upper lift with
\(x_m\) times a constant depending only on the outer infinite variables.
Integrating this constant and multiplying by \([J^\circ_0:J_M]\)
gives \(c_0\). It follows that \(\langle x_m,L_\xi\mathcal P\phi\rangle\)
is \(c_0\) times the pairing with the restricted upper lift for the datum
\[
 \mathcal A_\xi
   =\sum_{\boldsymbol c\in\prod_{j<m}\Omega_j}
         \left(\prod_{j<m}m_j(c_j)^{-1}\right)\mathcal E(\xi;\boldsymbol c).
\]
Since the equality holds for every test form \(x_m\), it identifies
the lifts themselves.
\leansrc{New/UpperBound/Seesaw/SeesawIdentity}{Top branching map and equality of the two chains}.

\emph{Clearing the denominators of the finite data.}
Set \(\mathsf B=M^{f n_m^2}\), with \(n_j=\dim_E X_j\).
Finite Fourier expansion expresses each evaluated datum as a sum of
glued translated data. The denominator for one pair divides
\(\mathsf B^4\), and the factor relating the vacuum to the integral
vacuum is independent of the twist.
Thus \(\mathcal A_\xi=\sum_t r_t C\,\mathcal V_t\), with \(\mathcal V_t\)
the glued translates of the integral vacuum, \(C\) independent of
the twist, and \(\mathsf B^{2m}r_t\) integral at places prime to
\(\mathrm NQ(\xi)\).
For the trivial twist the sum of the absolute values of the coefficients
of each evaluated tensor is at most \(1\), and each sum over the
classes contributes at most \(\mathsf B\). Thus
\(\sum_t|\sigma(r_t)|\le\mathsf B^m\).
The level actions permute the tuples of classes and preserve their
multiplicities. This gives joint invariance of \(\mathcal A_\xi\)
under the intermediate levels of the upper chain.
\leansrc{New/UpperBound/Seesaw/GoodDatum}{Finite Fourier expansion of the glued data}.

If \(c_0C/\mathsf B^{2m}\ne0\), take it for \(c_{\mathrm{ss}}\) and
put \(\Psi_\xi(t)=\mathsf B^{2m}r_t\).
If \(c_0C=0\), the composite map vanishes and one may take
\(c_{\mathrm{ss}}=1\) and \(\Psi_\xi=0\).
At places prime to \(M\mathrm NQ(\xi)\), the translations can be
replaced by zero with the same effect on the vacuum.
Writing \(B_{\xi,\Psi_\xi}\) for the corresponding upper composite
lift, we obtain
\begin{equation}\label{eq:upper-seesaw}
 L_\xi(c_{\mathrm{ss}}^{-1}\mathcal P\phi)=B_{\xi,\Psi_\xi}\phi
 \qquad(\phi\in S).
\end{equation}
The translations are supported above \(M\mathrm NQ(\xi)\), and
multiplication by \(\overline{\pi_M}\) makes them integral at places
prime to \(\mathrm NQ(\xi)\).
The coefficients are integral at the same places. For the trivial twist,
\(\sum_t|\sigma(\Psi_1(t))|\le\mathsf B^{3m}=M^{3fmn_m^2}\).
The glued tensor is jointly invariant under the intermediate levels.\leanref{exists_isAdmissibleCombination_seesawDatum}

For the estimates on the composite lift, expand each gluing as a
finite sum modulo \(N\). Each term adds the last column \(-y/N\) to
one translation and the last row \(-\overline y/N\) to the next.
Collecting the terms gives a sum of unglued translated data with
the same integrality condition on the coefficients.
The translations may have denominators at \(N\). For the trivial
twist the sum of the absolute values of the coefficients is at most
\(\exp(x^{60})\).
\leansrc{New/UpperBound/Divisibility/ChainDatumExpansion}{Expansion of the glued tensor}.

After projection, the lift is a sum of right translates of \(L_\xi z\).
For a principal-level form \(z\) and a finite adelic point \(u\) of \(U(Y_m)\),
\begin{equation}\label{eq:upper-project}
 L_\xi\bigl(|\overline G|e_0\rho(g)z\bigr)(u)
 =
 \sum_{v\in\overline G}|\overline G|e_0(v)\,
            (L_\xi z)\bigl(u\,\widetilde{vg}_{\,\xi}\bigr).
\end{equation}
Here \(\widetilde{vg}_{\,\xi}\) is the lift of \(vg\) to the first block at
the corresponding level. Equation \eqref{eq:upper-project} follows from
the equivariance of \(L_\xi\) for the group action.
The expression of the central projector in irreducible characters of
\(\overline G\) shows that \(|\overline G|e_0(v)\) is an algebraic integer
and \(|\sigma(e_0(v))|\le 1\).
Right translation preserves the lattices used to test integrality.
The ultrametric triangle inequality therefore shows that whenever
\(cL_\xi z\) is integral, so is \(cL_\xi(|\overline G|e_0\rho(g)z)\).
If the squared norm of each value is at most \(\exp(B)\), the projected
form has pointwise norm at most \(|\overline G|^2\exp(B/2)\).
For the trivial twist the number of classes is at most \(\exp(x^{30})\).
Squaring the pointwise estimate and summing over classes gives an upper
bound of \(\exp(B+x^{99})\).
\leansrc{New/UpperBound/Size/ArchimedeanSize}{Integrality and norm estimates for the projection}.

\subsection{Theta derivatives and contributions from the finite places}

Put \(K=\mathbb Q(\sqrt{-d})\) and normalize the valuation at \(p\)
by \(v_p(p)=1\). If \(p\) is nonsplit in \(K\), set
\[
  \pi_p=\begin{cases}p & (p\ne d),\\ \sqrt{-d}&(p=d),\end{cases}
  \qquad
  \rho_p=\begin{cases}(p^2-1)^{-1}&(p\ne d),\\
                        (2p-2)^{-1}&(p=d).\end{cases}
\]
Let \(I_p^0\subseteq\mathbb Z_{\ge0}\) be the positive odd integers if \(p\)
is inert and all nonnegative integers if \(p=d\).
Set \(I_p=I_p^0\setminus\{0\}\). At a split prime both sets are empty.
Write \(\tau_e=(-p^e)\tau\) for the action on the labels of the infinite places.
As in the formalization, the multiplier is the identity when \(-p^e\)
does not represent a unit modulo the cyclotomic modulus.

\paragraph{Normal form for the theta series}
Let \(p\) be nonsplit. Fix a \(p\)-adic embedding and write its images in
\(\mathbb C_p\). Write \(\mathcal O\) for the valuation ring and put
\(q=p^2\) in the inert case and \(q=p\) in the ramified case, so that
\(\rho_p=v_p(\pi_p)/(q-1)\). Let \(\theta\) denote the theta series of
the CM elliptic curve, and let \(y=z+O(z^2)\) be the normalized local
parameter. For the chosen point \(t\) of order two, the translated
coordinate \(x_{\rm CM}(z)=\wp(z+t)-\wp(t)\) belongs to
\(\mathcal O[[y]]\) and satisfies \(x_{\rm CM}=z^2+O(z^3)\).

\emph{Rescaling the formal parameter.}
Complex multiplication by an integer of \(K\) preserves
\(\mathcal O[[y]]\). Passing to the \(p\)-adic closure of this action gives
\(y(\zeta z)\in\mathcal O[[y]]\) for \(\zeta^{q-1}=1\). Set
\[
 T(z)=\frac{1}{q-1}\sum_{\zeta^{q-1}=1}\zeta^{-1}y(\zeta z).
\]
This is the part of \(y\) in degrees \(1\bmod(q-1)\).
Since \(q-1\) is a unit and \(T=y+O(y^2)\), the coefficients of the inverse
substitution can be solved for recursively in \(\mathcal O\).
Thus \(\mathcal O[[T]]=\mathcal O[[y]]\).
The relation \(T(\zeta z)=\zeta T(z)\) gives
\[
 T(\pi_p z)=\pi_p T(z)+\sum_{i\ge1}b_i T(z)^{1+i(q-1)}
 \qquad (b_i\in\mathcal O).
\]
The order of vanishing of this series in the residue field equals the
inseparable degree of \([\pi_p]\), which is at most its degree \(q\).
The linear term reduces to zero, and there are no terms of degrees
\(2,\ldots,q-1\). Hence \(b_1\) is a unit.
\leansrc{Lib/CM/FormalGroup/InseparableDegree}{Reduction and inseparable degree of the CM action}.

Choose \(\varpi^{q-1}=-\pi_p/b_1\) and put \(Y(z)=\varpi^{-1}T(\varpi z)\).
Then \(v_p(\varpi)=\rho_p\), and substitution in the preceding identity gives
\[
 \frac{Y(\pi_p z)}{\pi_p}=Y(z)-Y(z)^q+\pi_p H_1(Y(z)),
 \qquad H_1\in z^2\mathcal O[[z]].
\]
Indeed, the coefficient of the term indexed by \(i\ge2\) is
\(b_i\varpi^{i(q-1)}/\pi_p\in\pi_p\mathcal O\).
We have \(Y(0)=0\) and \([z]Y=1\). In degree \(n\ge2\),
\((\pi_p^{n-1}-1)[z^n]Y\) is an integral polynomial in the earlier
coefficients. Since \(\pi_p^{n-1}-1\) is a unit, induction yields
\(Y\in\mathcal O[[z]]\). Reduction in the ring \(\mathcal O/(\pi_p)\) of
characteristic \(p\) gives \(Y-Y^q=z\). There is a unique solution with
zero constant term, determined coefficient by coefficient, and it is
\begin{equation}\label{eq:upper-formalparam}
 Y(z)\equiv\sum_{\nu\ge0}z^{q^\nu}\pmod{\pi_p}.
\end{equation}
For any \(H\in\mathcal O[[y]]=\mathcal O[[T]]\), substituting
\(T(\varpi z)=\varpi Y(z)\) shows that \(H(\varpi z)\in\mathcal O[[z]]\).
Even before rescaling, the coefficients of \(H\) in degrees below \(q\)
are integral, since \(T=z+O(z^q)\).
\leansrc{Lib/CM/FormalGroup/FormalParameter}{Construction and congruence of the formal parameter}.

\emph{The exponent in one variable.}
Put \(g=-\log\theta\). Pullback by \([\pi_p]\) makes \(\theta\) an even
section of degree \(q\) and preserves integrality at torsion points.
The expansion of even sections gives
\(\theta(\pi_p z)=\theta(z)^qR(x_{\rm CM}(z))\) with \(2\deg R<q\),
where the inequality is strict because \(q\) is odd.
\leansrc{Lib/CM/Integrality/IntegralSections}{Polynomial expansion of even sections}.
This identity gives
\begin{equation}\label{eq:upper-frobenius}
 q g(z)-g(\pi_p z)=\pi_p H_2,\qquad H_2\in\mathcal O[[y]].
\end{equation}
To see this, write \(g_n=[z^n]g\). We have \(g_0=g_1=0\) and
\(\pi_p g_n\in\mathcal O\) for \(n<q\).
For \(n=2\), this follows from the integrality of \(\sqrt{-d}\,g_2\)
and \(v_p(\sqrt{-d})\le v_p(\pi_p)\).
\leansrc{Lib/CM/Integrality/LogTheta}{Integrality of the quadratic coefficient of the theta logarithm}.
For \(n>2\), the identity \(g''=g''(0)+x_{\rm CM}\) gives
\(n(n-1)g_n=[z^{n-2}]x_{\rm CM}\). The right-hand side is integral
if \(n<q\), and then \(v_p(n(n-1))\le1\) in the inert case and
\(v_p(n(n-1))=0\) in the ramified case.
Since \(q=\pi_p^2\) or \(q=-\pi_p^2\),
\((q-\pi_p^n)g_n\in\pi_p\mathcal O\) in degrees below \(q\).

It follows that all coefficients of \(\exp(qg-g(\pi_p z))-1\) in degrees
below \(q\) are divisible by \(\pi_p\). Indeed,
\(v_p(j!)\le(j-1)/(p-1)\le (j-1)v_p(\pi_p)\) for \(j\ge1\), and the
exponent has constant term zero. The exponential is \(R(x_{\rm CM})\).
The first term of \(x_{\rm CM}^i\) is \(z^{2i}\), so comparison of the
coefficients in degree \(2i\), starting at \(i=0\), yields
\(R-1\in\pi_p\mathcal O[X]\). Hence \(R(x_{\rm CM})=1+\pi_p H\) with
\(H\in\mathcal O[[y]]\) and \(H(0)=0\).
The logarithm belongs to \(\pi_p\mathcal O[[y]]\), since its \(j\)-th term
has scalar \(\pi_p^j/j\in\pi_p\mathcal O\).
\leansrc{New/Derivatives/FormalGroup/FrobeniusCongruence}{Frobenius congruence for the theta logarithm}.

Let \(g^\flat\) be the part in degrees \(p+1\bmod(q-1)\).
Write \(i_0\) for the least nonnegative representative of this class.
It equals \(p+1\) in the inert case and \(2\) in the ramified case.
The integrality of the theta transformation under the CM action gives
\[
 \exp\bigl(\zeta^{-(p+1)}g(\zeta z)-g(z)\bigr)\in\mathcal O[[y]]
 \qquad(\zeta^{q-1}=1).
\]
This follows by passing to the limit of the transformation at CM integers
and their norms, using the pair \((\zeta^{-1},\zeta^{-(p+1)})\), and
replacing the variable by \(\zeta z\).
\leansrc{Lib/CM/FormalGroup/FormalParameter}{Theta transformation at roots of unity}.
Multiplying these identities over all \(\zeta\) gives
\(\exp(g^\flat-g)^{q-1}\in\mathcal O[[y]]\) by the averaging formula
for this degree class. The \((q-1)\)-st root with constant term \(1\)
is obtained recursively by dividing the coefficient at each new degree
by \(q-1\). Hence \(\exp(g^\flat-g)\) itself belongs to \(\mathcal O[[y]]\).
With \(c_\varpi=\varpi^{p+1}/p\), this gives
\begin{equation}\label{eq:upper-one-variable}
 \begin{aligned}
 \theta(\varpi z)&=W_1(z)\exp(c_\varpi\Xi(z)),\\
 W_1(z)&=\exp\bigl((g^\flat-g)(\varpi z)\bigr)\in\mathcal O[[z]],\\
 \Xi(z)&=-g^\flat(\varpi z)/c_\varpi .
 \end{aligned}
\end{equation}

Take the part of \eqref{eq:upper-frobenius} in degrees
\(i_0\bmod(q-1)\). Since \(T(\zeta z)=\zeta T(z)\), one may take the
corresponding degree class directly in the expression for \(H_2\) in
\(\mathcal O[[T]]\). This part is
\(\lambda' T^{i_0}+T^{i_0+q-1}H_3\) with
\(\lambda'\in\mathcal O\) and \(H_3\in\mathcal O[[T]]\).
Replace \(z\) by \(\varpi z\) and divide by \(-q c_\varpi\).
The scalar \(-\pi_p\varpi^{i_0}/(q c_\varpi)\) is a unit, and the
second term retains a factor \(\varpi^{q-1}\in\pi_p\mathcal O\).
For some \(\lambda\in\mathcal O\) we thus obtain
\[
 (1-\pi_p^n/q)[z^n]\Xi\equiv\lambda[z^n]Y^{i_0}\pmod{\pi_p\mathcal O}.
\]
The coefficients for \(n<i_0\) vanish. If \(n>2\),
\(1-\pi_p^n/q\) is a unit congruent to \(1\) modulo \(\pi_p\).
For \(n=2\) in the ramified case, this factor equals \(2\).
Thus \(\Xi\) is integral, with reduction
\[
 \Xi\equiv
 \begin{cases}
    \lambda Y^{p+1}& (p\text{ inert}),\\
    \lambda(Y^2-z^2/2)& (p=d)
 \end{cases}
 \pmod{\pi_p}.
\]
In the inert case, expanding \(Y^pY\) using \eqref{eq:upper-formalparam}
gives exponents \(p^{2i+1}+p^{2j}\). The two exponents of \(p\) have
positive odd difference. Each pair consisting of the smaller exponent
\(\nu\) and the difference \(e\) occurs once. In the ramified case,
the diagonal terms of \(Y^2\) have coefficient \(1\), and the terms
with distinct exponents have coefficient \(2\). The correction changes
the coefficient of \(z^2\) to \(1/2\).
This gives
\[
  \Xi(z)\equiv
  \sum_{e\in I_p^0}\sum_{\nu\ge0}\lambda_{e,\nu}(z^{p^e}z)^{p^\nu}
  \pmod{\pi_p}.
\]
In the inert case all coefficients \(\lambda_{e,\nu}\) equal \(\lambda\).
In the ramified case \(\lambda_{0,0}=\lambda/2\),
\(\lambda_{0,\nu}=\lambda\) for \(\nu>0\), and \(\lambda_{e,\nu}=2\lambda\)
for \(e>0\).
\leansrc{New/Derivatives/FormalGroup/Exponent}{Construction and coefficient congruence of the one-variable exponent}.

\emph{The multivariable theta product.}
Suppose that the vectors \(u_i\) are integral at all primes above \(p\)
and that the multiplicities \(\mu_i\in\mathbb Z_{\ge0}\) satisfy
\(\sum_i\mu_i u_i u_i^\ast=1\).
Write \(L_i(z)=u_i^\ast z=\sum_{\tau,c}\overline{\tau(u_{i,c})}z_{\tau,c}\)
for the full linear form. Let \(F\) be the product of
\(\prod_i\theta(L_i(z))^{\mu_i}\) and an integral power series in
the \(y(L_i(z))\). Substitute each \(L_i\) in \eqref{eq:upper-one-variable}.
To simplify the exponent, use
\(\overline{\tau(u_{i,c})}^{\,p^e}\equiv\tau_e(u_{i,c})\pmod{\pi_p}\).
For algebraic integers, this congruence follows by comparing Frobenius
and the partner action on the generating root of unity and on
\(\sqrt{-d}\). At an inert prime, \(e\) is odd, so conjugation and the
\(p^e\)-th power act with the same sign on \(\sqrt{-d}\).
At a ramified prime \(\sqrt{-d}\equiv0\pmod{\pi_p}\).
For a vector entry integral above \(p\), one can clear denominators by a
rational integer prime to \(p\), whose \(p^e\)-th power is congruent to
itself. Thus
\[
 \begin{aligned}
 \sum_i\mu_i L_i(z)^{p^e}L_i(z)
 &\equiv \sum_{\tau,c}z_{\tau,c}^{p^e}
              \sum_i\mu_i\tau_e(u_{i,c}) L_i(z)\\
 &=\sum_{\tau,c} z_{\tau,c}^{p^e}z_{\tau_e,c}\pmod{\pi_p}.
 \end{aligned}
\]
The last equality is the \((\tau_e,c)\) coordinate of the resolution
of the identity \(\sum_i\mu_i u_i u_i^\ast=1\).
Moreover, taking \(p^\nu\)-th powers commutes with summation after
reduction, and \(\mu_i^{p^\nu}\equiv\mu_i\).
Consequently \(\sum_i\mu_i\Xi(L_i(z))\) differs from the sum in the
following exponent by an element of \(\pi_p\mathcal O[[z]]\) with zero
constant term.
\begin{equation}\label{eq:upper-normal}
 F(\varpi z)=W(z)\exp\!\left(
       \frac{\varpi^{p+1}}{p}
       \sum_{e\in I_p^0}\sum_{\nu\ge0}\lambda_{e,\nu}\Gamma_e(z)^{p^\nu}
                              \right),
 \qquad
 \Gamma_e(z)=\sum_{\tau,c}z_{\tau,c}^{p^e}z_{\tau_e,c}.
\end{equation}
Here \(W\in\mathcal O[[z]]\). Indeed, the exponential of the error
multiplied by \(c_\varpi\) is integral because
\(v_p(c_\varpi\pi_p)=1/(p-1)\) and \(v_p(j!)\le j/(p-1)\).
Multiplying this factor by \(\prod_i W_1(L_i(z))^{\mu_i}\)
and the integral series in the \(y(L_i(z))\) rescaled at \(\varpi z\)
gives \(W\). Each coefficient of the exponent receives only finitely
many contributions.\leanref{exists_hasNormalForm}

\paragraph{Differentiation along products of minors}
For a polynomial \(Q\), write \(J_Q(F)=(Q(\partial)F)(0)\).
Suppose \(Q\) is homogeneous of degree \(B\) with integral coefficients and
that all coefficients of \(G\) of degree \(B\) are integral.
If there are \(w\) variables, then
\begin{equation}\label{eq:upper-legendre}
  v_p(J_Q(G))\ge
  \frac{B}{p-1}-w\bigl(1+\lfloor\log_p\max(B,1)\rfloor\bigr).
\end{equation}
Indeed, the derivative for a monomial of degree \(|D|=B\) multiplies
its coefficient by \(D!=\prod_c D_c!\). Legendre's formula gives
\(v_p(D_c!)=(D_c-s_p^{\mathrm{dig}}(D_c))/(p-1)\), where
\(s_p^{\mathrm{dig}}\) is the sum of the base-\(p\) digits.
Each digit is at most \(p-1\) and there are at most
\(1+\lfloor\log_p\max(B,1)\rfloor\) digits. Apply the ultrametric
triangle inequality when summing the terms.

\emph{From a Schur matrix entry to a product of minors.}
Write \(n=n_j,\ n'=n_{j+1}\) for the source and target ranks in one step.
The diagram at each place \(\tau\) has no column of height strictly between
\(n\) and \(n'\). Fix a standard filling \(I'\) at the target. Let
\(Z_\tau=(z_{\tau,b,a})\) be an \(n'\)-by-\(n\) matrix of variables,
and \(Y_\tau=(y_{\tau,a,k})\) a matrix of integral directions.
Let \(Q_Y(Z)\) be the product of minors of \(Z_\tau Y_\tau\) obtained
by taking, for each column of height \(h<n'\) at \(\tau\), the \(h\) rows
specified by \(I'\) and the first \(h\) columns.
Write \(P_{I'I}(Z)\) for the \((I',I)\) entry of the Schur matrix.
Its defining expansion in the standard fillings gives
\[
 Q_Y(Z)=\sum_I B_I(Y)P_{I'I}(Z).
\]
Here \(B_I(Y)\) is the product of minors for the source filling \(I\)
times a common determinant factor from columns of height equal to
the source rank when \(n<n'\). The determinant factor is a nonzero
polynomial over the residue field, and the standard-filling basis
is linearly independent over that field. Hence the reductions of the
\(B_I\) are linearly independent.
\leansrc{New/Derivatives/Divisibility/Count}{Expansion of the Schur matrix in the standard-filling basis}.

Put \(\alpha_I=J_{P_{I'I}}(F)\). If all \(\alpha_I\) vanish, the bound is
immediate. Otherwise choose \(|\alpha_{I_0}|_p\) maximal. The reduction
of \(\sum_I(\alpha_I/\alpha_{I_0})B_I\) is a nonzero polynomial.
The residue field is infinite, so choose a point at which the reduction
is nonzero and lift its coordinates to \(\mathcal O\).
At the resulting \(Y\), the norm of
\(J_{Q_Y}(F)=\sum_I\alpha_I B_I(Y)\) is \(|\alpha_{I_0}|_p\).
A lower bound on \(v_p(J_{Q_Y}(F))\) for all integral \(Y\) thus
gives the same lower bound for every Schur derivative entry.\leanref{exists_norm_schurJet_le_norm_jet_minorProd}

\emph{Congruences on the directions.}
Let \(N_{\tau h}\) be the number of columns of height \(h<n'\).
Write \(t_\tau\) for the number of columns of height \(1\),
\(r_\tau\) for the number of the remaining columns, and \(B_\tau\) for
their number of boxes. The total number of boxes is
\(B=\sum_\tau(t_\tau+B_\tau)\). Let \(A\) be the set of marked places.
In counting congruences, put \(t_\tau(A)=0\) at \(\tau\in A\) and
\(t_\tau(A)=t_\tau\) elsewhere.
At each \(\tau,h\), an \(h\)-by-\(h\) matrix \(g_{\tau h}\) changes the
directions to
\[
 \beta_{\tau h k,a}=\sum_{\ell<h}(g_{\tau h})_{k\ell}\,y_{\tau,a,\ell}.
\]
For \(e\in I_p^0\) with \(e\le\lfloor\log_p\max(B,1)\rfloor\) and
\(\tau'=\tau_e\), the required congruence is
\begin{equation}\label{eq:upper-direction-pairing}
 \sum_a\beta_{\tau h k,a}^{\,p^e}\beta_{\tau' h' k',a}\equiv 0
   \pmod{\pi_p}.
\end{equation}
A pair of heights is exempt when \(h=1\) and either \(\tau\in A\)
or both \(h'=1\) and \(\tau'\in A\).

The conditions on the rows of \(g_{\tau h}\) are linear.
Choose \(\alpha_e^{p^e}=\pi_p\). Use modulus \(\alpha_e\) on the left
and modulus \(\pi_p\) on the right. If every entry of a row on the
left is divisible by \(\alpha_e\), each component of the new direction
has \(p^e\)-th power divisible by \(\pi_p\).
On the right, require each entry of the row to be divisible by \(\pi_p\).
If the height on the left is \(1\), its new direction is a scalar
multiple of the original direction \(y_{\tau,\cdot,0}\).
For each row on the right it suffices to impose one congruence,
\(\sum_a y_{\tau,a,0}^{p^e}\beta_{\tau' h' k',a}\equiv0\pmod{\pi_p}\).
If the height on the right is \(1\), choose a \(p^e\)-th root
\(\eta_a\in\mathcal O\) of each \(y_{\tau',a,0}\) and impose
\(\sum_a\eta_a\beta_{\tau h k,a}\equiv0\pmod{\alpha_e}\).
Raising to the \(p^e\)-th power gives
\(\sum_a y_{\tau',a,0}\beta_{\tau h k,a}^{p^e}\equiv0\pmod{\pi_p}\),
since \(\pi_p\mid p\) makes the \(p^e\)-th power additive after reduction.
Thus a pairing condition uses only the original direction on the
opposite side.

Allocate the conditions as follows.

\begin{enumerate}
\item If the height on the left is \(1\), impose entrywise conditions
at height \(1\) on the left, or impose pairing conditions on the right
against the original direction on the left. The respective costs are
\(t_\tau(A)/p^e\) and \(r_{\tau'}+t_{\tau'}(A)\). Choose the smaller cost.
Omit marked columns of height \(1\) on the chosen side.
\item If the height on the left exceeds \(1\) and the height on the
right is \(1\), impose pairing conditions on the left against the
original direction on the right, at cost \(r_\tau/p^e\).
\item If both heights exceed \(1\), impose entrywise conditions on the
relevant rows on one side, taking the smaller of the costs
\(B_\tau/p^e\) and \(B_{\tau'}\).
\end{enumerate}

The unit of cost is \(v_p(\pi_p)\). At each height a pairing condition
is one linear form, and an entrywise condition gives as many linear
forms as the height. The modulus on the left contributes
\(v_p(\alpha_e)=v_p(\pi_p)/p^e\), accounting for the factor \(p^{-e}\).
The factor given by the height counts the linear forms for the
entrywise conditions. Requiring divisibility of one form on every
row costs only one copy of its modulus in the determinant.

Indeed, let \(a_k\) be the values of a linear form on the rows.
If some value is nonzero, choose \(a_{k_0}\) with minimal valuation
and subtract \(a_k/a_{k_0}\) times row \(k_0\) from each other row.
The quotient is integral, and the values on those rows become zero.
Multiply the remaining row by the modulus. This multiplies the
determinant by that modulus. If all values are zero, the identity
operation suffices. Starting from the identity matrix, impose the
conditions one form at a time. The new rows are integral linear
combinations of the preceding ones, so every earlier divisibility
condition is preserved. Thus \(\det g_{\tau h}\) divides the product
of the moduli imposed at its place and height. The same matrix occurs
once for each column of that height, so
\[
 \begin{aligned}
  \sum_{\tau}\sum_{h<n'}N_{\tau h}v_p(\det g_{\tau h})
   &\le v_p(\pi_p)\mathcal D_p(A),\\
  \mathcal D_p(A)
   &=\sum_{\substack{e\in I_p^0\\e\le\lfloor\log_p\max(B,1)\rfloor}}
       \sum_{\tau\in\Phi}
       \biggl\{\min\!\left(\frac{t_\tau(A)}{p^e},\,r_{\tau_e}+t_{\tau_e}(A)\right)\\
   &\hphantom{=\sum_e\sum_\tau\biggl\{}{}
          +\frac{r_\tau}{p^e}
          +\min\!\left(\frac{B_\tau}{p^e},\,B_{\tau_e}\right)\biggr\}.
 \end{aligned}
\]
This establishes \eqref{eq:upper-direction-pairing} for all the pairs
that are subject to it, simultaneously.\leanref{exists_newDirection_dvd}

\emph{The valuation bound by the chain rule.}
Introduce variables \(s_{\tau h k b}\) and put
\[
 z_{\tau,b,a}(s)=\sum_{h<n'}\sum_{k<h}s_{\tau h k b}\,p^{w_{\tau h b}}\beta_{\tau h k,a}.
\]
When no places are marked, put \(w_{\tau h b}=0\).
When places are marked, choose nonnegative integers \(w_{\tau h b}\)
so that at least one of \(w_{\tau 1 b}\) and \(w_{\tau' h' b}\) is
positive for each exempt pair and each row \(b\).
Then \(\Gamma_e(z(s))\) is divisible by \(\pi_p\). Indeed, after
distributing the \(p^e\)-th power over the sum in the quotient ring,
each term has coefficient equal to the pairing in
\eqref{eq:upper-direction-pairing} times
\(p^{p^e w_{\tau h b}+w_{\tau' h' b}}\).
The exempt pairs also vanish by this factor.

Let \(Q_*(s)\) be the product of the determinants
\(\det(s_{\tau h k b})\) on the rows \(b\) chosen by \(I'\) and
directions \(k<h\) for each column. Let \(W_{\rm ch}\) be the sum of
\(w_{\tau h b}\) over the boxes, using the row index \(b\) of each box.
The chain rule gives
\begin{equation}\label{eq:upper-minor-chain-rule}
 J_{Q_*}\bigl(F(\varpi z(s))\bigr)
 = \varpi^B p^{W_{\rm ch}}\!
    \left(\prod_{\tau,h<n'}(\det g_{\tau h})^{N_{\tau h}}\right)J_{Q_Y}(F).
\end{equation}
In one column, the derivative matrix along the new directions is the
derivative matrix for \(Y_\tau\) multiplied on the right by
\(g_{\tau h}^{\,t}\) and on the left by the diagonal matrix with row
factors \(\varpi p^{w_{\tau h b}}\). Taking the product of determinants
gives \eqref{eq:upper-minor-chain-rule}.

The polynomial \(Q_*\) is integral of degree \(B\) in at most
\(f(n')^3\) variables. In \eqref{eq:upper-normal} after substitution,
the terms with \(e>\lfloor\log_p\max(B,1)\rfloor\) have degree
greater than \(B\). All remaining couplings are divisible by \(\pi_p\),
so the exponential calculation above shows that \(F(\varpi z(s))\)
has integral coefficients through degree \(B\).
Equations \eqref{eq:upper-legendre} and \eqref{eq:upper-minor-chain-rule} give
\begin{equation}\label{eq:upper-jet}
 \begin{aligned}
 v_p(J_{Q_Y}(F))+W_{\rm ch}
   &\ge \gamma_p B-v_p(\pi_p)\mathcal D_p(A)\\
   &\quad{}-f n_{j+1}^3\bigl(1+\lfloor\log_p\max(B,1)\rfloor\bigr),\\
 \gamma_p&=\frac1{p-1}-\rho_p .
 \end{aligned}.
\end{equation}
Here \(B\) includes the boxes of marked columns.\leanref{norm_jet_minorProd_le_of_charges}
For the positive exponents one has
\(\gamma_p=v_p(\pi_p)\sum_{e\in I_p}p^{-e}\).
Subtracting the deductions for \(e>0\) in \(\mathcal D_p(A)\) from
\(\gamma_p B\) gives the following contributions in units of \(v_p(\pi_p)\).
\[
 \left(\frac{t_\tau(A)}{p^e}-t_{\tau_e}(A)-r_{\tau_e}\right)_+
 +\left(\frac{B_\tau}{p^e}-B_{\tau_e}\right)_+
 -\frac{r_\tau}{p^e}.
\]
This uses \(x-\min(x,y)=(x-y)_+\) for the two minima.
Extending the deductions to all of \(I_p\) still gives a lower bound,
since the additional terms are nonnegative.
Add \(\gamma_p\) for each marked box. At \(p=d\), also subtract the
deduction for \(e=0\).

\paragraph{Gain in one step}
Define the local prime sum and the gain over the full chain by
\begin{align}
 s_p(y,\tau)
  &=v_p(\pi_p)\sum_{e\in I_p}
       \left(\frac{y}{p^e}-\kappa((-p^e)\tau)\right)_+,
       \label{eq:upper-sp}\\
 G_p
  &=m\left\{\sum_{\tau\in\Phi_a}s_p(a,\tau)
           +k\sum_{\tau\in\Phi}s_p(\kappa(\tau),\tau)\right\}.
       \label{eq:upper-gain}
\end{align}
Set \(s_p=G_p=0\) at split primes.
For \(0\le y\le a\), all the positive parts also vanish if \(p\ge x_1\).
If \(p\) is inert, \(1-\chi_K(p^e)\) is \(2\) for odd \(e\) and \(0\)
for even \(e\). If \(p=d\), this coefficient is \(1\) and
\(v_p(\pi_p)=1/2\). Grouping the terms by prime powers gives
\begin{equation}\label{eq:upper-gain-sum}
  2\sum_p G_p\log p =
    m\sum_{\tau\in\Phi_a}s_{\mathrm{fin}}(a,\tau)
       +mk\sum_{\tau\in\Phi}s_{\mathrm{fin}}(\kappa(\tau),\tau).
\end{equation}

There are \(a\) columns of height \(1\) at places in \(\Phi_a\) and none
elsewhere. The number of remaining columns is at most
\(\kappa(\tau)+x^{60}\), and the number of their boxes is between
\(k\kappa(\tau)\) and \(k\kappa(\tau)+x^{60}\).
Insert these bounds into the positive parts above.
Write \(D_{p,j}(A)\) for the valuation bound for one step, where \(A\subseteq\Phi\)
is the set of places at which the height-\(1\) columns are marked.
It is the predicted valuation of the diagram with those columns removed,
plus \(\gamma_p\) for each box removed, less the derivative error
\(f n_{j+1}^3(1+\log_p\max(B,1))\). When \(p<x_1\) is nonsplit,
\begin{equation}\label{eq:upper-one-step}
 D_{p,j}(A)\ \ge\ \frac{G_p}{m}
  - C_p(A)
  +v_p(\pi_p)|\Phi_a\cap A|\,\frac{a}{x_1}
  -\frac{fa(\log b)^2}{mp},
 \qquad
 C_p(A)=\sum_{\substack{\tau\in\Phi,\ e\in I_p\\
                 \tau,\,(-p^e)\tau\in\Phi_a\setminus A}}\frac{a}{p^e}.
\end{equation}
When a column of height \(1\) remains at the partner place, the overlap
costs at most \(a/p^e\).
The error in the other columns is controlled by \(x^{60}\) as above.
The deduction for \(p=d,e=0\) is bounded using that \(\tau,-\tau\)
cannot both lie in \(\Phi_a\) and cannot both lie in \(\Phi_b\).
In units of \(a/(mp)\), the error per place is at most \(2\) from
comparing the positive parts, \(2\) from the ramified-prime deduction,
and \(1\) from the logarithmic error in differentiation.
The condition \(5\le(\log b)^2\) absorbs these errors.
\leanref{le_divisibilityBound}

\paragraph{The local kernel and its lattice}
The local kernel of a lift \(X_j\to X_{j+1}\) uses Schwartz data on
the standard space of \(n'\)-by-\(n\) matrices.
Standard data of level \(1\) are finite \(p\)-integral combinations
of translations of the integral vacuum by points integral above \(p\).
A nonsplit prime \(p\) is prime to the cyclotomic modulus \(q_1\).
We show that the theta series of any standard datum of level \(1\)
has the normal form \eqref{eq:upper-normal}.

\emph{The theta product for the vacuum.}
Choose elements \(v_i\) of the self-dual lattice \(\mathfrak a_0\) and an
integer \(b_0\) prime to \(p\) with
\(\sum_i \tau(v_i)\overline{\sigma(v_i)}=b_0^2\delta_{\tau,\sigma}\).
The tight-frame construction makes a product of copies of \(\theta\)
along the linear forms given by placing \(v_i/b_0\) in each coordinate
in turn. The associated Weil vector \(F_0\) has theta value \(1\),
and its theta function has the prescribed signed periods along
\(b_0\mathfrak a_0\) in each coordinate.
\leansrc{New/Derivatives/Vacuum/TightFrames}{Construction of the theta product from a tight frame}.

Apply finite Fourier inversion for this sublattice to express the
integral vacuum as a sum of translates of \(F_0\).
For a representative \(u\) of a residue class, the corresponding
Fourier coefficient times the vacuum equals the sum of translates
of \(F_0\), first by \(-u\) and then by the representatives of the
period lattice. Applying the theta functional to the inversion
identity shows that the sum of these Fourier coefficients weighted
by the corresponding theta values of translates of the vacuum is
\((b_0^2)^{f n n'}\). Thus the integral vacuum is a finite sum of
translates of \(F_0\). The denominator in the coefficients is
\((b_0^2)^{f n n'}\), and the numerators are products of values of the
integral vacuum and the scalar factors of the Heisenberg action.
The coefficients are therefore integral at \(p\).
The translation points lie in the standard lattice after
multiplication by \(b_0\).
\leansrc{New/Derivatives/Vacuum/StandardDataProducts}{Finite Fourier inversion for the vacuum and its denominators}.

Apply this expansion to each translate in the standard datum.
The composition law for the Heisenberg action adds the translation
points and multiplies by an integral phase.
Each resulting translate of \(F_0\) can be represented by a
rational point \(\gamma\) inducing the same action on \(F_0\),
with \(a_0\gamma\) in the lattice in each coordinate for some
integer \(a_0\) prime to \(p\).
\leansrc{New/Derivatives/Vacuum/StandardDataProducts}{Choosing rational translation representatives}.
Each shift \(t_i\) in the theta product for \(\gamma\) satisfies
\(a_0 b_0 t_i\in\Omega_0\mathcal O_K\).
Indeed, the shift is \(-\Omega_0/b_0\) times the Hermitian pairing
between \(v_i\) and the corresponding coordinate of \(\gamma\).
The pairing of \(v_i\) and \(a_0\gamma\) lies in \(\mathcal O_K\)
by self-duality. Choose a common multiple \(k_0\) of the integers
\(a_0 b_0\) for all these translates, still prime to \(p\).

Place \(v_i/(b_0 k_0)\) in each coordinate in turn to obtain
the vectors \(u_i\) with multiplicity \(k_0^2\). Dividing the
frame identity by \(b_0^2\) gives \(\sum_i k_0^2u_i u_i^\ast=1\).
Since \(k_0 t_i\) is a period, the one-variable translated theta
series after pullback by \([k_0]\) is the product of \(\theta(z)^{k_0^2}\)
and an element of \(\mathcal O[[y]]\). Here the pulled-back
section has degree \(k_0^2\).
\leansrc{Lib/CM/Integrality/IsogenyIntegrality}{Pullback of the theta series translated by a torsion point}.
Since \(k_0\) was chosen simultaneously for all terms, the finite sum
has the common factor \(\prod_i\theta(u_i^\ast z)^{k_0^2}\).
The remaining factor lies in the ring of integral series in the
\(y(u_i^\ast z)\).

The construction first gives integrality of the adelic image
\(\alpha^{-1}u_i\). The idele \(\alpha\) that makes the standard
lattice self-dual satisfies \(q_1\alpha\) integral, and is therefore
integral at the places above \(p\) since \(p\nmid q_1\).
Thus \(u_i=\alpha(\alpha^{-1}u_i)\) is integral at all places above \(p\).
The expansion applies to the theta function with variable scaled by
\(\sqrt{2\pi/\sqrt d}/\Omega_0\). The kernel uses an additional
period factor and a factor \(\pi_M\), with \(\pi_M\overline{\pi_M}=M\).
The period factor is a unit at \(p\) since \(p\nmid q_1\).
If \(c\) is the product of the two additional factors, \(v_p(c)\ge0\).
Replacing \(z\) by \(cz\) in \eqref{eq:upper-normal} preserves the
integrality of \(W(cz)\) and multiplies the exponent coefficients by
\(c^{(p^e+1)p^\nu}\). The kernel therefore has the same normal form
at its prescribed scale.
\leansrc{New/Derivatives/Vacuum/Divisibility}{Normalization and integrality of the theta series of the local kernel}.

\emph{The bound in lattice coordinates.}
Use the standard-filling basis of each Weyl module.
The local norm \(\|\eta\|_g\) of a vector \(\eta\) at a point \(g\)
is the maximum \(p\)-adic absolute value of its coordinates after
acting by \(g^{-1}\) in the weight representation.
The weight action at \(\tau\) is read through the embedding
\(x\mapsto\iota_p(\tau(\overline x))\). The norm is unchanged by
right multiplication of the point by a lattice-preserving element,
or by the simultaneous action of a rational point on \(g\) and \(\eta\).

For a standard datum \(\psi\), write \(K_j(\psi;g,h)\) for the
matrix of Schur derivatives of \(\omega_j(g,h)\psi\) under the
twisted Weil action. Here \(g\) and \(h\) are finite adelic points
of \(X_j\) and \(X_{j+1}\), respectively.
Consider also replacing \(\psi\) by \(\omega_j(1,x d_\ast y)\psi\),
where \(x,y\) are integral above \(p\) and \(d_\ast\) is read at
each weight embedding as a diagonal matrix with entries \(u_{\tau,b}\).
Weak approximation gives rational points \(\gamma,\epsilon\) such
that \(g_0=\gamma g,\ h_0=\epsilon h x d_\ast\), and
\(d_\ast h_0 d_\ast^{-1}\) are integral above \(p\).
The points \(g_0,h_0 y\) preserve the standard lattices and their
Weil action preserves the span of integral translates of the vacuum.
The character factors of the twist are units at these points.
Thus \(\psi'=\omega_j(g_0,h_0 y)\psi\) is standard.

By covariance of the kernel and invariance of the local norms,
the comparison at the input point \(g\) and output point \(h\)
reduces to the derivative matrix of the standard datum \(\psi'\)
at the identity, with input norm at \(1\) and output norm at
\(d_\ast^{-1}\). Indeed, \(g_0\) is integral and the last two
factors of \(\epsilon h=d_\ast^{-1}(d_\ast h_0d_\ast^{-1})x^{-1}\)
are integral. The row indexed by \(I'\) is now multiplied in
absolute value by \(|u_{\tau,b}|_p\) for each box whose entry at
\(\tau\) is \(b\), and by the determinant factor
\(\prod_\tau(\prod_b|u_{\tau,b}|_p)^{e'_\tau}\), where \(e'_\tau\)
is the last entry of the target weight.
The operator norm is bounded by the largest absolute value of an
entry, so it suffices to bound these factors times the corresponding
Schur derivative.
\leansrc{New/UpperBound/Kernels/ThetaLiftGauge}{Twisted Weil action and the kernel bound in lattice coordinates}.

For unmarked standard data, take \(d_\ast=1\), so the row factors
are all \(1\). In \eqref{eq:upper-jet}, put \(W_{\rm ch}=0\),
extend the deductions to all \(e\), and relax the error term by
replacing the floor of the logarithm by the logarithm.
Each entry has valuation at least \(D_{p,j}(\varnothing)\).
Consequently
\begin{equation}\label{eq:upper-kernel-gauge}
 \|K_j(\psi;g,h)\eta\|_h
 \le p^{-D_{p,j}(\varnothing)}\|\eta\|_g
\end{equation}
for all \(g,h,\eta\).
\leansrc{New/UpperBound/Divisibility/ValuationAtNonSplitPrimes}{Valuation of the kernel for one lift}.

\emph{Adding valuations through the chain.}
Let \(\psi_{t,j}\) be the datum in step \(j\), and suppose its kernel
multiplies local norms by at most \(p^{c_j}\), uniformly for
\(t\) in the support of \(\Psi\).
Joint invariance at the intermediate levels expresses the value
of the composite lift as the finite sum, over \(t\) and class
representatives \(g_j\), of products
\[
  \Psi(t)\,K_{m-1}(\psi_{t,m-1};g_{m-1},g_m)
        \cdots K_0(\psi_{t,0};g_0,g_1)\phi(g_0).
\]
This expression has no denominator from the order of a class stabilizer.
The output lattice at \(g_{j+1}\) is the input lattice of the next
kernel, so the operator bounds multiply. If \(\phi\) has local norm
at most \(1\) and the coefficients \(\Psi(t)\) are \(p\)-integral,
the ultrametric triangle inequality bounds the sum by
\(p^{\sum_{j<m}c_j}\). Applying linear operators that commute with
the level actions to the data preserves joint invariance.
The final restriction to \(Y_m\subset X_m\) uses an integral branching
matrix and the \(p\)-adic unit \(\overline{\pi_N}\) when \(p\ne N\).
It preserves the bound.
\leansrc{New/UpperBound/Divisibility/ChainValuation}{Finite class sum and restriction bounds for the chain}.

\subsection{Colength at the finite places}

Let \(\iota\) be the isomorphism defining a valuation above \(p\), and let
\(\mathcal O_\iota\) be its valuation ring. Write \(M_\iota\) for the
module of \(\iota\)-integral forms in \(S_0\).
We say that a submodule has colength dividing \(p^j\) if for every finite
family in \(M_\iota\) a square matrix over \(\mathcal O_\iota\) sends that
family into the submodule and has determinant dividing \(p^j\).
In particular, the valuation of the determinant is at most \(j\).
This definition works directly with finite families in the module.

Let \(\mathcal U\) be the set of nonsplit primes below \(x_1\), and put
\[
 \ell_p=\frac{fa(\log b)^2}{p}
      \begin{cases} m & (p\equiv\pm1\bmod m),\\1 & (\text{otherwise}) \end{cases},
 \qquad
 \varepsilon_p=
      \begin{cases}
       2\ell_p &(p\in\mathcal U),\\
       2mfb &(p=N),\\
       0 &(\text{otherwise}).
      \end{cases}
\]
Here \(N\ge x_1\). For each \(p,\iota\), we will show that there is a
submodule \(M'_\iota\subseteq M_\iota\cap S\) of colength dividing
\(p^{j_\iota}\) in \(M_\iota\cap S\), with
\begin{equation}\label{eq:upper-local}
 j_\iota\le\varepsilon_p\dim S_0,\qquad
 v_p(B_{\xi,\Psi}\phi)\ge G_p-\varepsilon_p
  \quad(\phi\in M'_\iota,\ (p,\mathrm NQ(\xi))=1).
\end{equation}
For a form, a lower valuation bound \(v\) means that all its values have
\(p\)-adic norm at most \(p^{-v}\) in the corresponding lattice coordinates.
Here \(\Psi\) is any datum satisfying the conditions imposed in
\eqref{eq:upper-seesaw}. The choice of \(M'_\iota\) is independent of
\(\xi,\Psi\).

For \(p\notin\mathcal U\cup\{N\}\), one has \(G_p=0\).
Since \(d<x_1\), we also have \(p\ne d\), so the kernel at each step
relative to the standard lattice and the final restriction both preserve
integrality.
For \(p\in\mathcal U\), first apply \eqref{eq:upper-one-step} at every
step with no marks. The transformation laws of the kernels allow the
lower valuation bounds to be added along the chain. The coefficients
of the composite lift are integral, so its finite sum preserves the
resulting bound. This gives \(G_p-\ell_p-mC_p(\varnothing)\) on \(M_\iota\).
If every \(e\in I_p\) with \(\tau,(-p^e)\tau\in\Phi_a\) satisfies \(e\ge e_0\),
the geometric series gives \(C_p(\varnothing)\le 2fa/p^{e_0}\).
When \(p\equiv\pm1\bmod m\), take \(e_0=1\).
In the other case, a collision can occur only for inert \(p\), and then
\(e\) is an odd multiple of \(o_0\), the order of the place label of
\(-p\) modulo \(q_0\). If \(o_0>\lceil\log_p m\rceil\), one can take
\(p^{e_0}\ge pm\). In these cases \(mC_p(\varnothing)\le\ell_p\)
and \(M'_\iota=M_\iota\cap S\) suffices.

The remaining exceptional case is when \(p\) is inert,
\(p\not\equiv\pm1\bmod m\), and \(o_0\) is odd with
\(o_0\le\lceil\log_p m\rceil\).
Among the finite places of \(E\) induced by \(\Phi_a\) and their conjugates,
one can select one place from each conjugate pair to obtain a full
marking \(\Lambda\). If \(A\) is the set of infinite places inducing
places in \(\Lambda\), then \(C_p(A)=0\).
\leanref{exists_isFullMark}

Put \(r_p=\lceil\log_p m\rceil\) and \(\nu_p=r_pf(n_m+r_p)\).
For \(w\in\Lambda\), let \(a_w\) be \(a\) times the number of members of
\(\Phi_a\) inducing \(w\).
Let \(Y\) be the product over \(w\in\Lambda\) of the operators \(T_w\)
obtained by multiplying the Hecke operator of index \(r_p\) at \(w\)
by \(p^{a_w}\).
\emph{The marked estimate on neighboring lattices.}
Consider translating the kernel on the right at each \(w\in\Lambda\)
by an element of a Hecke double coset of index at most \(r_p\).
Up to integral factors on both sides, its component at \(w\)
is diagonal with \(p\) on a set of rows \(J_w\), \(|J_w|\le r_p\),
and \(1\) elsewhere. The conjugate component is the inverse.
Since the weight embedding reads the conjugate component, the row
factor has absolute value \(p^{\eta_{\tau,b}}\).
Here \(\eta_{\tau,b}\) is \(1\) if \(\tau\) induces \(w\in\Lambda\)
and \(b\in J_w\), is \(-1\) if \(\tau\) induces \(\overline w\)
and \(b\in J_w\), and is otherwise \(0\).
The determinant factor has exponent \(\sum_\tau e'_\tau\sum_b\eta_{\tau,b}\),
which vanishes by the equality of the sums of the last weight
entries above \(w\) and \(\overline w\).

For \eqref{eq:upper-jet}, take \(w_{\tau h b}=1\) if \(\tau\) induces
\(w\) and \(h=1,\ b\notin J_w\), or if \(\tau\) induces \(\overline w\)
and \(b\in J_w\). Take it to be \(0\) elsewhere.
The prime \(p\) is inert, and each relevant \(e\) is odd, so
\(\tau_e\) induces the conjugate of the prime induced by \(\tau\).
In each exempt pair, if \(b\in J_w\) then the conjugate side
receives a factor of \(p\). If \(b\notin J_w\), the marked column
of height \(1\) receives that factor. Thus the terms exempted from
\eqref{eq:upper-direction-pairing} also have the required divisibility.

The sum \(\eta_{\tau,b}+w_{\tau h b}\) for a box is zero at a place
inducing \(\overline w\). At a place inducing \(w\), a box in a
column of height \(1\) contributes \(1\). In any other column,
the row indices in a standard filling are distinct, so there are
at most \(r_p\) boxes with index in \(J_w\).
The sum over boxes is therefore at most
\(\sum_{w\in\Lambda}a_w+r_p\sum_{\tau\in A}r_\tau\).
Subtracting \(D_{p,j}(A)\) gives the upper bound on the loss
for the kernel on the neighboring lattices.
\leansrc{New/UpperBound/Divisibility/ValuationAtNonSplitPrimes}{Weight factors and powers of \(p\) for the neighboring lattices}.

\emph{Transferring the Hecke operators into the data.}
Let \(C_{\boldsymbol x,j}\) be the composite of the first \(j\) lifts
with the translation tuple \(\boldsymbol x\). For \(w\in\Lambda\),
write \(Q_w=z_w^2\) for the order of its residue field, with
\(z_w=p^{h_w}\) and \(2h_w\le f\). Denote the Hecke operator on
\(X_j\) of index \(i\) by \(T_{w,i}^{(j)}\).
Write \(t_{\rm rk}+1\) for the rank increment \(n_{j+1}-n_j\).
The character of the twist is unramified at this place, and the
local datum is the vacuum. For the Hecke polynomial
\[
 {\cal H}_{w,j}(U)=\sum_{i=0}^{n_j}(-U)^i Q_w^{\binom{i}{2}}T_{w,i}^{(j)}
\]
the identity for an unramified theta lift \(L_j\) is
\begin{equation}\label{eq:upper-hecke-polynomial}
 {\cal H}_{w,j+1}(U)L_j
    =R_j(U)L_j{\cal H}_{w,j}(z_w^{t_{\rm rk}+1}U).
\end{equation}
Here \(\chi_{w,j+1}\) is the value of the base change of the lifting
character at a uniformizer of \(w\), and
\(R_j(U)=\prod_{0\le k\le t_{\rm rk}}(1-Q_w^k z_w^{n_j}\chi_{w,j+1}U)\).
The equality of the sums of the character weights above \(w\) and
\(\overline w\) makes \(\chi_{w,j+1}\) a \(p\)-adic unit.
Hence \(R_j\in\mathcal O[U]\) and \(R_j(0)=1\).
\leansrc{Lib/Weil/Unramified/HeckeTransfer}{Hecke polynomial identity for the unramified theta lift}.

Multiply \eqref{eq:upper-hecke-polynomial} by \(R_j^{-1}\in\mathcal O[[U]]\)
and compare coefficients of \(U^i\).
Then \(z_w^{(t_{\rm rk}+1)i}Q_w^{\binom{i}{2}}L_j T_{w,i}^{(j)}\)
is an integral linear combination of \(T_{w,i'}^{(j+1)}L_j\)
for \(i'\le i\).
The prefactor is \(p^{h_w i(t_{\rm rk}+i)}\).
For \(i\le r_p\le t_{\rm rk}\), it divides \(p^{f r_p t_{\rm rk}}\),
since \(2h_w\le f\). Paying this power of \(p\) at each step thus
keeps the transfer integral and preserves the upper bound \(r_p\)
on the index. Iterate through \(j+1\) steps and multiply by the
remaining power of \(p\), using \((j+1)f r_p t_{\rm rk}\le\nu_p\).
This gives \(p\)-integral coefficients \(c_{w,i}\) with
\begin{equation}\label{eq:upper-hecke-transfer}
 \sum_{i\le r_p}c_{w,i}\,T_{w,i}^{(j+1)}C_{\boldsymbol x,j+1}
       = p^{\nu_p} C_{\boldsymbol x,j+1} T_{w,r_p}^{(0)} .
\end{equation}
The identity holds on forms spherical at \(w\), and its coefficients
are independent of \(\boldsymbol x\).
\leansrc{New/UpperBound/Divisibility/ValuationAtNonSplitPrimes}{Iterating the Hecke transfer through the lifts}.

Apply the operator on the left of \eqref{eq:upper-hecke-transfer}
to the datum and multiply by \(p^{a_w-\nu_p}\).
Do this at each of the distinct places \(w\).
Covariance moves the target Hecke action into the datum.
The datum is spherical at each \(w\), so the resulting operator
commutes with the target level and with the source action.
Writing \(L'_{\boldsymbol x,j}\) for the lift with the changed datum,
we have \(L'_{\boldsymbol x,j}C_{\boldsymbol x,j}
=C_{\boldsymbol x,j+1}Y\).
Induction gives \(C_{\boldsymbol x,m}Y^m\) for the composite of the
\(m\) changed lifts.

Each \(T_{w,i}^{(j+1)}\) is a finite sum of translations by one-sided
coset representatives in the double coset.
They all satisfy the conditions of the neighboring-lattice estimate.
The ultrametric triangle inequality preserves the bound in a
finite sum with the integral coefficients \(c_{w,i}\).
Multiplication by \(p^{a_w-\nu_p}\) adds a loss of \(\nu_p-a_w\)
for each \(w\). Hence the total loss in step \(j\) is at most
\[
  |\Lambda|\nu_p+r_p\sum_{\tau\in A}r_\tau-D_{p,j}(A).
\]
Commutation with the levels preserves joint invariance of the data
in the chain, so these losses can be summed over the steps.\leanref{gaugeLe_scaledHeckeOperator_pow}

The cost at the marked places satisfies
\[
 |\Lambda|\nu_p+r_p\sum_{\tau\in A}r_\tau
 \le \frac{\ell_p}{m}+|\Phi_a\cap A|\frac a{x_1}.
\]
\leansrc{New/UpperBound/Divisibility/MarkedPrimes}{Cost at marked places}.
At an exceptional prime \(v_p(\pi_p)=1\), so the terms involving
\(a/x_1\) cancel, and the loss per step is at most
\((2\ell_p-G_p)/m\).
Set \(j_p=\lceil 2mfa/p\rceil\).
On \(p^{j_p}M_\iota\), the valuation of the image is at least
\(G_p-\ell_p\), by the unmarked estimate and
\(mC_p(\varnothing)\le 2mfa/p\le j_p\).
Thus the bound \(G_p-2\ell_p\) holds on
\(Y^m M_\iota+p^{j_p}M_\iota\).

The normalized Hecke operator \(T_w\) preserves \(M_\iota\), and the
colength of \(T_wM_\iota+p^{j_p}M_\iota\) divides \(p^{c_w}\) for some
\[
  c_w\le\left(\frac{12a_w}{mp}+\frac{j_p}{m^2f}\right)\dim S_0 .
\]
\leansrc{New/UpperBound/Divisibility/MarkOperator}{Elementary divisors of the normalized Hecke operator}.
To see the numerical bound, the places of \(\Phi_a\) inducing
\(w\) or \(\overline w\) form one orbit under \(-p^{o_0}\).
Let its length be \(2u_w\) and put \(P_p=p^{o_0}\).
Order the embeddings of the orbit into the valuation ring, composing with
complex conjugation on the conjugate side. Successive embeddings are
related by the \(P_p\)-power map modulo \(p\). The even embeddings induce
\(w\), so \(a_w=au_w\).
The elementary-divisor theorem for a Frobenius cycle factors the Hecke operator on the lattice classes into a restriction
to residue classes and a pairing of matrices, and compares the
colengths of the images of these matrices.
For the coordinate model and numerical conditions
\[
 \begin{gathered}
  0<r_p\le k/2,\qquad a^4\le P_p^{2u_w},\qquad \theta_0=(m^2f)^{-1},\\
  (2u_w k\log(P_p^{2u_w})/\theta_0)^5 P_p^{3r_p}\le a
 \end{gathered}
\]
the bound per coordinate is \(3C_{\mathrm{div}}+\theta_0j_p\), where
\(C_{\mathrm{div}}=2u_wa(P_p^{-2r_p}+r_p^2/(kP_p))\).
\leansrc{New/UpperBound/ElementaryDivisors/ElementaryDivisorTheorem}{Elementary divisors for a Frobenius cycle}.
For the present Hecke operator, the hypotheses follow from its coordinate
model and the scale conditions.
\leansrc{New/UpperBound/Divisibility/MarkOperator}{Application of the elementary-divisor theorem}.
From \(m\le p^{r_p}\) and \(p\le P_p\) one gets
\(P_p^{-2r_p}\le1/(mp)\). From \(mr_p^2\le k\) and \(p\le P_p\) the
second term is at most \(1/(mp)\) as well.
Hence \(3C_{\mathrm{div}}\le12a_w/(mp)\).

Colength exponents add under composition of the operators.
First send a finite family into \(T_1M_\iota+p^{j_p}M_\iota\), then apply
a correcting matrix for \(T_2\) to the preimages under \(T_1\) of the
\(T_1M_\iota\) components. Since \(T_1\) preserves \(p^{j_p}M_\iota\),
the composite matrix sends the family into
\(T_1T_2M_\iota+p^{j_p}M_\iota\). Its determinant is the product of the
two determinants. Repeating this argument and using
\(\sum_w a_w\le af\), \(|\Lambda|\le f\), and \(j_p\le2mfa/p+1\), the
colength exponent of \(Y^mM_\iota+p^{j_p}M_\iota\) is bounded above by
\[
  m\sum_{w\in\Lambda}
      \left(\frac{12a_w}{mp}+\frac{j_p}{m^2f}\right)\dim S_0
    \le \frac{14fa+p/m}{p}\dim S_0
    \le \ell_p\dim S_0 .
\]
Intersecting with \(S\) preserves this bound because the original family
can be taken in \(S\) and its linear combinations then remain in \(S\).

Finally, for \(p=N\), take the submodule on which
\(\mathrm{Up}^{m/2}\phi\) is integral in the modified standard lattice.
Its colength exponent is at most \(mfb\dim S_0\).
\leansrc{New/UpperBound/Divisibility/ColengthAtLevel}{Colength at the level prime}.
The ramification condition expresses each glued pair of lifts as a
composite of \(\mathrm{Up}\) and direct lifts. On this submodule the
resulting valuation is at least \(-2mfb\).
\leansrc{New/UpperBound/Divisibility/ValuationAtLevel}{Valuation at the level prime}.
Intersecting with \(S\) preserves the colength bound.
Since \(G_N=0\), this proves \eqref{eq:upper-local} in all cases.

The local errors satisfy
\begin{equation}\label{eq:upper-error-sum}
  \sum_p \varepsilon_p\log p\le E_3.
\end{equation}
For \(p\in\mathcal U\), use \(\log p\le\log x_1\).
When \(p\equiv\pm1\bmod m\), write \(p=mc\pm1\) with \(1\le c\le p\)
and use \(m/p\le2/c\). Otherwise take \(c=p\).
For the term \(p=N\), use \(c=1\) and
\(2mfb\log N\le4fa(\log b)^2\log x_1\).
At most four terms correspond to a given \(c\), from
\(N,mc-1,c,mc+1\). Thus the sum is at most
\(16fa(\log b)^2\log x_1(1+\log x_1)\), which is bounded by \(E_3\)
by the scale conditions. \leanref{sum_allowedLoss_mul_log_le}

\subsection{Archimedean norms and normalization}

For step \(j\) of the upper chain, let \(l_{j,\tau,i}\), \(i<n_j\),
be the row lengths used at \(\tau\) and put
\(U_j=\sum_\tau\sum_{i<n_j}s_\infty(l_{j,\tau,i})\).
Recall that \(s_\infty(q)=\log(q!)+q\log(M/\Omega)\), where \(\Omega\)
is the period. The row sum \(U_j\) gives the factorial contribution
in the norm of the kernel.
Use the Gram norm on the Weyl modules and test the Schur matrix
against row vectors of dual norm at most \(1\).
Let \(\rho_{\tau,i}\ge0\) be the squared singular values of
\(Z_\tau\), in decreasing order. The squared operator norm of the Schur matrix is
at most \(\prod_{\tau,i}\rho_{\tau,i}^{l_{j,\tau,i}}\).
Indeed, the unitary factors preserve the Gram norm, and for the
diagonal factor the highest weight gives the exponents \(l_{j,\tau,i}\).
The sum of the squared singular values is the sum of squared
absolute values of the entries of \(Z_\tau\).
\leansrc{Lib/GLn/SchurModule/SchurNorm}{Singular-value bound for the Schur matrix}.

Put \(\lambda=a+1\), \(u=\log(M/\Omega)\), and \(v=u+1/\lambda\).
Taking the term of degree \(l\) in the exponential series gives
\(\rho^l\le l!\exp(lv+\rho e^{-v})\).
If \(\sum_{\tau,i}\rho_{\tau,i}\le e^u(2\lambda+2)R\), then
\(e^{-v}\sum_{\tau,i}\rho_{\tau,i}\le2\lambda R\), since
\(e^{1/\lambda}\ge1+1/\lambda\).
Thus the Schur polynomial satisfies a growth bound with factor
\[
 \exp\!\left(\frac{U_j}{2}
       +\frac1{2\lambda}\sum_{\tau,i}l_{j,\tau,i}+\lambda R\right).
\]
In the conjugated Gaussian integral, the scalar multiplying the
variables has squared modulus \(M/\Omega\).

The conjugated derivative is a Gaussian integral of the Schur
polynomial, with its places permuted, against a conjugate
translated vacuum. With \(s_j=f n_j n_{j+1}\) complex variables \(y\),
put \(R=\sum|y|^2/(2\lambda+2)\).
The factor \(\exp(\lambda R)\) from polynomial growth times
\(\exp(-\sum|y|^2/2)\) from the vacuum is \(\exp(-R)\).
The measure is Lebesgue measure divided by \(\pi\) in each complex
coordinate, so the integral of this Gaussian is \((2\lambda+2)^{s_j}\).
It suffices to use the integral representation for conjugates fixing
\(K\), since composing with complex conjugation preserves the
absolute values for the other conjugates.
\leansrc{New/UpperBound/Kernels/JetConjugatesAndBounds}{Gaussian integral and the row factorials}.

The logarithmic constant from the bound on the conjugate of
the integral vacuum is \(\mathrm{const}(P)(q_1 n_j n_{j+1}+1)^2\).
The sum of this constant, \(s_j\log(2\lambda+2)\), and
\((2\lambda)^{-1}\sum_{\tau,i}l_{j,\tau,i}\) is at most \(x^{50}\)
by the bound \(2a\) on the row lengths and the scale conditions.
Thus a single lift increases the norm by at most
\(\exp(x^{50}+U_j/2)\) for every conjugate.\leanref{sqrt_weylNormSq_jetMap_le}

The principal contributions to the row lengths are \(a\,1_{\Phi_a}(\tau)\)
in the first row and \(\kappa(\tau)\) in each of the next \(k-1\) rows.
Each additional box contributes at most \(25\log b\). Indeed, all the
row lengths are at most \(2a\) and adding a box increases the bound by at
most \(\log(2a+1)+\log(M/\Omega)\le25\log b\).
At each step the total number of additional boxes is at most
\(f(b+a/x_1+n_m r^{16})\).
The scale conditions give
\begin{equation}\label{eq:upper-rows}
 \sum_{0\le j<m}U_j
 \le m\left\{|\Phi_a|s_\infty(a)+k\sum_\tau s_\infty(\kappa(\tau))\right\}
         +\frac{E_3}{2}.
\end{equation}
\leanref{sum_rowSizeSum_le}

In the expansion of the glued tensor for the trivial twist, each step
has at most \(\exp(x^{30})\) classes, and the sum of the absolute values
of the conjugated translation coefficients is at most \(\exp(x^{60})\).
If all values of \(\phi\) have squared norm at most \(\exp(B)\) after
conjugation, summing the bounds for products of the \(m\) kernels
bounds the squared norms before the final restriction by
\[
  \exp\!\left(B+\sum_j U_j
             +2(mx^{30}+x^{60}+mx^{50})\right).
\]
The restriction at the last step uses a branching matrix with
integer entries, which multiplies the
squared norm by at most \(\exp(x^{70})\).
Write \(u_\tau\) for the difference of the sums of the weights of
\(X_m,Y_m\) at \(\tau\). The restriction to \(Y_m\subset X_m\)
also multiplies by \(\prod_\tau\tau(\overline{\pi_N})^{-u_\tau}\).
Since \(\pi_N\overline{\pi_N}=N\), this factor has squared absolute value
\(N^{-\sum_\tau u_\tau}\) at every conjugate. The exponent therefore
decreases by at least \((|\Phi_a|a-fmr^8)\log N\).
\leansrc{New/UpperBound/Divisibility/ChainValuation}{The final restriction}.

The remainder is
\(\mathcal R=\sum_{\tau\in\Phi_a}(s_{\mathrm{fin}}(a,\tau)-s_{\mathrm{main}})\),
where \(s_{\mathrm{main}}\) is the main term for \(s_{\mathrm{fin}}(a,\tau)\).
Writing \(\overline{\log\kappa}=f^{-1}\sum_\tau\log\kappa(\tau)\) and
\(D_E=\operatorname{rd}(E)\), the condition on the level is
\[
 \log N\ge m\!\left(\overline{\log\kappa}
                 +\log \frac{M D_E}{2\pi}\right)
          -\log\frac{M D_E}{2\pi}-\tfrac{15}2\log b.
\]
The period formula
\(\log(M/\Omega)=1+c_1+\log(M D_E/(2\pi))\),
the identity \(s_{\mathrm{main}}=a(\log a-\overline{\log\kappa}+c_1)\),
and Stirling's upper bound
\(\log(a!)\le a\log a-a+1+\frac12\log a\) give
\[
 s_\infty(a)-s_{\mathrm{main}}
 \le a\!\left(\overline{\log\kappa}+\log \frac{M D_E}{2\pi}\right)
        +1+\tfrac12\log a .
\]
Consequently \(m(s_\infty(a)-s_{\mathrm{main}})-a\log N\le32a\log b\).
This uses \(m(1+\frac12\log a)\le a\log b\) and
\(\log(M D_E/(2\pi))\le23\log b\).
Substituting \eqref{eq:upper-gain-sum} gives
\[
 H_0-m\mathcal R+2\sum_p G_p\log p
 =m|\Phi_a|s_{\mathrm{main}}+mk\sum_\tau s_\infty(\kappa(\tau)).
\]
Use this identity and \eqref{eq:upper-rows} and absorb
\(32fa\log b+x^{70}+fmr^8\log N+
  2(mx^{30}+x^{60}+mx^{50})\) into \(E_3/2\).
All conjugated values of \(B_{1,\Psi}\phi\) then have squared norms at most
\begin{equation}\label{eq:upper-arch}
  \exp\!\left(B+H_0-m\mathcal R+2\sum_p G_p\log p+E_3\right).
\end{equation}
\leanref{weylNormSq_restrictedChain_le}

\paragraph{A common rational normalization}
Put \(g_p=\lfloor G_p-\varepsilon_p\rfloor\) and set
\[
 C_B=\prod_{p\in\mathcal U\cup\{N\}}p^{-g_p},\qquad
 \mathfrak p_g\phi
   = C_B|\overline G|e_0\rho(g)c_{\mathrm{ss}}^{-1}\mathcal P\phi .
\]
The factor \(C_B\) is a positive rational number.
By \eqref{eq:upper-local}, \(C_B B_{\xi,\Psi_\xi}\phi\) is integral on
\(M'_\iota\), and \eqref{eq:upper-seesaw} and \eqref{eq:upper-project}
show that \(\Theta_\xi\mathfrak p_g\phi\) is integral as well.
The same \(M'_\iota\) works for every allowed twist.

The rounding costs at most \(\sum_{p\in\mathcal U\cup\{N\}}\log p\le b\).
Using \eqref{eq:upper-error-sum} we get
\begin{equation}\label{eq:upper-cb}
 \log C_B\le-\sum_p G_p\log p+E_3+b .
\end{equation}
Multiplication of the squared norm by \(C_B^2\) combines with the
finite-place gains in the exponent as follows.
\[
 2\log C_B+2\sum_p G_p\log p
    =2\sum_{p\in\mathcal U\cup\{N\}}(G_p-g_p)\log p
    \le 2(E_3+b).
\]
Thus the principal finite-place term in \eqref{eq:upper-arch} cancels,
leaving this error and the projection loss \(x^{99}\).

\begin{proof}[Proof of Theorem~\ref{thm:upper-bound}]
The space \(S_0\) is spanned by forms integral at all finite places
whose values, after any conjugation, have squared norm at most
\(\exp(E_3)/4^f\). To obtain these forms, use that the stabilizer of
each lattice class is trivial and take forms supported on one class
at a time. For a class \(c\), choose short bases for its lattice and
for the lattices entering the determinant part of the weight.
The evaluation vectors of the Weyl module at roots of unity span the
module, and applying the chosen basis matrices preserves their span.
The bases are locally integral and the roots of unity have absolute
value \(1\), so the resulting forms are integral at all finite places.
The archimedean constant for the short bases is \(C=6(q_1d)(fk)^2\).
At all places together, columns of height less than \(k\) contain at
most \(fa\) boxes. The squared-norm bound is then \(C^{6(fk+fa)}\).
Admissibility gives \(4^f C^{6(fk+fa)}\le\exp(E_3)\).
\leanref{span_isShort_eq_top}

Let \(z_w\) be the generator of the scalar action at a specified place
\(w\) above \(N\). The image of \(\prod_w(1-z_w)\) is \(S\).
Scalar translation preserves integrality, and each difference
\(1-z_w\) increases the squared-norm bound by a factor of at most \(4\).
There are \(f\) specified places, so the images of the preceding forms
span \(S\) and satisfy the bound with \(B=E_3\).
Combine the exponent in \eqref{eq:upper-arch} with the projection
loss \(x^{99}\) and twice the bound on \(\log C_B\) in
\eqref{eq:upper-cb} to obtain
\[
 \|\Theta_1\mathfrak p_g\phi\|_\sigma^2
 \le \exp\bigl(H_0-m\mathcal R+x^{99}+4E_3+2b\bigr).
\]
The admissibility inequality \(x^{99}+4E_3+2b\le E_4\) proves (1).

For (2), first fix \(g_i=g\).
For each \(\iota\), apply the correcting matrix from
\eqref{eq:upper-local} to \((\phi_i)_i\) to obtain a family in \(M'_\iota\).
The same matrix sends \((\mathfrak p_g\phi_i)_i\) to a family integral
for every allowed twist. The valuation of its determinant is at most
\(\varepsilon_p\dim S_0\).
Imposing ramification loses at most \(\frac{f}{N-1}\dim S_0\)
dimensions, and \(\frac{f}{N-1}\le1/2\), so \(\dim S_0\le2\dim S\).
By \eqref{eq:upper-error-sum}, the logarithmic index for a fixed \(g\)
is at most \(2(\dim S)E_3\).
For a general family, group its members according to \(g_i\) and take
the block diagonal matrix with the corresponding correcting matrices
as blocks. Its determinant is the product of the determinants of the
blocks. Adding their logarithmic index bounds gives
\(2|\mathcal T|(\dim S)E_3\).
\end{proof}

\section{Spanning by the Poincar\'e map}\label{sec:spanning}

\begin{theorem}[Spanning theorem B]\label{thm:spanning}
The dimension \(h\) of the target space \(V\) is positive. There is a finite
set \(T\subset\bar G\) such that
\[
  |T|\dim S\leq 5h\log(1+2b^2),\qquad
  \dim\sum_{g\in T}\mathfrak p_g S\geq (1-4/b^2)h.
\]
\end{theorem}

To control the span of the Poincar\'e series, we estimate inner products between
translates by the finite level group. We first work with the full source space
\(S_0\) and then impose the ramification condition.
Let \(\mathcal H\) be the space of forms of principal congruence level on \(Y_0\),
with the unitary action \(\rho\colon\bar G\to U(\mathcal H)\). The orthogonal
projection \(e_0\) onto \(V=e_0\mathcal H\) commutes with this action.
Put \(D_X=\dim S_0>0\). We equip \(S_0\) with the inner product coming from
the double cosets and Bessel coordinates of the formalization. The norm
\(\|\cdot\|_{\mathrm{HS}}\) is the Hilbert--Schmidt norm.

\subsection{Unfolding and local bounds for the second moment}

Let \(\deg_B\) be the product of the formal degrees of the Bessel systems at
the infinite places, and let \(\mathrm{mass}_r\) be the mass of the genus of
the standard Hermitian lattice of rank \(r\). The mass is the sum over
isomorphism classes of lattices, weighted by the inverse orders of their
automorphism groups. For the target space we have
\begin{equation}\label{eq:span-dimension-mass}
 h\leq |\bar G|\deg_B\,\mathrm{mass}_r.
\end{equation}
Equality of multiplicities identifies \(h\) with the dimension of the
corresponding projected space on the definite space \(W_m\).
\leanref{targetDim_eq_finrank_goodForms_spaceW}
On \(W_m\) the classes at principal congruence level have trivial
automorphism groups. Their number is therefore the mass \(\mathrm{mass}_r\)
at maximal compact level times the level index \(|\bar G|\). By Weyl's dimension
formula, the product of the dimensions of the algebraic representations
attached to each class is \(\deg_B\). Hence the dimension before projection
is the right-hand side of \eqref{eq:span-dimension-mass}.
\leanref{finrank_principalForms_spaceW}

Let \(J\) be the principal congruence level and \(\nu\) the covolume of the
real unipotent lattice, which is independent of the class. In the coordinates
of the formalization, scale the Poincar\'e map by \((\operatorname{vol}(J)\nu)^{-1/2}\)
to obtain \(P\colon S_0\to\mathcal H\), and put
\[
 P_s=P^*\rho(s)P \qquad(s\in\bar G).
\]
Let \(c,c'\) be classes of the source space on \(U(X_0)\), let \(h_c,h_{c'}\)
be the images of their representatives in \(U(Y_0)\), and write \(W\) for the
Bessel function and \(\beta_c,\beta_{c'}\) for the class cutoffs. Unfolding the
inner product of the Poincar\'e series gives the \((c,c')\) block of \(P_s\)
as the entrywise complex conjugate of
\begin{equation}\label{eq:span-unfolded}
 \frac1\nu
 \sum_{\substack{\gamma\in U(Y_0)(F)\\
       h_c^{-1}\gamma h_{c'}\in K,\!
       \operatorname{red}(h_c^{-1}\gamma h_{c'})=s}}
 \int
    \beta_c(y)\beta_{c'}(\gamma^{-1}y)\,
    W(y)W(\gamma^{-1}y)^*\,dy .
\end{equation}
Here \(F\) is the maximal totally real subfield of \(E\), \(K\) is the lattice
stabilizer in the finite adeles, and the integral is over the product at the
infinite places. \leanref{hasAmplitudeBlocks_scaledPoincare}

Let \(R_0\) be the rational Bessel subgroup, and let \(\bar R_0\subset\bar G\)
be the reduction of its integral points. Restricting
\eqref{eq:span-unfolded} to \(\gamma\in R_0\) gives the unitary representation
\(r_\psi\) of \(\bar R_0\), extended by zero to \(\bar G\). Set
\[
  E_s=P_s-r_\psi(s),\qquad
  D_{\mathrm{id}}=|\bar R_0|D_X,\qquad
  D_{\mathrm{main}}=\frac{D_X^2}{\deg_B\,\mathrm{mass}_r}.
\]
Write \(\bar G=G_0\times G_1\), where \(G_i\) is the product over the chosen
places above \(p_i\), viewed also as a subgroup of \(\bar G\).
The calculation of the terms on \(R_0\) gives
\begin{equation}\label{eq:span-equivariance}
 P_t=P_1r_\psi(t),\qquad E_{ts}=r_\psi(t)E_s
 \quad(t\in\bar R_0,\ s\in\bar G).
\end{equation}
Indeed, the action of \(r_\psi(t)\) on the source reindexes the terms of
the unfolded expression. \leanref{besselTranslate_mul_left}

To bound the second moment on larger sets, group the terms in
\eqref{eq:span-unfolded} by their correspondences attached to the isotropic
part \(I\) of the Bessel frame. The correspondence of a term \(\gamma\) is
the subspace consisting of the pairs of images of \(y\) and \(\gamma^{-1}y\)
in \(I^\perp/I\), for \(y\) orthogonal to both \(I\) and \(\gamma I\).
For fixed classes, reduction \(s\in\Omega\), and correspondence, sum the
amplitudes of the truncated Bessel function. Let \(\mathcal L(\Omega)\) be
the sum of the squared norms of these sums divided by \(\nu^2\).
The terms included are those for which
\(\mathfrak d(\gamma)=\Delta\overline{\Delta}-N^2\Delta'\overline{\Delta'}\ne0\),
where \(\Delta\) is the determinant of the frame pairing matrix and
\(\Delta'\) is its preceding leading principal minor.
This is the normalized sum over links.
\leanref{linkSecondMoment}

We use the primitive \(5040\)-th root of unity in \(E\) and the Weil splitting
character chosen in the formalization. The Bessel system and level constructed
from admissible data satisfy the size and level hypotheses of the following
three estimates. \leanref{filling}

\begin{enumerate}
\item
If \(\Omega\) is disjoint from the cosets of \(\bar R_0\) obtained by
multiplying by nontrivial roots of unity, then
\begin{equation}\label{eq:span-cutoff}
 \Bigl(\sum_{s\in\Omega}\|E_s\|_{\mathrm{HS}}^2\Bigr)^{1/2}
 \leq \mathcal L(\Omega)^{1/2}+b^{-15}D_{\mathrm{id}}^{1/2}.
\end{equation}
Section \ref{sec:span-pair-count} bounds the cross terms in the sum of squares
and the error from truncation.
\leanref{sqrt_secondMoment_compl_besselS_le}\leanref{secondMoment_remainingTerms_le}
\item
Among the chosen places and their conjugates, take those at which the
span of the reduced frame has a nonzero vector in common with its image
under \(s\). Let \(\operatorname{Meet}(s)\) be the product of their norms.
If a positive integer \(z\) is at most \(\operatorname{Meet}(s)\) for every
\(s\in\Omega\), then
\begin{equation}\label{eq:span-intersect}
 \mathcal L(\Omega)\leq
       \left(\frac{b^{2f}}z\right)^k D_{\mathrm{main}}.
\end{equation}
Section \ref{sec:span-pair-count} counts pairs in the same correspondence
using the lattice in the intersection. \leanref{linkSecondMoment_le_of_le_meetNorm}
\item
Let \(A\) be a set of chosen places above \(p_0\). For each \(w\in A\),
specify a nonzero vector in the span of the frame at \(w\) or its conjugate.
Write \(\Gamma_A\) for the subgroup fixing these vectors.
Call a reduction \(s\) open if at some chosen place above \(p_0\) all
proper leading principal minors of positive order of the frame pairing
matrix are nonzero. Let \(\Omega_A\subset\Gamma_A\) be the open reductions.
Then
\begin{equation}\label{eq:span-open-input}
 \mathcal L(\Omega_A)\leq (1+b^{-4})b^{-6k|A|}D_{\mathrm{main}}.
\end{equation}
Section \ref{sec:span-open-proof} applies Siegel's estimate and Poisson
summation, deriving a factor \(b^{-6k}\) per fixed-vector condition from
the local mass. \leanref{linkSecondMoment_isOpenElement_le}
\end{enumerate}

\subsection{Counting with the common lattice}\label{sec:span-pair-count}

Put \(n=\varrho+1\), \(n_X=\#\{c\}\), \(\tau=b^{kf/4}\), and
\(n_p=|\bar R_0|n_X/N^{kf}\). Group the terms of the truncated Bessel function
into double orbits under the integral unipotent groups. A term
\(\gamma\) is near if there is a point \(y\) in the product at the infinite
places such that both \(y\) and \(\gamma^{-1}y\) lie in the region of reduced
points with bounded Bessel lengths. In each orbit containing a near term,
choose one of them as representative. A reduction \(s\) matches the left and
right unipotent characters if their values agree on integral unipotent
elements satisfying \(us=su'\) after reduction. If this condition fails,
the change of variables by such a pair multiplies the amplitude by a
nontrivial character ratio and forces it to vanish.

The terms to be counted are those with \(\mathfrak d(\gamma)\ne0\) and those
with \(\mathfrak d(\gamma)=0\), \(\gamma I\ne I\), and matching rational
unipotent characters. The latter matching condition means equality of
the character values of rational unipotent elements satisfying
\(u\gamma=\gamma u'\). This set is stable under the rational unipotent
actions on both sides.
\leansrc{New/Spanning/Unfolding/TermStructure}{The remaining terms and rational character matching}

The lattice \(L_\gamma\) of a near term \(\gamma\) consists of the vectors
in \(X_0\) in the lattice of class \(c\) that can be made orthogonal to the
frame of \(\gamma\) by adding \(M\) times an integral isotropic vector.
For two representatives \(\gamma,\gamma'\) over \((c,s,c')\), take the
intersection of their correspondences and impose the vanishing of the
last frame component in both coordinates. Project to \(X_0\) by the first
coordinate and restrict \(L_\gamma\cap L_{\gamma'}\) to this subspace.
Denote the resulting lattice by \(L_{12}\), its corank by \(l\), and its
discriminant by \(d\). The discriminant is the square root of the generalized
index of the Gram determinant of a lattice basis. Each correspondence has
dimension \(k+1\). If their intersection has codimension \(i\) in each
correspondence, the two last-coordinate conditions give \(l\le i+1\).
When \(i=0\), vanishing of one \(\mathfrak d\) implies vanishing of both, and
then \(l=0\). The bounds on lengths and lattices for near terms give
\(l\le n^3\) and \(1\le d\le b^{n^3 f}\).
\leansrc{New/Spanning/RemainingTerms/Gain}{Rank and discriminant bounds for the common lattice}

For a term of cell rank \(j\), define \(G_\gamma\) by taking the larger of
\(N^f\) times the normalized norm of the leading minor and the norm of
\(\mathfrak d(\gamma)\), multiplying by \(M^{2\min(\varrho,j)f}\), and dividing
by the product of the two indices of the isotropic row lattices.
For \(1\le z\le\operatorname{Meet}(s)\), set \(\zeta=z\) when \(i=0\) and
\(\mathfrak d(\gamma)\ne0\), and set \(\zeta=1\) otherwise. The discriminant
comparison for the two lattices gives
\begin{equation}\label{eq:span-pair-gain}
 (d\zeta)^2\le (2/M)^{2if}G_\gamma G_{\gamma'}.
\end{equation}
The factor for \(\zeta=z\) comes from a determinant. Over a residue field
where the frames meet, express a common vector in both frames and take
the largest index at which its two coefficients are not both zero.
Character matching allows the same pair of coefficients to occur at
the next index in a vector common to the two flags. Indeed, all linear
forms killing the columns before the next index in both frames take the
same value on the next columns weighted by these two coefficients.
Iteration gives a relation between the columns of both frames with at
least one nonzero last coefficient.
\leansrc{New/Spanning/Unfolding/Climbing}{Passing the coefficients of a common vector up the flags}
Pair this relation with the two integral lifts orthogonal to the earlier
part of the frames. The two last columns then have a nonzero kernel.
The resulting determinant \(\delta\) vanishes at every place counted in
\(\operatorname{Meet}(s)\). The quotient
\(\widehat{\mathcal O}_E/\delta\widehat{\mathcal O}_E\) maps onto the product
of the residue fields at these places, so the generalized index of
\(\delta\) is at least \(\operatorname{Meet}(s)\).
\leansrc{New/Spanning/RemainingTerms/Cut}{The determinant at places where the frames meet}

To obtain the comparison, split the common part of the correspondences
into the subspace cut out by the two last-coordinate conditions and its
complement. Then take the orthogonal complement of the common part in
each correspondence. Write \(d_\gamma^0\) for the discriminant of the
part of \(L_\gamma\) cut out by the last-coordinate conditions,
\(a_\gamma\) for the torsion index from the isotropic rows, and
\(\Delta_\gamma\) for the squared discriminant of the lattice outside the
common part. The determinant factor and the factorization of Gram
determinants give
\((a_\gamma d_\gamma^0\zeta)^2\Delta_\gamma\le G_\gamma^2\).
The two outer complements pair nondegenerately. Since the reductions
agree, the pairing defined by the difference of the Hermitian forms in
the two coordinates is integral after division by \(M\). Estimating its
determinant by the two Gram determinants gives
\(1\le(2/M)^{4if}\Delta_\gamma\Delta_{\gamma'}\).
Imposing the conditions of \(L_{\gamma'}\) on the common part of
\(L_\gamma\) increases its index by at most \(a_{\gamma'}\). Interchanging
the two terms gives the other comparison. Multiply
\(d\le d_\gamma^0a_{\gamma'}\) and \(d\le d_{\gamma'}^0a_\gamma\), then
eliminate the complementary discriminants to obtain
\eqref{eq:span-pair-gain}.
\leansrc{New/Spanning/RemainingTerms/Gain}{Discriminant comparison for two correspondences}

Fix the corank \(l\). A term of a given label in \(\bar R_0\) is
determined by its frame modulo the level-\(M\) unipotent moves of class
\(c\). To count frames with \(L\subset L_\gamma\), first fix the cell rank
\(j\) and the matrix \(B\) of frame products. Let \(B_j\) have the
first \(j\) rows of \(B\) and the remaining rows of the identity.
Its determinant \(\delta_j\) is the nonzero leading minor of order \(j\). Write \([a]\) for the index
\(|\widehat{\mathcal O}_E/a\widehat{\mathcal O}_E|\) of an integral finite idele \(a\).

Write a basis of \(L\) as \(Y=yA_L\), with \(k-l\) independent rational
columns in \(y\) and \(A_L\) invertible over the finite adeles.
The hypothesis \(L\subset L_\gamma\) gives a lift \(w\) satisfying
\[
 X_\gamma^*y+M B_{\mathrm{iso}}^*w=0,\qquad
 wA_L\text{ is integral},
\]
where \(X_\gamma\) is the matrix of \(X_0\)-components of the frame and
\(B_{\mathrm{iso}}\) consists of the isotropic rows of \(B\).
Here \({}^*\) denotes conjugate transpose. Reduce each row of \(wA_L\)
modulo the row lattice of \(Y^*Y\). That lattice has index
\(\operatorname{disc}(L)^2\), giving at most
\(\operatorname{disc}(L)^{2\varrho}\) classes for all \(\varrho\) rows.
Within one class a unipotent move makes the right-hand side of the
orthogonality condition for \(X_\gamma^*y\) the same for all frames.

The \(n\) columns of a difference of the matrices \(X_\gamma\) now lie
in the rank-\(l\) lattice orthogonal to \(L\). Reduction modulo
\(M\delta_j\) gives \([M\delta_j]^{ln}\) classes. Within one such class
unipotent moves make the \(X_0\)-components coincide. For the remaining
isotropic components, first fix the integral matrix of adjugate
entries recording the common vectors from stage \(j\) onwards. The
near condition bounds each entry by \(b^{182n^2}\) at infinity, giving
at most \(b^{366n^4 f}\) such matrices. Finally reduce the \(\varrho n\)
entries of the isotropic block modulo \(M\det(B_j B_j^*)\).
Frames in the same class coincide after a unipotent move.
\leansrc{New/Spanning/RemainingTerms/FrameParts}{The lift equation and unipotent moves}
\leansrc{New/Spanning/RemainingTerms/FramesOverLattice}{Congruence classes of frames over a fixed lattice}

For fixed \(j,B\), the number of these classes is at most
\[
 \operatorname{disc}(L)^{2\varrho}\,[M\delta_j]^{ln}
       b^{366n^4 f}[M\det(B_j B_j^*)]^{n\varrho}.
\]
Both modulus indices and \(\operatorname{disc}(L)^2\) are at most
\(b^{2n^3f}\). There are at most \(b^{2f}\) choices of \(j\) and
\(b^{270n^3f}\) choices of \(B\). Thus there are at most
\(b^{2000n^6(l+1)f}\) near representatives whose lattices contain \(L\).
A label can be chosen for each reduction with matching characters,
and the same label applies to both terms. There are consequently at
most \(b^{4000n^6(l+1)f}\le\tau\) ordered pairs.

To sum the weights \(d^{-k}\) for lattices of rank \(k-l\), average over
the cosets at maximal level, of which there are at least \((N^k/b)^f\).
In one unimodular lattice, extend a basis of the saturation of \(L\)
to a basis of the whole lattice. Write the basis of \(L\) as the
basis of its saturation multiplied by an integral matrix \(A_{\mathrm{sat}}\).
For the specified row to vanish on \(L\) modulo \(\pi_N\), its first
\(k-l\) values \(a\) must satisfy \(aA_{\mathrm{sat}}=\pi_N w\).
The class of \(w\) modulo the row lattice of \(A_{\mathrm{sat}}\) determines
\(a\) modulo \(\pi_N\), giving at most \([\det A_{\mathrm{sat}}]\) choices.
The other \(l\) values contribute
\(\bigl|\widehat{\mathcal O}_E/\pi_N\widehat{\mathcal O}_E\bigr|^l=N^{lf}\).
Distinct cosets have distinct rows modulo \(\pi_N\), so at most
\([\det A_{\mathrm{sat}}]N^{lf}\) cosets contribute for this lattice. Summed over the cosets, the weighted
lattice count with \(d\le D_\ast=b^{n^3f}\) is at most
\[
  N^{lf}D_\ast^{2\varrho+6}
       (200fk\,\operatorname{rd}_E)^{f(k+1)l}2^k
  \le b^{-f}\tau^{l+1},
\]
where \(\operatorname{rd}_E\) is the root discriminant of \(E\). Dividing by
the number of cosets bounds the average weight for one class by
\(N^{-kf}\tau^{l+1}\). Including all \(n_X\) classes, the \(|\bar R_0|\)
labels, and the multiplicity of pairs for each lattice gives
\begin{equation}\label{eq:span-pair-count}
 \sum_{\substack{\text{pairs of near representatives}\\
                            \operatorname{corank} L_{12}=l}}d^{-k}
 \le \tau^{l+2}n_p .
\end{equation}
\leansrc{New/Spanning/RemainingTerms/LatticesOfCorank}{Averaging lattice counts over cosets}
\leansrc{New/Spanning/RemainingTerms/Count}{Counting pairs weighted by their common lattice}

Schur's estimate for the orbit sum gives the bound
\[
 \sigma_\gamma=\tau^2\left(\frac{D_j}{n_pG_\gamma^k}\right)^{1/2},
 \qquad
 D_j=\begin{cases}
       D_{\mathrm{main}} & j=\varrho+1,\\
       \tau^{-16}D_{\mathrm{id}} & j\le\varrho
     \end{cases}
\]
for each near representative. More precisely, for a subset of the
remaining terms stable under both rational unipotent actions, the norm
of its truncated
amplitude divided by \(\nu\) is at most \(\sum_\gamma\sigma_\gamma\) over
the chosen near representatives. The case \(j=\varrho+1\) is precisely
\(\mathfrak d(\gamma)\ne0\).
\leansrc{New/Spanning/RemainingTerms/Size}{Amplitude bound for the double-orbit sum}
If \((d\eta)^2\le G_\gamma G_{\gamma'}\) and \(D\ge D_j,D_{j'}\), then
\[
 \sigma_\gamma\sigma_{\gamma'}
 \le \frac{\tau^4D}{n_p\eta^k}\,d^{-k}.
\]
Thus if \(\tau^{l+6}D/\eta^k\le C\) for each corank \(l\),
\eqref{eq:span-pair-count} bounds the sum of these weights at each corank
by \(C\). There are at most \(n^3+1\le\tau\) coranks, giving the bound \(\tau C\).

For equal correspondences with nonzero discriminants, \(l\le1\), and
\eqref{eq:span-pair-gain} allows \(\eta=z\) and \(D=D_{\mathrm{main}}\).
Taking \(C=\tau^7D_{\mathrm{main}}/z^k\) gives
\(\tau^8D_{\mathrm{main}}/z^k=(b^{2f}/z)^kD_{\mathrm{main}}\).
Each pair determines its reduction and its correspondence, so this
proves \eqref{eq:span-intersect}.

For equal correspondences with vanishing discriminants, take
\(l=0\), \(\eta=1\), and \(D=\tau^{-16}D_{\mathrm{id}}\).
Here \(C=\tau^{-10}D_{\mathrm{id}}\) suffices. For distinct correspondences,
\(i\ge1\) and \(l\le i+1\). Take \(\eta=(M/2)^{if}\) and
\(D=D_{\mathrm{main}}\). Since \(\rho_\ast=\tau(2/M)^{kf}\le1\) and
\(D_{\mathrm{main}}\le(M/b^6)^{kf}D_{\mathrm{id}}\), we have
\[
 \frac{\tau^{l+6}D_{\mathrm{main}}}{\eta^k}
 \le \tau^7\rho_\ast^iD_{\mathrm{main}}
 \le \tau^8(2/b^6)^{kf}D_{\mathrm{id}}
 \le \tau^{-10}D_{\mathrm{id}} .
\]
The last inequality uses \(2\le b\) and \(\tau=b^{kf/4}\).
Both cases therefore fit the same bound \(C=\tau^{-10}D_{\mathrm{id}}\),
and their total is at most \(\tau^{-9}D_{\mathrm{id}}\le b^{-kf}D_{\mathrm{id}}\).
\leansrc{New/Spanning/RemainingTerms/Sum}{Squares for one correspondence and the cross terms}

To apply the count, split each block into one piece per correspondence
with nonzero discriminant and a single piece of all terms with vanishing
discriminant. Square the sum of the norms of the pieces. Summing the
diagonal terms for the pieces of nonzero discriminant gives
\(\mathcal L(\Omega)\). Use the products
\(\sigma_\gamma\sigma_{\gamma'}\) for the other terms. A pair belongs to
only one reduction, so the same bound applies after summing over the blocks.

For the truncation error, the Bessel tail bound gives
\(q=\nu/(2b^{k^3 f})\) as an upper bound on the sum of the norms of the
amplitude differences in each row \((c,s)\). The sum of squares in the
row is at most \(q^2\). The number of rows and the inequalities
\(|\bar G|b^{40}\le(2b^{k^3 f})^2\) and \(n_X\le D_{\mathrm{id}}\) show that
the sum over all rows divided by \(\nu^2\) is at most
\(b^{-40}D_{\mathrm{id}}\). This estimate compares the full and truncated
functions on the same set of terms.
\leansrc{New/Spanning/RemainingTerms/Size}{Row estimate for the truncation error}
On the allowed set \(\Omega\), the full amplitude for \(E_s\) equals the
sum over the remaining terms. Outside the Bessel subgroup and the
remaining terms, those failing the rational character matching condition
give a vanishing sum over a stable set. Those with matching rational
characters have a norm-one coefficient pair from the first stage, and
their reductions lie in the excluded scalar cosets.
\leansrc{New/Spanning/Unfolding/SquareExpansion}{The amplitude outside the Bessel subgroup}
The triangle inequality for the square root of the sum of squares and
\(kf\ge40\) give
\[
 \Bigl(\sum_{s\in\Omega}\|E_s\|_{\mathrm{HS}}^2\Bigr)^{1/2}
 \le \mathcal L(\Omega)^{1/2}
       +(b^{-kf/2}+b^{-20})D_{\mathrm{id}}^{1/2}
 \le \mathcal L(\Omega)^{1/2}+b^{-15}D_{\mathrm{id}}^{1/2}.
\]
This proves \eqref{eq:span-cutoff}.

\subsection{The second moment on open reductions}\label{sec:span-open-proof}

\paragraph{Local mass under fixed-vector conditions}
Choose a conjugation-stable finite set of finite places \(\mathcal T\)
containing the support of the level \(M\). Write \(K_{\mathcal T}\) for the
integral unitary group over the product of completions and
\(K_{\mathcal T}(g)\) for its principal congruence subgroup of local level
\(g\). For \(K_{\mathcal T}(g)\subset Z\subset K_{\mathcal T}\) and Haar
measure \(\mu\), put
\[
 \mathfrak M(Z)=\int
    \left|\int_Z \chi_g(yxy^{-1})\,d\mu(y)\right|^2d\mu(x),
 \qquad
 m_Z=\frac{\mu(Z)\mu(K_{\mathcal T}(g))}{u_g^2}.
\]
Here \(\chi_g\) is the unipotent character extended to the product of its
subgroup with \(K_{\mathcal T}(g)\) and by zero elsewhere, and \(u_g\) is
the associated unipotent measure constant. For \(x\in K_{\mathcal T}(g)\),
\(\chi_g(yxy^{-1})=1\), so the contribution from this subgroup is exactly
\(\mu(Z)^2\mu(K_{\mathcal T}(g))\).

Suppose the local lattice data have the following property for \(Q\ge2\).
An invertible integral element whose inverse is nonintegral has modulus
at most \(Q^{-1}\). The estimate outside the congruence subgroup divides
the integral into ranges of the index of the denominator ideal. For
\(a+10^8n^4\le k\), the integral of the zeroth-stage majorant there is at most
\(m_Z2^{10^9n^4}\min(1,Q^{-a})\). If no coordinate at a position \(x\) outside the congruence subgroup
had depth at least \(Q\), one could choose integral units among the
stable frame products. This would make \((x-1)/g\) integral, a contradiction.
Thus the bound for a deep vector applies. The factor in the integral
estimate for the number of coordinates is at most \(2^{99\cdot13n}\).
The sum over dyadic ranges of the denominator index incurs at most
\(2^{4\cdot10^8n^4}\). Their product is at most \(2^{10^9n^4}\).
Taking \(a=t+10^9n^4\) absorbs the extra factor by
\(2/Q\le1\). Multiplying the majorant by \(u_g^2\) and adding the principal
congruence contribution gives, for \(t+2\cdot10^9n^4\le k\),
\begin{equation}\label{eq:span-local-mass}
 \mathfrak M(Z)
 \le\mu(Z)\mu(K_{\mathcal T}(g))\,\bigl(\mu(Z)+Q^{-t}\bigr).
\end{equation}
\leansrc{New/Spanning/Frames/FiniteShells}{The local integral estimated by denominator ideals}

On the factor for a single \(w\in A\) and its conjugate, denote the
integral group by \(K_w\) and the inverse image of the fixed-vector
condition by \(Z_w\). Reduction onto \(\mathrm{GL}_r(\mathbf F_{p_0})\) is
surjective, and the conjugate factor acts by the inverse adjoint.
Transitivity on nonzero vectors in either action gives
\([K_w:Z_w]=p_0^r-1\). Put \(B=b^9\), so \(2\le B\le p_0\).
The additive volume of nonunits in the product over the selected places
is at most \(B^{-1}\). Each residue field has cardinality at least
\(b^{10}\) and contributes at most its reciprocal, with at most \(b\) places.
The same bound applies to the modulus of an element whose inverse is nonintegral.
The admissibility conditions give \(t_0=k-10^{10}n^4\) with
\[
 t_0+1+n+2\cdot10^9n^4\le k,\qquad
 6k\le9t_0,\qquad t_0+1<r.
\]
The index formula gives \(\mu(Z_w)\le B^{-t_0-1}\mu(K_w)\).
In the normalization of the local integral,
\(\mu(K_w)\ge2^{-n}\ge B^{-n}\). Apply \eqref{eq:span-local-mass} with
\(Q=B\) and \(t=t_0+1+n\) to obtain
\[
 \mathfrak M(Z_w)\le
 \mu(Z_w)\mu(K_w(g))
       \bigl(\mu(Z_w)+B^{-t_0-1}\mu(K_w)\bigr).
\]
Multiply by \(\mu(K_w)b^{6k}\) and use \(b^{6k}\le B^{t_0}\) and \(2/B\le1\).
The contribution to \(\mathfrak M(K_w)\) from \(K_w(g)\) alone is
\(\mu(K_w)^2\mu(K_w(g))\), so
\[
 \mathfrak M(Z_w)\mu(K_w)b^{6k}
 \le\mathfrak M(K_w)\mu(Z_w).
\]
The groups and characters factor over distinct conjugate pairs. Multiply
the mass integrals and volumes, taking the full integral group on the
remaining factors. This gives the same inequality with
\(K_w,Z_w,b^{6k}\) replaced by \(K_{\mathcal T},Z,b^{6k|A|}\).
The ratio comparison is invariant under a common rescaling of the
product Haar measure.
\leansrc{New/Spanning/Frames/FixedLines}{Mass ratio for each fixed-vector condition}

For \(Z=K_{\mathcal T}\), apply \eqref{eq:span-local-mass} with \(Q=2\)
and \(t=n\). The inequalities \(\mu(K_{\mathcal T})\le1+2^{-n}\) and
\(2^{-\varrho}\le b^{-6}\) give
\(\mathfrak M(K_{\mathcal T})\le(1+b^{-6})\mu(K_{\mathcal T})\mu(K_{\mathcal T}(g))\).
The zeroth-stage orbital form is a squared correlation averaged against
the character, and its unnormalized absolute value times \(u_g^2\) is
at most \(\mathfrak M(Z)\). Dividing by \(m_Z\) bounds the normalized
zeroth-stage finite-place kernel by
\begin{equation}\label{eq:span-local-zero}
 m_A=(1+b^{-6})b^{-6k|A|}.
\end{equation}
\leansrc{New/Spanning/MainTerm/FiniteSieveSystem}{Zeroth-stage finite-place density}

\paragraph{Siegel's estimate at a fixed last entry}
Take finitely many nonzero amplitude blocks \((c,s,c',\Lambda)\) over
\(\Omega_A\). All proper leading minors are nonzero at the rational
point, so the big-cell decomposition names each block by its last
entry \(t\in E^\times\), a representative \(\sigma\) fixing the preceding
hyperbolic basis vectors, and an integral lift \(n_s\).
They satisfy \(|\Phi_v(t)|^2>N^2\) at every infinite place \(v\), and
\((c,c',\operatorname{red}(n_s),\sigma)\) determines the block.
With \(\sigma\) fixed, the remaining leading entries are in a finite set
\(T_0\subset E^\times\). The block amplitude can be written
\(\mathcal A=\sum_{\mathbf p\in T_0^\varrho} a(\mathbf p)\), where
\(a(\mathbf p)\) is the truncated unipotent double-coset amplitude of the
Weyl representative \(w(\mathbf p,t)\sigma\). Enlarge \(\mathcal T\) so
that the elements of \(T_0\) and their inverses, \(M^{-1}\), \(N^{-1}\),
the lattice basis and its inverse, and the class translates of each
\(\sigma\) are integral outside \(\mathcal T\).
\leansrc{New/Spanning/MainTerm/EnergyEstimate}{Big-cell amplitude sums and reduction to finite sets}

At an infinite place \(v\), let
\(\lambda_{v,j}=\kappa_v^{2(\varrho-j)}/N\) be the scale of entry \(j\), where
\(\kappa_v=M\Phi_v(\varpi)/(2\pi\mathrm i)\) for the generator \(\varpi\) of the
different. Let \(R_t(\mathbf p,\mathbf q)\) be the product of the top
kernels at the infinite places at last entry \(t\). These are Gram kernels
of orbital integrals with weights given by the coordinate scales and the
densities of the preceding entries. The last-entry finite-place density
\(F_t(\mathbf p,\mathbf q)\) is an integral over a quadric of the Gram
pairings on fibers with a common last vector.

For a fixed \(t\), the \(X_0\)-component \(y\) of the image of the last
frame vector under \(\sigma\) satisfies \(N\langle y,y\rangle=t\bar t-N^2\).
Set
\[
 G_t=\left(\prod_v(|\Phi_v(t)|^2-N^2)\right)^{k-1}>0 .
\]
Write \(d_{B,v}\) for the number of Bessel coordinates at \(v\) and \(d_B\)
for their product. With the big-cell volume constant
\(c_G=\operatorname{vol}(U(k))^2/
(\operatorname{vol}(U(r))\operatorname{vol}(U(k-1)))\), the weight of the
archimedean orbital integral is
\[
 \rho_t(\mathbf p)=
 \prod_v \frac{\sqrt{c_G}}{d_{B,v}}
 \left(\frac{|\Phi_v(t)|^2}{N^2}-1\right)^{(k-1)/2}
 \prod_{j<\varrho}|\lambda_{v,j}\Phi_v(p_j)|^{k-1+2(\varrho-j)} .
\]
Let \(h_t(\mathbf p)\) be the corresponding finite-place density of
the head. Since \(p_j\) is a unit outside \(\mathcal T\), the product
formula identifies \(\prod_v|\Phi_v(p_j)|\) with the inverse square root
of the modulus of \(p_j\) at \(\mathcal T\). This gives
\(\rho_t(\mathbf p)h_t(\mathbf p)=r_0\sqrt{G_t}\) with
\[
 r_0=\frac{1}{d_B}\prod_v
       \left(\sqrt{c_G}\prod_{j<n}\lambda_{v,j}^{k-1+2(\varrho-j)}\right).
\]
Put \(D=|\bar R_0|M^{f\varrho(2k+2\varrho+1)}\) and \(C_S=D\,n_Xc_S\),
where \(c_S\) is the constant in the Siegel sum on the sphere.
We first sum the blocks with a common point \(x=(c,y)\).
By \(\Psi_t(x)\) we mean evaluation at the last local vector obtained
from \(y\) after transporting it by the inverse of the class representative.
Put \(\mathcal K=K_{\mathcal T}(g)\) and \(\mathcal C=Z/\mathcal K\).
The scaled orbital integral for a head \(\mathbf p\) has coordinates
in the finite set \(\mathcal C^2\). For a position \(\ell\) and a
coordinate \((a,a')\in\mathcal C^2\), denote the vector of these
coordinates, indexed by the heads, by \(v_{\ell,a,a'}\).
The sum of the finite-place pairings weighted by \(R_t\) is
\[
 Q_t(\ell)=\sum_{a,a'\in\mathcal C} v_{\ell,a,a'}^*R_t v_{\ell,a,a'}\ge0 .
\]
For fixed \(c,c',\sigma\), the block with reduction \(s\) occurs in
the coordinates where \(a'\) is the coset of \(n_s^{-1}a\).
Distinct reductions occupy disjoint sets of coordinate pairs.
Unfolding the orbital integral identifies each such summand with
\(\mu(\mathcal K)^2G_t r_0^2\|\mathcal A\|^2\). Summing over \(a\) gives
a factor \(|\mathcal C|\), and the sum over all the blocks is bounded by
\(Q_t(\ell)\).

Write \(H_{\mathrm{fix}}\) for the local stabilizer of the last frame
vector in the subgroup fixing the hyperbolic pairs. The positions
\(\ell\) for \((c',\sigma)\) with the same \(x\) lie in a single orbit
under left multiplication by \(H_{\mathrm{fix}}\). They lie in distinct
congruence cosets. Indeed, congruence at \(\mathcal T\) together with
integrality away from \(\mathcal T\) identifies both the source class
and the rational representative. The integrand \(Q_t\) is invariant
under left multiplication by \(H_{\mathrm{fix}}\cap\mathcal K\).
Multiplying the sum by the measure of a congruence coset thus bounds it
by the integral on the fiber. For the Haar measure \(\mu_H\) on
\(H_{\mathrm{fix}}\), the mass comparison gives
\(D\mu_H(H_{\mathrm{fix}}\cap\mathcal K)\mu(\mathcal K)^2|\mathcal C|\ge m_Z\).
The step function \(\Psi_t\) extending the finite-place fiber pairing
therefore satisfies
\[
 G_t(m_Z/D)\sum_{(c,y)=x} r_0^2\|\mathcal A\|^2\le\Psi_t(x).
\]
\leansrc{New/Spanning/Frames/TopFloor}{Finite quotient coordinates and the fiber pairing}

Put \(\beta_t=(t\bar t-N^2)/N\). Apply Siegel's formula for one vector in
the definite space of rank \(k\ge3\). The test function imposes
integrality away from \(\mathcal T\) and pulls back \(\Psi_t\) along the
last-vector map at \(\mathcal T\). It is compactly supported, locally
constant, and invariant under the source level. The norm \(\beta_t\)
is totally positive. The class sum is the finite adelic norm density
times
\[
 n_X|\operatorname{disc}(E)|^{1/4}
       \prod_v\frac{\pi^k(\Phi_v(\beta_t))^{k-1}}{(k-1)!}.
\]
\leansrc{Lib/SiegelWeil/OneVector/SiegelClassSum}{Siegel's formula at one norm}

The norm density integrates the test function times
\(\psi_f(u(w^*w-\beta_t)/2)\) over the finite vectors \(w\) and
self-adjoint finite adelic frequencies \(u\). Splitting the vectors
and frequencies by their components at \(\mathcal T\) factors the
density into the sphere density there and an integral of Gauss sums
from the integrality condition outside \(\mathcal T\). Write
\[
 F_{u,\Gamma}(w)=\psi_f\!\left(u(w^*w-\Gamma)/2\right),
 \qquad w\in\widehat{\mathcal O}_E^{\,k}.
\]
Here \(\Gamma\) is the shift from the prescribed norm that remains in
the Gauss sum, and the bound is uniform in \(\Gamma\).
The square of the absolute value of \(\int F_{u,\Gamma}(w)\,dw\)
is the integral of \(\overline{F_{u,\Gamma}(w)}F_{u,\Gamma}(w+h)\)
over \(w,h\). Translating coordinate \(j\) of \(w\) by \(\bar a\)
multiplies the correlation by \(\psi_f(u h_j a)\).
Consequently the integral over \(w\) vanishes unless every \(h_j\)
lies in the conductor ideal
\[
 I_u=\{h\in\widehat{\mathcal O}_E:
            \psi_f(uha)=1\text{ for all }a\in\widehat{\mathcal O}_E\}.
\]
The integrand has modulus \(1\), so the remaining correlation is
bounded by the product of the volumes. With
\(q_u=|\widehat{\mathcal O}_E/I_u|\), it follows that
\[
 \left|\int F_{u,\Gamma}(w)\,dw\right|
 \le \operatorname{vol}(\widehat{\mathcal O}_E)^k q_u^{-k/2}.
\]
For a fixed conductor \(I\), the frequency classes modulo the
self-adjoint codifferent are distinguished by the characters of
the finite additive group \(\widehat{\mathcal O}_E/I\). Each class has
measure \(1\), and there are at most \([\widehat{\mathcal O}_E:I]\) of them.
With \(\mu_0=3\varrho+4\) and \(k\ge2(\mu_0+1)\), integration over the
frequencies is therefore bounded by
\(\operatorname{vol}(\widehat{\mathcal O}_E)^k
\sum_I[\widehat{\mathcal O}_E:I]^{-\mu_0}\).
The sum over finite-index ideals is at most
\(1+2^{-\varrho}\le1+b^{-6}\), by comparison with the Dedekind zeta
Euler product using \([E:\mathbf Q]\le2^{2\varrho}\).
\leansrc{Lib/SiegelWeil/Mass}{Character orthogonality and the Gauss-sum bound by conductor}

Finite-place Fourier inversion bounds the density at \(\mathcal T\)
by the integral of the test function on the sphere
\(\langle y,y\rangle=\beta_t\). The density \(F_t\) divides this integral by
the modulus \(N^{-f}\) of \(\pi_N\). Expressing the integral in terms of
\(F_t\) multiplies the archimedean norm factor \(N^{-f(k-1)}\) by \(N^{-f}\),
giving
\[
 c_S=\left(\frac{\pi^k}{(k-1)!N^k}\right)^f
       |\operatorname{disc}(E)|^{1/4}
       \operatorname{vol}(\widehat{\mathcal O}_E)^k.
\]
Summing over the distinct points \(x\) gives
\[
 \sum_x\Psi_t(x)\le(1+b^{-6})n_X c_S G_t\,
             \operatorname{Re}\sum_{\mathbf p,\mathbf q}R_t(\mathbf p,\mathbf q)
                                                       F_t(\mathbf p,\mathbf q).
\]
Cancel \(G_t\), divide by \(m_Z\), and sum over \(t\) to obtain
\begin{equation}\label{eq:span-siegel-sum}
 r_0^2\sum_{\mathrm{blocks}}\|\mathcal A\|^2
 \le (1+b^{-6})C_S\sum_t\operatorname{Re}\sum_{\mathbf p,\mathbf q}
                         R_t(\mathbf p,\mathbf q)\frac{F_t(\mathbf p,\mathbf q)}{m_Z}.
\end{equation}
\leansrc{New/Spanning/MainTerm/SiegelTopBound}{Siegel's estimate for the class sum and its local densities}

\paragraph{Removing the entries by Poisson summation}
Let \(J_i\) be the product of the archimedean kernels and the normalized
finite-place kernel at stage \(i\). Let \(\mathcal S_{i+1}\) be the real part of
the sum of \(J_{i+1}\) over two heads \(\mathbf p,\mathbf q\in T_0^i\)
followed by a common entry \(t\in T_0\). Entries after the heads are filled
with \(1\), and each entry at infinity is multiplied by its scale
\(\lambda_{v,j}\). At stage zero \(\mathcal S_0=\operatorname{Re}J_0\) on the empty
heads. The sum for each common entry is nonnegative by the Gram property,
so the kernel sum on the right of \eqref{eq:span-siegel-sum} is at most \(\mathcal S_n\).

Apply the one-sided Poisson summation bound to entry \(i\).
The sum of the kernels with that entry held in common is a nonnegative
lower semicontinuous function. Multiply it by the indicator of the
integral adeles outside \(\mathcal T\). By the integral identity for the
orbital forms, the zero frequency gives the sum of \(J_i\) over all
pairs of heads, with multiplier
\[
  q_i=\frac{\operatorname{vol}(\widehat{\mathcal O}_E)}
                 {\prod_v\lambda_{v,i}^2}.
\]
The other frequencies are rational points integral outside \(\mathcal T\)
and scaled by \(\varpi^{-1}\). Let \(\epsilon\) bound the sum of the
products of the local Fourier decay factors at the nonzero frequencies,
and let \(W_i\) be the sum of the nonnegative weights of heads of length \(i\).
The contribution from those frequencies is at most \(q_i\epsilon W_i^2\).
In the sum of \(J_i\) produced by the zero frequency, the pairs with
equal last entries give \(\mathcal S_i\). For unequal last entries \(u,v\), the
finite-place frequency \(M^{-1}(u^{-1}-v^{-1})\) is nonzero and integral
outside \(\mathcal T\). The same decay estimate bounds these pairs by
\(\epsilon W_i^2\), giving
\[
 \mathcal S_{i+1}\le q_i(\mathcal S_i+2\epsilon W_i^2).
\]

For the frequency bound put \(K_0=k-n^7\).
At infinity the decay function is \(d_{K_0}(a)=a^{K_0}\min(1,a^6)\).
The factor at \(v\) and its volume ratio are
\[
 c_{\mathrm{vol}}b^{2000n^9}d_{K_0}(|\xi_v|^{-1}),
 \qquad
 c_{\mathrm{vol}}=\pi^{(2\varrho+1)k}
            \frac{\operatorname{vol}(U(k))}{\operatorname{vol}(U(r))}.
\]
At the finite places the factor is the generalized index of the frequency
to power \(K_0/2\), divided by its denominator index cubed, times \(b^{n^9 f}\).
The frequencies associated to a nonzero \(\mathcal T\)-integer \(y\)
have the form \(\xi_v=(c_v/\lambda_{v,j})\Phi_v(y)\) at infinity.
In both error terms the scales satisfy \(|c_v|^{-1}\le|\kappa_v|/M\), so
\[
 |c_v/\lambda_{v,j}|^{-1}\le
 \theta_v:=|\kappa_v|^{2\varrho+1}/(MN)
 \le B_F:=b^{23(2\varrho+1)}.
\]
Write \(H(y)=\operatorname{Norm}(\mathfrak a_y)\prod_v\max(1,|\Phi_v(y)|)^2\)
for the height of \(y\), where \(\mathfrak a_y\) is its denominator ideal.
Since \(y\) is integral outside \(\mathcal T\), the product formula bounds
the square root of the finite-place generalized index by
\(\prod_v|\Phi_v(y)|\). This upper bound cancels the factors
\(|\Phi_v(y)|^{-K_0}\) at infinity.
Since \(\min(1,|\xi_v|^{-1})\le B_F/\max(1,|\Phi_v(y)|)\) and the
finite-place denominator index is at least \(\operatorname{Norm}(\mathfrak a_y)\),
the contribution of one frequency, with the powers of \(b\) and volume
ratios factored out, is at most \(H(y)^{-3}\prod_v(\theta_v^{K_0}B_F^6)\).
The inverse cubes of the heights sum to at most
\((64\cdot38\exp(4))^f\le b^{3f}\).
Together with \((bB_F^2)^3\le b^{n^9}\), this gives
\[
 \epsilon\le b^{2002n^9f}\mathcal D_\ast,\qquad
 \mathcal D_\ast=\prod_v c_{\mathrm{vol}}\theta_v^{K_0}.
\]
\leansrc{New/Spanning/MainTerm/GlobalSieve}{Summing frequencies using the product formula and height}

For the head weight, each prefix product contributes its finite-place
generalized index to power \(n^6/2\). At each infinite place, it
contributes the larger of \(1\) and the inverse modulus of the scaled
prefix product, raised to \(n^6\). The weight is zero if a prefix product
is nonintegral or if its scaled modulus exceeds \(b^{60n^2}\).
A head of nonzero weight is determined by its sequence of prefix
products, which are algebraic integers of modulus at most \(b^{120n^2}\)
in every embedding. The lattice count bounds the number of such heads
by \((1+2b^{120n^2})^{2fi}\). By the product formula each weight is at
most \(b^{120n^8fi}\), giving \(W_i\le b^{130n^9f}\) for \(i\le\varrho\).
Also \(\lambda_{v,j}\le b^{46\varrho}\) and
\(\operatorname{vol}(\widehat{\mathcal O}_E)=(2/\operatorname{rd}_E)^f\)
give \(\prod_{j<i}q_j\ge b^{-100n^2f}\). Hence
\[
 W_i^2\le b^{260n^9f+100n^2f}\prod_{j<i}q_j
       \le b^{400n^9f}\prod_{j<i}q_j .
\]
\leansrc{New/Spanning/MainTerm/GlobalSieve}{Counting prefix products and bounding head weights}
Write \(\deg_v\) for the formal degree of the Bessel system at \(v\).
Using \(\mathcal S_0\le\prod_v((1+b^{-10})/\deg_v)\,m_A\) and
\(2n\le b^{n^9f}\), iterate the recurrence \(n\) times to obtain
\begin{equation}\label{eq:span-sieve-sum}
 \mathcal S_n\le \left(\prod_{j<n}q_j\right)
  \left[
    \left(\prod_v\frac{1+b^{-10}}{\deg_v}\right)m_A
    + b^{2403n^9f}\mathcal D_\ast
  \right].
\end{equation}

\paragraph{Comparison of constants}
Let \(\Lambda_\ast\) be the product of the entry scales to powers
\(k+2(\varrho-j)\). The formulas for \(r_0\) and the Poisson multipliers
\(q_j\) give
\[
 V_\ast=\frac{\prod_{j<n}q_j}{r_0^2}
 =\frac{\operatorname{vol}(\widehat{\mathcal O}_E)^n}{c_G^f}
    \left(\frac{d_B}{\Lambda_\ast}\right)^2 .
\]
Substituting the covolume formula for \(\nu\) into the comparison of
Siegel constants gives
\[
 \begin{split}
 C_S V_\ast
  &\le(1+b^{-9})D_{\mathrm{main}}\deg_B\nu^2,\\
 C_S V_\ast\mathcal D_\ast
  &\le b^{n^9f}d_BD_{\mathrm{id}}\nu^2 .
 \end{split}
\]
The first inequality substitutes the class count and volumes into the
main term. For the second, put \(s_d=\sqrt{\operatorname{rd}_E}\),
\(Q_d=Ms_d^2/(2\pi)\), and \(x_d=Q_d^{2\varrho+1}/(MN)\).
Use \(\prod_v|\kappa_v|=Q_d^f\) and cancel the \(k\)-th power in
\(x_d^{k-n^7}\) against the powers in the Siegel constant. The ratio
\(C_SV_\ast\mathcal D_\ast/(d_BD_{\mathrm{id}}\nu^2)\) is the \(f\)-th
power of
\[
 \frac{(2/s_d^2)^{n+2\varrho^2}(s_d/s_d^{2\varrho})}
      {2M^{\varrho(2\varrho+1)}}\,x_d^{-n^7}.
\]
Since \(2/s_d^2\le b\), \(s_d/s_d^{2\varrho}\le1\), and
\(x_d^{-1}\le MN\le b^{23+60n}\), this factor is at most \(b^{n^9}\).
\leansrc{New/Spanning/MainTerm/EnergyEstimate}{Comparing the covolume and decay constants}

Put \(Y=b^{-6k|A|}D_{\mathrm{main}}\nu^2\). The admissibility inequalities
\(d_B\le b^{2kf}\), \(b^{11kf}D_{\mathrm{id}}\le D_{\mathrm{main}}\),
\(|A|\le f\), and \(512000n^9\le k\) give
\(2404n^9f+2kf+10+6k|A|\le11kf\). Thus the error in
\eqref{eq:span-sieve-sum}, after accounting for its normalization, is
\(C_SV_\ast b^{2403n^9f}\mathcal D_\ast\le b^{-10}Y\).
For the main term use \((1+b^{-10})^f\le1+b^{-9}\) and
\(\prod_v\deg_v=\deg_B\). The former follows from \(f\le b/1000\), which
gives \((1+b^{-10})^f\le\exp(b^{-9}/1000)\le1+b^{-9}\). Combining \eqref{eq:span-siegel-sum},
\eqref{eq:span-local-zero}, and \eqref{eq:span-sieve-sum} yields
\[
 \sum_{\mathrm{blocks}}\|\mathcal A\|^2
 \le (1+b^{-6})\left[(1+b^{-9})^2(1+b^{-6})+b^{-10}\right]Y
 \le (1+b^{-4})Y .
\]
For the last inequality, \(b\ge1000\) implies
\(b^{-6},b^{-9},b^{-10}\le s_\ast=b^{-4}/100\le1\), so the coefficient
is at most \((1+s_\ast)^4+s_\ast(1+s_\ast)\le1+17s_\ast\).
Taking the supremum over the finite sums and dividing by \(\nu^2\)
proves \eqref{eq:span-open-input}.

An element of \(\bar R_0G_i\) agrees with an element of \(\bar R_0\) above
the other prime. The frames therefore meet at all \(f\) chosen places
above that prime and at their conjugates, so \(z\) may be taken to be
that prime to the power \(2f\). The set also avoids the cosets of
\(\bar R_0\) obtained from nontrivial roots of unity. Using \(p_j^2\geq M/b\)
for both primes and \(D_{\mathrm{main}}\le (M/b^6)^{kf}D_{\mathrm{id}}\), the
right-hand side of \eqref{eq:span-intersect} is at most
\(b^{-3kf}D_{\mathrm{id}}\). Under \(kf\geq 40\) and \(b\ge 1000\),
\eqref{eq:span-cutoff} gives
\begin{equation}\label{eq:span-bessel-near}
 \sum_{s\in\bar R_0G_i}\|E_s\|_{\mathrm{HS}}^2
 \leq b^{-12}D_{\mathrm{id}}.
\end{equation}
\leanref{secondMoment_compl_besselS_le}

\subsection{Second moments before and after projection}

For \(t\in\bar R_0\), \eqref{eq:span-equivariance} gives
\(\|E_t\|_{\mathrm{HS}}=\|P^*P-1\|_{\mathrm{HS}}\).
Restricting \eqref{eq:span-bessel-near} to \(\bar R_0\) thus yields
\begin{equation}\label{eq:span-gram-isometry}
 \|P^*P-1\|_{\mathrm{HS}}^2\leq b^{-12}D_X.
\end{equation}
For the orthogonal projection \(Q_U\) onto a subspace \(U\subset S_0\),
\[
 \|P Q_U\|_{\mathrm{HS}}^2
 = \dim U+\operatorname{tr}\bigl(Q_U(P^*P-1)\bigr)
 \geq \dim U-b^{-6}D_X.
\]
To estimate the loss from projecting to the target, factor
\(e_0=e_{p_0}e_{p_1}\) and let \(a_i(t)\) be the coefficient of \(t\) in the
complementary projection \(1-e_{p_i}\) in \(\mathbf C[G_i]\).
Counting the types of the finite general linear groups specified by
polynomials and using their degrees gives
\[
 \sum_{t\in G_i}|a_i(t)|^2=a_i(1)\leq \tfrac12 b^{-8}.
\]
The equality follows by comparing coefficients of the identity in the
square of the self-adjoint idempotent. \leanref{bad_types_rare}

Since \(1-e_{p_i}\) is an orthogonal projection,
\[
 \|(1-e_{p_i})P\|_{\mathrm{HS}}^2
  =\operatorname{tr}\sum_{t\in G_i} a_i(t)\bigl(r_\psi(t)+E_t\bigr).
\]
The trace of the Bessel term is \(D_X\) times the reduced unipotent
character. Since \(a_i\) is a class function, this character can be
replaced by its conjugacy average \(\psi_U^\natural\). The bound
\(\sum_{s\in\bar G}|\psi_U^\natural(s)|^2\le2\) and Cauchy--Schwarz
show that the absolute value of this contribution is at most \(b^{-4}D_X\).
\leanref{trace_besselTranslate_eq_mul_reducedChar}

For the remaining term, extend \(s\mapsto E_s\) by zero outside \(\bar R_0G_i\)
and apply \(f(s)\mapsto\sum_t a_i(t)f(st)\). As an orthogonal projection on
\(\ell^2(\bar G)\), this operation does not increase the sum of squares.
At \(s\in\bar R_0\) its value is \(r_\psi(s)\sum_t a_i(t)E_t\) by
\eqref{eq:span-equivariance}. Hence \eqref{eq:span-bessel-near}, after division by
\(|\bar R_0|\), gives \(\|\sum_t a_i(t)E_t\|_{\mathrm{HS}}^2\le b^{-12}D_X\).
Its trace has absolute value at most \(b^{-6}D_X\le b^{-4}D_X\).
Applying this at the two primes and using the orthogonal decomposition
\(1-e_{p_0}e_{p_1}=(1-e_{p_0})+e_{p_0}(1-e_{p_1})\) gives
\(\|(1-e_0)P\|_{\mathrm{HS}}^2\le4b^{-4}D_X\). Therefore
\begin{equation}\label{eq:span-near-isometry}
 \|e_0P Q_U\|_{\mathrm{HS}}^2\geq \dim U-5b^{-4}D_X
       \qquad(U\subset S_0).
\end{equation}
\leanref{isNearIsometry_piTensorElem_comp}
Taking \(U=S_0\) in \eqref{eq:span-near-isometry} gives \(h>0\).

We next estimate the second moment over the full group. Put
\(\beta(s)=\|P_s\|_{\mathrm{HS}}^2\) and \(X_t=\rho(t)PP^*\rho(t)^*\).
Cyclicity of the trace gives
\[
 \langle X_t,X_u\rangle_{\mathrm{HS}}=\beta(t^{-1}u).
\]
Thus \(\beta\) is of positive type, and \(\beta(s)\leq\beta(1)\).
By \eqref{eq:span-gram-isometry}, we also have
\(\beta(1)\le(1+b^{-6})^2D_X\le2D_X\).

Among the finitely many choices of \(A\) and of its specified vectors,
choose one maximizing
\[
 b^{6k|A|}\sum_{s\in\Gamma_A}\beta(s).
\]
Put \(Q=\sum_{s\in\Gamma_A}\beta(s)\). First suppose \(A\) includes all the
chosen places above \(p_0\), and put \(C=b^{-6kf}D_{\mathrm{main}}\).
The union of the scalar cosets of \(\bar R_0\) obtained from integral adeles
has at most \(M^f|\bar R_0|\) elements. Its contribution is at most
\(2 M^fD_{\mathrm{id}}\le C/2\), by the parameter comparison
\(8 M^f D_{\mathrm{id}}\le b^{-6kf}D_{\mathrm{main}}\).
For the remaining elements we can take \(z=p_0^f\).
Since \(p_0\ge b^{10}\), the right-hand side of \eqref{eq:span-intersect}
is at most \(C/8\). The square of the error term in
\eqref{eq:span-cutoff} is also at most \(C/8\), so the square of the sum
of the two square roots is at most \(C/2\). It follows that \(Q\le C\).
\leanref{secondMoment_fixer_univ_le}

Now suppose there is a chosen \(w\notin A\). By maximality, the sum of
\(\beta\) over the elements fixing any specified nonzero frame vector
at \(w\) or its conjugate is at most \(b^{-6k}Q\).
We use this to bound the mass of the singular minors at \(w\).
Let \(\pi\) be the matrix representation of \(\Gamma_A\) over the residue
field of order \(p_0\). Write the \(\ell\)-th leading principal minor as
\(\det Z_s\), where \(Z_s=A_1^{\mathsf T}\pi(s)A_0\). The columns of
\(A_0\) are the first \(\ell\) frame vectors, and the columns of \(A_1\)
represent the corresponding linear forms. Orthogonality for a nontrivial
additive character \(\psi\) gives
\[
 p_0 \sum_{\det Z_s=0}\beta(s)
 \le \sum_{s\in\Gamma_A}\beta(s)\#\{y:y^{\mathsf T}Z_s=0\}
 = \sum_{y\in\mathbf F_{p_0}^\ell}\frac{1}{p_0^\ell}
    \sum_{v\in\mathbf F_{p_0}^\ell}\ \sum_{s\in\Gamma_A}
        \psi(y^{\mathsf T}Z_s v)\beta(s).
\]
The total contribution from \(y=0\) or \(v=0\) is at most \(2Q\).
For \(y,v\ne0\), use the orthonormal basis of point masses and the
orthonormal basis of additive characters. Tensor the point mass at
\(\pi(s)A_0v\) and, separately, the normalized character at
\(\pi(s)^{-\mathsf T}A_1y\) with \(X_s\) and sum each over \(s\).
The squared norm of each sum is \(|\Gamma_A|\) times the sum of \(\beta\)
over the corresponding stabilizer. Expanding the inner product of the two
sums gives the innermost character sum above times \(|\Gamma_A|p_0^{-r/2}\).
Cauchy--Schwarz therefore bounds the real part of that character sum
by \(p_0^{r/2}b^{-6k}Q\). Together with \(p_0^{\ell+r/2}\le b^{6k}\)
this gives
\[
  \sum_{\det Z_s=0}\beta(s)
  \le p_0^{-1}\bigl(2+p_0^{\ell+r/2}b^{-6k}\bigr)Q
  \le 3b^{-10}Q .
\]
Allowing the conjugate place supplies the stabilizer condition for the
linear forms under the contragredient of \(\pi\).
\leanref{card_mul_sum_det_eq_zero_le}
The frame has length at most \(b\). Summing over the minors bounds the
contribution from the complement of the open reductions by \(b^{-4}Q\).
\leanref{sum_not_isOpenElement_le}

The open reductions avoid the scalar cosets, and \(r_\psi(s)=0\) there.
Apply \eqref{eq:span-cutoff} and \eqref{eq:span-open-input} with
\(D_{\mathrm{id}}\le b^{-6k|A|}D_{\mathrm{main}}\) and
\(b^{-15}\le\tfrac12 b^{-4}\). The square root of the contribution of the
open reductions is then at most
\((1+b^{-4})(b^{-6k|A|}D_{\mathrm{main}})^{1/2}\).
Combining the two contributions gives
\[
 Q \leq (1+4b^{-4})b^{-6k|A|}D_{\mathrm{main}}+b^{-4}Q,
 \qquad
 Q\leq (1+9b^{-4})b^{-6k|A|}D_{\mathrm{main}}.
\]
The second inequality follows from \((1-u)(1+9u)\ge1+4u\) at \(u=b^{-4}\).
It holds also when \(A\) consists of all the chosen places above \(p_0\).
Comparing the maximum to the value at \(A=\varnothing\) proves
\begin{equation}\label{eq:span-second-moment}
 \sum_{s\in\bar G}\|P_s\|_{\mathrm{HS}}^2
    \leq (1+9b^{-4})D_{\mathrm{main}}.
\end{equation}
\leanref{sum_secondMoment_le}

\subsection{From the second moment to a rank bound}

Put \(\epsilon=(8b^2)^{-1}\). Combining \eqref{eq:span-dimension-mass}
with \eqref{eq:span-second-moment} cancels the factor
\(\deg_B\,\mathrm{mass}_r\). Since \(9b^{-4}\le\epsilon\), we obtain
\begin{equation}\label{eq:span-moment-comparison}
 \frac{h}{|\bar G|}\sum_{s\in\bar G}\|P_s\|_{\mathrm{HS}}^2
    \leq (1+\epsilon)D_X^2.
\end{equation}
For every subspace \(U\subset S_0\), \eqref{eq:span-near-isometry}
also gives \(\|e_0P Q_U\|_{\mathrm{HS}}^2\geq\dim U-\epsilon D_X\).

Let \(F_1\subset V\) be a subspace and set \(\alpha=\dim F_1/h\).
Write \(Q_{F_1}\) for the orthogonal projection in \(\mathcal H\) onto \(F_1\).
The second moment gives
\begin{align*}
 \frac1{|\bar G|}\sum_{s\in\bar G}\|Q_{F_1}\rho(s)P\|_{\mathrm{HS}}^2
 &\leq
 \left(\dim F_1\cdot
        \frac1{|\bar G|}\sum_{s\in\bar G}\|P_s\|_{\mathrm{HS}}^2\right)^{1/2}\\
 &\leq \sqrt{\alpha(1+\epsilon)}D_X
 \leq \frac{1+\alpha+\epsilon}{2}D_X.
\end{align*}
The first line applies Cauchy--Schwarz to the Hilbert--Schmidt inner
product of \(Q_{F_1}\) with \(\sum_s X_s\). Indeed,
\(\|\sum_s X_s\|_{\mathrm{HS}}^2=|\bar G|\sum_s\beta(s)\).
Some \(s\) consequently attains a value no greater than this average.
Set \(B_s=e_0\rho(s)P\colon S_0\to V\) and \(U=B_s^{-1}(F_1)\).
On \(U\) the map \(B_s\) agrees with \(Q_{F_1}\rho(s)P\). Since \(\rho(s)\)
is unitary and commutes with \(e_0\), we have
\[
 \dim U-\epsilon D_X
 \ \leq\ \|e_0P Q_U\|_{\mathrm{HS}}^2
 \ \leq\ \|Q_{F_1}\rho(s)P\|_{\mathrm{HS}}^2.
\]
Therefore \(\dim U\le(1+\alpha+3\epsilon)D_X/2\).
\leanref{isSpread_of_second_moment}

The map to the quotient \(V/F_1\) has rank \(D_X-\dim U\).
While \(\operatorname{codim}F_1>16\epsilon h\), adjoining
\(\operatorname{im}B_s\) accounts for at least a fraction \(2D_X/(5h)\)
of the missing dimensions, since
\[
 \frac{1-\alpha-3\epsilon}{2}D_X
 \ge \frac25(1-\alpha)D_X
 \qquad (1-\alpha>16\epsilon).
\]
Write \(\lambda=2D_X/(5h)>0\). At \(F_1=0\) the gain is at most \(h\),
so \(\lambda\le1\). Continue choosing translates while the missing fraction
exceeds \(16\epsilon\). At each step the codimension decreases by a
factor at most \(1-\lambda\le e^{-\lambda}\). Within
\(\lceil\lambda^{-1}\log(1/(16\epsilon))\rceil\) steps it is at most
\(16\epsilon h=2h/b^2\). The set \(T\) of chosen translates satisfies
\[
  |T|D_X\le\tfrac52\bigl(1+\log(1/(16\epsilon))\bigr)h
               \le 5h\log(1+2b^2).
\]
\leanref{IsSpread.exists_few_le_finrank_iSup_range}

\subsection{Restriction to the ramified source}

Write \(N=\pi_N\overline{\pi_N}\), and for each \(w\mid\pi_N\) let \(z_w\)
be the finite adelic scalar action chosen in the formalization.
The space \(S\) consists of the elements killed by \(\sum_{j=0}^{N-2}z_w^j\)
for every \(w\). Since \(z_w^{N-1}=1\), the rank of this averaging operator
equals the dimension of the fixed space of \(z_w\).
The level generated by adjoining \(z_w\) has index at least \(N-1\) over
the original level. The class automorphism groups at both levels are
trivial, so the fixed space has dimension at most \(D_X/(N-1)\).
As there are \(f\) places \(w\mid\pi_N\), this gives
\[
 \operatorname{codim}_{S_0} S\le \frac{f}{N-1}D_X\le b^{-8}D_X.
\]
\leanref{finrank_ker_one_sub_scalingGenerator_le}

Restricting each map to \(S\) loses at most the dimension of the image
of a complement of \(S\). Since all the normalizing scalars are nonzero,
\(\operatorname{im}(B_s|_S)=\mathfrak p_sS\). The total loss is at most
\[
   |T|\operatorname{codim}_{S_0}S
   \le 5b^{-8}\log(1+2b^2)h\le 2h/b^2.
\]
Adding the earlier codimension bound \(2h/b^2\) gives a codimension
at most \(4h/b^2\). This and \(\dim S\le D_X\) prove the two inequalities
in Theorem~B. \leanref{enough_forms}

\section{Determinant estimates by twisting}
\label{sec:determinant-estimates}

The Archimedean estimates concern the Gram determinant for the trivial
twist. Twisting the chain allows us to compare this determinant with
another Gram determinant at each finite place. The choice of twist
controls both the denominators of projections onto spaces of theta lifts
and the Hecke \(L\)-values that enter the comparison. These two bounds
give a determinant ratio at each stage and then a ratio for the full chain.

\subsection{Choosing twists and clearing projection denominators}

Set \(r=\dim_E Y_0\) and \(t=5r\), and write \(W_j,V_j\) for the
Hermitian spaces of ranks \(r,r+t\), respectively, used in the comparison
at stage \(j\). Let \(F_j\) be the image of \(e_0\) on the forms of
principal level on \(W_j\). After projection, the multiplicity of each
Hecke eigensystem at split places away from the level is the same on
\(W_j\) and \(Y_0\). In particular, \(\dim_{\mathbb C} F_j=h\).
\leansrc{New/Spanning/Multiplicities/EqualMultiplicities}{Equality of multiplicities}.

A twist \(\xi\) is an automorphic character of \(U(1)\) of weight \(0\) at
the infinite places. Write \(Q_\xi\) for its set of ramified places.
For every twist under consideration, \(\xi^j\) is ramified at every
place of \(Q_\xi\) for \(1\le j\le m\). The modulus norm appearing
at the level is therefore \(N_\xi=\prod_{w\in Q_\xi}\mathrm Nw\) for
every \(\xi^j\). The twist \(\xi\) is admissible at \(\ell\) if \(N_\xi\)
is coprime to \(\ell\).

Consider a character \(\nu\) of weight \(-c(\tau)\) and an integer \(u\)
with \(1\le u\le t\) and \(u\le c(\tau)\le4b\). They form a critical
pair if the tame part of \(\nu\) is trivial away from the chosen level
places. A critical pair is called ramified if its tame part is nontrivial
at every chosen level place. We use the ratio
\[
 \mathcal L_{\nu,u}(\vartheta;\varepsilon)
 =
 \frac{|L(u,(\nu\vartheta^\varepsilon)^{-1})|}
      {|L(u,\nu^{-1})|},
 \qquad \varepsilon\in\{1,-1\}.
\]
Here \(L\) is a Hecke \(L\)-function over \(E\) with character pulled back
from \(U(1)\) along \(z\mapsto z/\bar z\). Both numerator and denominator
omit Euler factors at the chosen level places and their conjugates.
The numerator also omits the factors at \(Q_\vartheta\) and its conjugate.
A main pair for \(j\) is a critical pair with \(m<u\le r\), with tame
part equal to that of the character at stage \(j\), and with
\(w_j(\tau)+r-t\le c(\tau)\le w_j(\tau)\) at the infinite places, where
\(-w_j(\tau)\) is the weight of that stage's character. For a main pair,
we decompose the associated space of rank \(2u\) into four blocks of
ranks \(12,4,4,2u-20\).
\leansrc{New/Determinant/Twist/FitTwists}{Critical pairs and main pairs}.

Write \(u_{j,\vartheta}\) for the theta lift on \(F_j\) and
\(P_{j,\vartheta}\) for the orthogonal projection onto its image
\(I_{j,\vartheta}\). At a split place \(w\) away from the level, specialize
the Hecke polynomial operator at the value associated with the lift at
stage \(j\). The false lifts at \(w\) form the intersection of the kernel
of this operator, the image of \(e_0\), and \(I_{j,\vartheta}^{\perp}\).
Orthogonality to false lifts means orthogonality at every such \(w\).
\leansrc{New/OneLift/TestedForms}{Definition of false lifts}.

Put
\[
 B_{\mathrm{den}}=\frac{E_3}{mk}.
\]
The existence theorem for regular twists supplies, for each \(\iota\),
a twist \(\xi_\iota\) admissible at \(\ell\), with the three properties
below. As before, we write \(|v|_\iota=|\iota(v)|_\ell\).
\begin{enumerate}
\item For every ramified pair, every \(1\le j\le m\), and every
  \(\varepsilon=\pm1\), the family
  \(\bigl(\mathcal L_{\nu,u}(\xi_\iota^j;\varepsilon)^{-1}\bigr)_\iota\)
  has height at most \(x_1\).
\item For each \(j\), there are nonnegative reals \(\beta_\ell\) indexed
  by primes, with \(\sum_{\ell\in T}\beta_\ell\le B_{\mathrm{den}}\) for
  every finite set \(T\) of primes. For \(\vartheta=\xi_\iota^j\), there
  is an endomorphism \(\Lambda\) of \(I_{j,\vartheta}\) such that
  \[
    |\det\Lambda|_\iota\ge\exp(-hr\beta_\ell).
  \]
  The operator \(\Lambda P_{j,\vartheta}\) takes forms that are integral
  at \(\iota\) and orthogonal to the false lifts to forms integral
  at \(\iota\).
\item The same \(\beta_\ell\) controls projection onto the line generated
  by the lift from rank \(1\) associated with a main pair. Write \(P\)
  for the orthogonal projection onto that line. If \(\Xi,\Xi'\) are
  test forms integral at all places above \(\ell\) and orthogonal to
  the false lifts for the rank-one descent, then
  \[
       |\langle P\Xi,P\Xi'\rangle|_\iota\le\exp(\beta_\ell).
  \]
\end{enumerate}

The existence theorem uses the Galois representations of rank \(r\)
attached to the finitely many Hecke eigensystems in the \(e_0\)-part of
the forms on \(W_m\), together with the rank-one representation attached
to each main pair. The twisting theorem is applied to the main summands
whose traces are obtained from those of the twisted lifts by removing a
string of characters. For the first \(m\) powers of one anticyclotomic
character it bounds both the lengths of the Selmer groups with divisible
coefficients and the valuations of the \(L\)-value ratios for ramified pairs.
Let \(C_{\mathrm{lev}}\) be the product of level group orders used in the
formal proof, and \(C_1\) the number of classes of \(U(1)\) at the chosen
principal level. For the inverse ratio of a ramified pair the bound for
\(\log|\cdot|_\iota\) is
\[
  \big((4f+3)r\,v_\ell(C_{\mathrm{lev}})+2^{-\ell}\big)\log\ell.
\]
The denominator estimate obtained by Ribet's method combines the Selmer
group length with local and level corrections in the exponent \(c_\ell\).
\leanref{clearingBound}
The inner product of rank-one projections also requires an allowance of
\(v_\ell(C_1)\log\ell\), so both estimates use
\(\beta_\ell=(c_\ell+v_\ell(C_1))\log\ell\).
The two sums over primes are bounded by \(x_1\) and \(B_{\mathrm{den}}\),
respectively. We use the twisting theorem and the passage from Selmer
group lengths to projection denominators in this form.
\leansrc{New/Congruences/Selmer/TwistingTheorem}{Twisting theorem for Selmer groups},
\leansrc{New/Congruences/FirstOrder/DenominatorBound}{Projection denominator bound},
\leansrc{New/Congruences/FitTwistsExist}{Existence of regular families}.

\subsection{Hecke values and theta derivatives for main pairs}

Fix a main pair \(p=(\nu,u)\) and write \(\phi_\vartheta\) for the normalized
lift from rank \(1\) to rank \(2u\) associated with it. Comparing the norm
formula for the character lift at the two twists gives
\begin{equation}
 \mathcal L_p(\vartheta;1)^{-1}
   =N_\vartheta^{2u-1}\frac{\|\phi_1\|^2}{\|\phi_\vartheta\|^2}.
 \label{eq:main-pair-norm}
\end{equation}
To estimate the ratio at finite places, we project two test forms onto the
line spanned by \(\phi_\vartheta\).

First, consider the height of \(\|\phi_1\|^2\) as a constant family.
Let \(\phi^{(0)}_1\) be the lift corresponding to a critical pair of the
same rank and derivative order \(0\). Its squared norm becomes integral
at every finite place after multiplication by an integer \(n_L\).
Choose \(n_L\|\phi^{(0)}_1\|^2/\|\phi_1\|^2\) as the common multiplier
in the height bound for \(\|\phi_1\|^2\). This makes the finite-place
contribution zero. At infinity, the two logarithmic squared norms have
a common term that cancels. The derivative order for the main pair is
\(c(\tau)-u\ge\kappa(\tau)\), so the contribution from the derivative
orders differs from that for order \(0\) by at least
\(\sum_\tau s_\infty(\kappa(\tau))\). The error in the norm formula
together with \(\log n_L\) is bounded by \(B_{\mathrm{den}}\), giving
\begin{equation}
 \operatorname{ht}\big((\|\phi_1\|^2)_\iota\big)
  \le B_{\mathrm{den}}-\sum_{\tau} s_\infty(\kappa(\tau)).
 \label{eq:untwisted-rankone-height}
\end{equation}
The norm estimate is first proved for conjugations fixing the quadratic
field \(K\). The lemma on absolute values under conjugation extends it
to every field automorphism of \(\mathbb C\).
\leansrc{New/OneLift/HeckeValues/AlgebraicPartSize}{Height of the untwisted norm}.

Write the required finite-place valuation bound for the theta derivative as
\[
 F(\ell)=\frac12\sum_{\tau\in\Phi}\sum_{e\ge 1}
     (1-\chi_K(\ell^e))
     \left(\frac{\kappa(\tau)}{\ell^e}-\kappa(-\ell^e\tau)\right)_+,
\]
where \(\chi_K=\chi_d\) is the quadratic character of \(K\).
Grouping the terms by prime powers gives, for every finite set \(T\)
of primes,
\begin{equation}
 2\sum_{\ell\in T}F(\ell)\log\ell
   \le \sum_{\tau\in\Phi}s_{\mathrm{fin}}(\kappa(\tau),\tau).
 \label{eq:local-prime-total-lower}
\end{equation}
\leansrc{New/UpperBound/Size/Estimates}{Sum of local prime contributions}.

The construction of the test form expands the lift of rank \(2u\)
according to the four blocks above. In the block of rank \(12\), choose a
theta-series coefficient that attains the valuation bound. This supplies
a term of valuation at most \(F(\ell)\) plus the valuation of a fixed
factor. In the second block, of rank \(4\), use the derivative-order-zero
lift of a ramified pair \(q_2\), with ratio
\(\rho_\vartheta=\mathcal L_{q_2}(\vartheta;1)\). The remaining two blocks
are assigned the two factors of a character of the class group of
\(U(1)\). Translate the class representatives so that the trivial class
is represented by the identity. Summing the normalized, character-weighted
class sums over all characters recovers the chosen coefficient by character
orthogonality. The ultrametric inequality then gives at least one class
sum whose \(\iota\)-absolute value is at least that of the coefficient.
Choose a factorization of the corresponding character so that the product
of the algebraic \(L\)-values from the last two blocks has
\(\iota\)-absolute value at least their geometric mean over all possible
factorizations. The product over the blocks defines the test form.
Expanding its pairing with the lift as a class sum yields these factors.

The resulting local estimate uses a scalar
\(c_0\ne0\), independent of the twist and the finite place, with
\(|\sigma(c_0)|\le\exp(fmb/k+3x^{34})\le\exp(B_{\mathrm{den}})\) for every
\(\sigma\). For each \(\iota\) and each \(\vartheta\) admissible at its
prime, there is a test form \(\Xi\) such that
\begin{equation}
 |\langle\Xi,\phi_\vartheta\rangle|_\iota
    \ge |c_0\rho_\vartheta|_\iota\,\ell^{-F(\ell)}.
 \label{eq:test-pairing-lower}
\end{equation}
Here \(c_0\) is the product of the fixed factor multiplying the theta
coefficient, the untwisted algebraic \(L\)-value from the second block,
and the geometric mean from the remaining two blocks. The same construction
makes the test form integral at all places above \(\ell\) and orthogonal
to the false lifts.
\leansrc{New/Derivatives/Sharpness/BoundAttained}{Theta coefficient valuations},
\leansrc{New/OneLift/HeckeValues/TwistedValueHeights}{Construction of the test form}.

Take \(\vartheta=\xi_\iota^j\). Choose test forms \(\Xi,\Xi'\)
as in \eqref{eq:test-pairing-lower} at \(\iota\) and
\(\iota\circ\mathrm{conj}\), respectively. Since \(\rho_\vartheta\)
is real, the product
\(A_\iota=\langle\Xi,\phi_\vartheta\rangle\langle\phi_\vartheta,\Xi'\rangle\)
satisfies
\[
 \left|c_0\bar c_0\,\frac{\rho_\vartheta^2}{A_\iota}\right|_\iota
      \le \ell^{2F(\ell)}.
\]
Both pairings are nonzero by \eqref{eq:test-pairing-lower}.
The orthogonal projection \(P\) onto \(\mathbb C\phi_\vartheta\) gives
\(\langle P\Xi,P\Xi'\rangle=A_\iota/\|\phi_\vartheta\|^2\).
Thus the right-hand side of \eqref{eq:main-pair-norm} is \(N_\vartheta^{2u-1}\)
times the product of four factors
\[
 \|\phi_1\|^2,\qquad
 \langle P\Xi,P\Xi'\rangle,\qquad
 \rho_\vartheta^{-2},\qquad
 \frac{\rho_\vartheta^2}{A_\iota}.
\]
The height bound for each factor follows, in order, from
\eqref{eq:untwisted-rankone-height}, the estimate for the projected
pairing, the inverse-ratio bound for ramified pairs, and
\eqref{eq:local-prime-total-lower} with common multiplier \(c_0\bar c_0\).
The four bounds are
\[
 \begin{gathered}
 B_{\mathrm{den}}-\sum_\tau s_\infty(\kappa(\tau)),\qquad
 B_{\mathrm{den}},\qquad 2x_1,\\
 2B_{\mathrm{den}}+\sum_\tau s_{\mathrm{fin}}(\kappa(\tau),\tau).
 \end{gathered}
\]
Moreover, \(N_\vartheta^{2u-1}\) is a unit at \(\iota\).
Adding the bounds and using \(x_1\le B_{\mathrm{den}}\) gives
\begin{equation}
 \operatorname{ht}\bigl(
   (\mathcal L_p(\xi_\iota^j;1)^{-1})_\iota\bigr)
 \le 6B_{\mathrm{den}}-A_0.
 \label{eq:main-pair-height-lower}
\end{equation}
\leansrc{New/OneLift/HeckeValues/TwistedValueHeights}{Height of the inverse ratio for a main pair}.

\subsection{Comparing determinants at a single stage}

Set \(Z_j(\vartheta)=\det(u_{j,\vartheta}^*u_{j,\vartheta})\). The
comparison operator \(U_{j,\vartheta}\) is a scalar multiple of
\(G_{j,\vartheta}G_{j,\vartheta}^*\), where \(G_{j,\vartheta}\) glues
tuples of forms on an auxiliary space into forms on \(V_j\). The scalar
is an integer from the level times the square of a hook number.
Three exchanges of theta integrals give a positive-definite operator
\(U_j^\circ\) on \(F_j\), independent of \(\vartheta\), such that
\begin{equation}
 u_{j,\vartheta}^*U_{j,\vartheta}u_{j,\vartheta}
      =\rho_j(\vartheta)U_j^\circ,\qquad
 \rho_j(\vartheta)=
 \frac{\mathcal H_j(\vartheta)}{\mathcal H_j(1)}.
 \label{eq:comparison-operator-lower}
\end{equation}
Here \(\mathcal H_j\) is the product of the three Hecke values that enter
the exchanges. \leanref{heckeValueProd}
The operator \(U_j^\circ\) is a positive scalar multiple of the Gram
operator of the untwisted period map. This map is injective, so
\(\det U_j^\circ\ne0\).

The ratio \(\rho_j(\vartheta)\) is nonzero. The operator \(U_{j,\vartheta}\)
carries theta lifts into forms orthogonal to the false lifts and preserves
integrality at primes coprime to \(N_\vartheta\). For the second property,
a tight frame of integral test vectors expresses
\(G_{j,\vartheta}G_{j,\vartheta}^*\) as a sum of rank-one operators.
The level integer clears the denominators of their pairings.
The Cauchy--Schwarz estimate for the same frame shows that the untwisted
operator increases the norm by at most \(\exp(3fmb(\log b)^2)\) at every
conjugate.
\leansrc{New/OneLift/ComparisonOperator/ThreeExchanges}{The three exchanges},
\leansrc{New/OneLift/LiftVolumeHeight}{The comparison operator and its bounds}.

Let \(B_{j,\vartheta}\) be the restriction of
\(P_{j,\vartheta}U_{j,\vartheta}\) to \(I_{j,\vartheta}\) and put
\(D_j(\vartheta)=\det B_{j,\vartheta}\). The lift \(u_{j,\vartheta}\)
is an isomorphism onto its image, and \(u_{j,\vartheta}^*\) vanishes
on \(I_{j,\vartheta}^\perp\). Hence
\[
 u_{j,\vartheta}^* U_{j,\vartheta} u_{j,\vartheta}
  =(u_{j,\vartheta}^*u_{j,\vartheta})
    (u_{j,\vartheta}^{-1} B_{j,\vartheta}u_{j,\vartheta}).
\]
Taking determinants and using \eqref{eq:comparison-operator-lower} gives
\[
 Z_j(\vartheta)D_j(\vartheta)
      =\det(U_j^\circ)\rho_j(\vartheta)^h,\qquad
 \frac{Z_j(1)}{Z_j(\vartheta)}
   =\frac{D_j(\vartheta)}{D_j(1)}\,\rho_j(\vartheta)^{-h}.
\]

At finite places, for \(\vartheta=\xi_\iota^j\) the operator
\(\Lambda B_{j,\vartheta}\) preserves the integral forms in
\(I_{j,\vartheta}\). Normalize each eigenvector so that it is integral
and rescaling it by a scalar with \(\iota\)-absolute value greater than
\(1\) takes it outside the integral module. The corresponding eigenvalue
has \(\iota\)-absolute value at most \(1\). Consequently,
\[
 |D_j(\xi_\iota^j)|_\iota
   \le|\det\Lambda|_\iota^{-1}
   \le\exp(hr\beta_\ell).
\]
At infinity, \(U_{j,1}(I_{j,1})\) is orthogonal to the false lifts.
Projecting forms in this image back onto \(I_{j,1}\) cannot increase
their norm at any conjugate. The projection estimate uses the adjoint
relation for the specializations of the Hecke operators under conjugation.
Put \(C_{\mathrm{arch}}=fmb(\log b)^2\). Multiplying the bounds for the
eigenvalues gives \(|\sigma(D_j(1))|\le\exp(3h C_{\mathrm{arch}})\).
Using \(D_j(1)\) as common multiplier gives
\[
 \operatorname{ht}\bigl((D_j(\xi_\iota^j)/D_j(1))_\iota\bigr)
     \le h(3C_{\mathrm{arch}}+rB_{\mathrm{den}}).
\]
\leansrc{New/OneLift/LiftVolumeHeight}{Finite and Archimedean bounds for the compressed determinant}.

Up to an integer power of \(N_\vartheta\), the ratio \(\rho_j(\vartheta)\)
in the three exchanges is the square of the product of the ratios for
\(k/2\) main pairs and \(5m\) ramified pairs. The main pairs occur at points
\(7m/2<u\le 7m/2+k/2\). The remaining points number \(7m/2+m+m/2=5m\).
The sign is \(\varepsilon=1\) for the main pairs and for the first \(7m/2\)
ramified pairs. Among the flat pairs it is \(\varepsilon=1\) for \(m\) of
them and \(\varepsilon=-1\) for \(m/2\) of them.
Adding the bounds in \eqref{eq:main-pair-height-lower} and those for the
ramified pairs, with the factor of two from the square, gives
\[
 \operatorname{ht}\bigl((\rho_j(\xi_\iota^j)^{-1})_\iota\bigr)
  \le k(6B_{\mathrm{den}}-A_0)+10mx_1.
\]
The admissible data satisfy
\(3C_{\mathrm{arch}}+(r+6k)B_{\mathrm{den}}+10mx_1
        \le kB_{\mathrm{den}}\log b\), so for every stage
\begin{equation}
 \operatorname{ht}\bigl((Z_j(1)/Z_j(\xi_\iota^j))_\iota\bigr)
   \le h\big(E_4/m-kA_0\big).
 \label{eq:single-relative-height}
\end{equation}
\leansrc{New/OneLift/HeckeValues/TwistedValueHeights}{Factorization of Hecke values},
\leansrc{New/OneLift/LiftVolumeHeight}{Height bound for a single stage}.

\subsection{The determinant of the chain}

To compute the norm of the chain, arrange in a triangle the lifts that
adjoin blocks of height \(t\) to the initial space. The vertex \((i,n)\)
consists of the common part and blocks \(i+1,\ldots,i+n\). The arrows
that add a block on the right are paired with transposed lifts that
remove a block on the left. Comparing the two paths obtained by switching
the order of these operations uses the exchange identity.

The matching over a finite ring works as follows. Split the row and column
indices into two groups. For an invertible matrix \(g\), write \(C(g)\)
for its first columns and \(R(g)\) for the last rows of \(g^{-1}\), with
the sizes specified by the split. The weights \(\Psi,\Phi\) are supported
on matrices of these respective forms and satisfy
\(\Phi(R(g))=\chi(\det g)\Psi(C(g))\).
Let \(K,K'\) be the kernels along the two paths, subject to the formal
proof's weight, ramification, signature, and rank conditions.
The exchange gives \(\sum\Psi K'=(\kappa_0/c)\sum\Phi K\).
The value \(|c|^2\) is determined by the product of powers of residue-field
cardinalities and squared absolute values of partial \(L\)-values.
The scalar \(\kappa_0\) is independent of the twist.
Applying the identity to point data and to their finite Fourier transforms
gives the two exchange laws needed for the triangle.
\leansrc{New/Determinant/Exchange/TheoremForPartners}{Exchange of theta integrals},
\leansrc{New/Determinant/Exchange/OneSquare}{Exchange across one square},
\leansrc{New/Determinant/Chain/ChainNorms}{Exchange laws and injectivity for the triangle}.

Write \(c_Q(\xi)\) for
\[
 N_\xi^{rt}\prod_{w\in Q_\xi}\prod_{t<s\le t+r}(1-(\mathrm Nw)^{-s}).
\]
Let \(R_j(\xi^j)\) be the Hecke operator on \(V\) from Rallis's formula
at stage \(j\). Applying Rallis's formula to a single lift and matching
the multiplicities of the Hecke eigensystems of \(W_j\) and \(Y_0\)
gives, for a constant \(A_j\ne0\) independent of the twist,
\[
 Z_j(\xi^j)=(A_jc_Q(\xi))^h\det R_j(\xi^j).
\]
For the chain, the exchange laws move the blocks one at a time.
Write \(X_{i,n}\) for the composition of the first \(n\) arrows in row
\(i\), weighting the arrow from stage \(s\) to \(s+1\) by the order
\(w_s\) of the general linear group at its source's finite level.
Let \(a_{i,n}\) be the arrow to the right and \(d_{i,n}\) the downward
arrow to the next row. The map \(j_i\) embeds the space of principal-level
forms isometrically at the beginning of row \(i\). Define \(x'\) by
\(j_{i+1}x'=d_{i,0}a_{i,0}j_i x\).
Composing the exchange laws for individual squares and using the
exchange law for adjoints gives
\begin{multline*}
  \langle X_{i,n+1}j_i x,X_{i,n+1}j_i y\rangle\\
   = \lambda_{i,n}\,\frac{w_0w_1w_n}{w_{n+1}}\,
     \langle X_{i+1,n}j_{i+1}x',X_{i+1,n}j_{i+1}y\rangle.
\end{multline*}
Up to a factor independent of the twist and a half-integral power of
\(N_\xi\), the absolute value of \(\lambda_{i,n}\) is
\(\prod_{i+1<j\le i+n+1}J_{i+1,j}(\xi)\).
On the right, the pairing with \(x'\) becomes a pairing with \(x\)
by applying the Gram operator of the first arrow in row \(i\),
canceling one \(w_0\). Starting the induction with the empty
composition, the Gram operator for row \(i\) is a scalar multiple of the
product of the Gram operators of the first arrows in the successive rows.
The scalar includes \(w_0w_1^n/w_n\) and the contributions from pairs
\(i<r'<j\le m\).
\leansrc{New/Determinant/Chain/TriangleNorms}{Recurrence for the weighted chain}.

For the chain starting in row \(0\), apply Rallis's formula to the
Gram operator of the first arrow in each row. Expressing the final level
index by a count of block upper triangular matrices cancels
\(w_0w_1^m/w_m\). With this normalization the Gram operator is
\[
 \Theta_\xi^*\Theta_\xi
       = c(\xi)c_Q(\xi)^m \prod_{j=1}^m R_j(\xi^j),
\]
where
\[
 |c(\xi)| = C\,N_\xi^{e/2}\prod_{1\le i<j\le m} J_{ij}(\xi),
 \qquad C>0,\quad e\in\mathbb Z.
\]
Write \(\eta_j=\mu_j^0\) for the untwisted character at stage \(j\).
Then
\[
 J_{ij}(\xi)=
   \prod_{u=1}^{t}\mathcal L_{\eta_j/\eta_i,u}(\xi^{j-i};1).
\]
Every pair \(i<j\) arises in moving the blocks. The tame part of
\(\eta_j/\eta_i\) at the chosen level is the \((j-i)\)-th power of a
fixed character. Since \(1\le j-i\le m\), the pairs in this product are
ramified.

Injectivity of the chain follows from that of the weighted chain in the
triangle. The exchange laws imply that when the weighted chain of length
\(n+1\) kills \(j_i v\), the chain of length \(n\) in the next row kills
\(d_{i,0}a_{i,0}j_i v\). This vector lies in the image of \(j_{i+1}\).
Induction on the length and injectivity of \(d_{i,0}a_{i,0}\) on the
image of \(j_i\) give \(v=0\). The positive level index and the positive
dimension of the Weyl module that enter the change of normalization
transfer this injectivity to \(\Theta_\xi\). Each individual lift is
also injective, so all the Gram determinants are positive. Taking
absolute values on both sides of the determinant identities retains
these positive Gram determinants and cancels \(c_Q(\xi)\) and the
determinants of the Hecke operators. Thus
\begin{equation}
 \frac{Z(1)}{Z(\xi)}
 =
 \left(N_\xi^{e/2}\prod_{1\le i<j\le m}J_{ij}(\xi)\right)^{-h}
       \prod_{j=1}^m\frac{Z_j(1)}{Z_j(\xi^j)}.
 \label{eq:chain-relative-factorization}
\end{equation}
If \(\ell\nmid N_\xi\), then \(N_\xi^{e/2}\) has \(\iota\)-absolute value
\(1\). Counting the ramified pairs gives height at most \(m^2t x_1\) for
the inverse of the parenthesized factor evaluated at \(\xi_\iota\), and
\(m^2t x_1\le b\). Adding \eqref{eq:single-relative-height} over the
\(m\) stages now gives
\begin{equation}
 \operatorname{ht}\bigl((Z(1)/Z(\xi_\iota))_\iota\bigr)
       \le h(b+E_4-H_0).
 \label{eq:chain-relative-height}
\end{equation}
\leansrc{New/Determinant/Chain/NormOfChain}{The chain's Gram operator and determinant factorization}.

\subsection{Forms for completing a basis}

To complete a basis, we span the space of principal-level forms before
applying \(e_0\) using forms of deep level. The formal proof uses two
grids. It transfers the period of the initial path along their boundary
to the period of the final path, making one turn in each of the \(m^2\)
squares. At each turn it applies the exchange identity to a sum over
finite matrix data. At a turn on the diagonal it first identifies the
relation to a rank-one kernel and then applies the same comparison.
The resulting nonzero scalar \(c_{\mathrm{path}}(\xi)\) satisfies
\begin{equation}
 \left|\frac{c_{\mathrm{path}}(\xi)}{c_{\mathrm{path}}(1)}\right|
   =N_\xi^{e'/2}\prod_{0\le i,j<m} T_{ij}(\xi),\qquad
 T_{ij}(\xi)=
    \prod_{u=1}^{t}\mathcal L_{\eta_{i+1}\eta_{j+1},u}(\xi^{j+1};1).
 \label{eq:turn-ratio-lower}
\end{equation}
The exponent of the tame character is \(i+j+2\), between \(2\) and \(2m\),
so the tame part is nontrivial at every chosen level place. Thus the
pairs in each \(T_{ij}\) are ramified.
\leansrc{New/Determinant/AuxiliaryForms/Turns}{The \(L\)-value ratio for a turn},
\leansrc{New/Determinant/AuxiliaryForms/DiagonalTurns}{Comparison of the full paths}.

The final path period sums the inner product of the lifts along the two
half-paths over tuples \(S\) of finite matrices, with coefficient \(\Psi(S)\).
With the normalization fixed, an integer independent of \(\xi\) clears the
denominators of the kernels along the edges and of the class sums when
\(\ell\nmid N_\xi\). For the trivial twist, the Weyl norm bounds for the
kernels and the number of summands give a bound for every conjugate of
the period. To recover the values of the original lift, pair it with dual
delta forms. For the matrix \(D\) of the dual pairing there are an integral
matrix \(C\) and a positive hook number \(H\) with \(CD=HI\).
If the pairing against each integral vector lies in \(H\) times the
valuation ring, testing against the rows of \(C\) and canceling \(H\) shows
that each coordinate of the original value is integral. At infinity, dividing
the bound for the pairing by the norm of the dual vector gives the Weyl
norm bound.

This gives a quantitative spanning result for the principal-level space.
There is a surjection \(s\) from the deep-level space to the principal-level
space and a constant \(c_{\mathrm{sc}}\ne0\). The deep-level space is
spanned by forms \(\psi\) integral at every finite place and with all
conjugates of their values of squared Weyl norm at most \(\exp(x^{99}b^2)\).
For each \(\psi\), scaling the lift of \(c_{\mathrm{sc}}s(\psi)\) by
\(c_{\mathrm{path}}(\xi)/c_{\mathrm{path}}(1)\) makes it integral at
admissible places. For the trivial twist, the squared Weyl norms of the
lift's values are bounded by the deep-level bound times
\(\exp(2x^{201}b^2)\), hence by \(\exp(x^{202}b^2)\).
These are the local and Archimedean kernel estimates used in the
construction.
\leansrc{New/Determinant/AuxiliaryForms/SpanningForms}{The final path period and construction of auxiliary forms}.

Send this spanning family to \(V\) by the projection multiplied by
the order of the level group. The scaled projection preserves
integrality and adds at most \(x^{99}\) to the logarithm of the
squared norm for the trivial twist. Since
\(x^{99}+x^{202}b^2\le2x_1^2b^2\), the resulting auxiliary forms \(z\)
satisfy
\[
 \|\Theta_1 z\|_\sigma^2\le \exp(2x_1^2b^2)
 \qquad(\sigma\in\operatorname{Aut}(\mathbb C)).
\]
Set \(\gamma(\xi)=c_{\mathrm{path}}(\xi)/c_{\mathrm{path}}(1)\).
Then \(\gamma(\xi)\Theta_\xi z\) is integral at every admissible place.
Equation \eqref{eq:turn-ratio-lower} and the bounds for ramified pairs give
\begin{equation}
 \operatorname{ht}\bigl((|\gamma(\xi_\iota)|^{-2})_\iota\bigr)
 \le 2m^2t x_1\le b.
 \label{eq:auxiliary-factor-height}
\end{equation}
\leansrc{New/Determinant/Twist/LongForms}{Bounds for auxiliary forms after projection to the target}.

\begin{proposition}\label{prop:determinant-inputs}
For all admissible data, \(\Theta_\xi\) is injective for each twist \(\xi\).
There exist a twist \(\xi_\iota\) admissible at the prime of each \(\iota\)
and constants \(\gamma(\xi)\ne0\), with \(\gamma(1)=1\), such that
\[
 \operatorname{ht}\bigl((Z(1)/Z(\xi_\iota))_\iota\bigr)
       \le h(b+E_4-H_0),\qquad
 \operatorname{ht}\bigl((|\gamma(\xi_\iota)|^{-2})_\iota\bigr)\le b.
\]
The space \(V\) is spanned by auxiliary forms \(z\) for which
\(\gamma(\xi)\Theta_\xi z\) is integral at every admissible place and
\(\|\Theta_1 z\|_\sigma^2\le\exp(2x_1^2b^2)\) for every \(\sigma\).
\end{proposition}

\begin{proof}
The comparison of lifts in the triangle gives injectivity.
For a regular family, \eqref{eq:single-relative-height} holds at
each stage. Substituting those bounds into
\eqref{eq:chain-relative-factorization} gives
\eqref{eq:chain-relative-height}. The surjection from deep level and
the projection by \(e_0\) give the spanning family of auxiliary forms,
with \eqref{eq:auxiliary-factor-height} for the same scalars \(\gamma(\xi)\).
\end{proof}

\section{The comparison inequality and the main theorem}\label{sec:conclusion}

\begin{proposition}[Lower bound for squared covolume and index]\label{prop:lower-bound}
Let \(y=(y_1,\ldots,y_n)\) be linearly independent in \(V\), with
\(n\ge(1-4/b^2)h\). If \(c\) bounds the logarithm of its squared
covolume and \(\delta\) bounds its logarithmic index, then
\begin{equation}\label{eq:lower-bound}
  h(H_0-3E_4)\le c+2\delta .
\end{equation}
\end{proposition}

\begin{proof}
Extend \(y\) to a basis using \(n'=h-n\) auxiliary forms from
Proposition~\ref{prop:determinant-inputs}. Use the twists \(\xi_\iota\)
and the scalars \(\gamma(\xi)\) of that proposition in
Proposition~\ref{prop:gram-product}, with \(B_{\mathrm{Ar}}=2x_1^2b^2\),
\(B_{\mathrm{fin}}=b\), and \(B_Z=h(b+E_4-H_0)\). This gives
\[
 0\le c+2\delta+n'(2x_1^2b^2+b)+h(b+E_4-H_0).
\]
Since \(n'\le4h/b^2\) and admissibility gives
\(4(2x_1^2+b^{-1})\le b\), the contribution of the extra forms is
\(n'(2x_1^2b^2+b)\le hb\).
Using \(b\le E_4\) to absorb both contributions of \(hb\) yields
\eqref{eq:lower-bound}.
\leanlinks{\leanref{integral_forms_long}}
\end{proof}

\begin{theorem}[The comparison inequality]\label{thm:main-inequality}
For all admissible data,
\begin{equation}\label{eq:main-ineq}
 m\mathcal R\le fa(\log b)^5 .
\end{equation}
\end{theorem}

\begin{proof}
Take the set \(T\) of translates from Theorem~\ref{thm:spanning}.
Since the short integral forms in Theorem~\ref{thm:upper-bound}(1)
span \(S\), we can choose a basis \(y_1,\ldots,y_n\) of
\(\sum_{g\in T}\mathfrak p_g S\) consisting of images \(\mathfrak p_g\phi\)
of these forms. Then \((1-4/b^2)h\le n\le h\) and \(h>0\).
The bound on the squared norms of the images and Hadamard's inequality
give the bound
\[
   c=n(H_0-m\mathcal R+E_4)
\]
for the logarithm of the squared covolume.
Theorem~\ref{thm:upper-bound}(2) gives a logarithmic index bound
\[
  \delta=2|T|(\dim S)E_3\le 10h\log(1+2b^2)E_3.
\]
Proposition~\ref{prop:lower-bound} and this last bound give
\[
 h(H_0-3E_4)
  \le n(H_0-m\mathcal R+E_4)+20h\log(1+2b^2)E_3.
\]
Since \(H_0+E_4\ge0\), we have \(n(H_0+E_4)\le h(H_0+E_4)\).
With this substitution, the \(hH_0\) terms cancel, leaving
\[
 n\,m\mathcal R
   \le h\bigl(4E_4+20\log(1+2b^2)E_3\bigr).
\]
Admissibility gives \(4/b^2\le1/2\), so that \(n\ge h/2>0\).
Twice the quantity in parentheses is at most \(fa(\log b)^5\).
Dividing by \(n\) gives \eqref{eq:main-ineq}.
\leanlinks{\leanref{star_inequality}}
\end{proof}

\begin{corollary}[The one-sided prime-sum bound]\label{cor:prime-bound}
Let \(q_0\) be a positive multiple of \(18480\), and let
\(d\equiv q_0/2-1\pmod{q_0}\) be prime.
For every \(\varepsilon>0\) and every sufficiently large real \(y\),
\begin{equation}\label{eq:prime-bound}
 R_{q_0}(y)\le y^{-1/2+\varepsilon}.
\end{equation}
\end{corollary}

\begin{proof}
Substitute \eqref{eq:total-remainder} into \eqref{eq:main-ineq} and
divide by \(ma|\Phi_b|\).
Since \(f=\varphi(q_0)|\Phi_a|\) and \(|\Phi_a|/|\Phi_b|\le\log b\),
\begin{equation}\label{eq:remainder-intermediate}
 R_{q_0}(x)
 \le\frac{\varphi(q_0)(\log b)^6}{m}
       +|R_{q_0}(x_1)|+\log b\,|R_1(x_1)|.
\end{equation}
Choose \(0<\eta<1/100\) with \(7\eta\le\varepsilon\) and apply
Proposition~\ref{prop:admissible-data} with \(\eta\).
For sufficiently large integers \(x\), we have
\(\log b\le2x^\eta\), \(m\ge\sqrt{x}\), and
\(\log x_1\ge x^{\eta/12}\).
Taking \(B\ge24/\eta\) in Lemma~\ref{lem:fixed-modulus-pnt} bounds
\(|R_{q_0}(x_1)|\) and \(|R_1(x_1)|\) by \(x^{-2}\).
Thus the right side of \eqref{eq:remainder-intermediate} is at most
\[
 64\varphi(q_0)x^{-1/2+6\eta}+x^{-2}+2x^{-2+\eta}
 \le (64\varphi(q_0)+3)x^{-1/2+6\eta}.
\]
For \(x=\lfloor y\rfloor\), the difference estimate in
Lemma~\ref{lem:fixed-modulus-pnt} gives \(R_{q_0}(y)-R_{q_0}(x)\le26/y\).
Since \(\theta=-1/2+6\eta\) belongs to \([-1,0]\) and \(x\ge y/2\),
we have \(x^\theta\le2y^\theta\). Consequently
\[
 R_{q_0}(y)\le (2(64\varphi(q_0)+3)+26)y^\theta .
\]
For large \(y\) we can absorb the constant into \(y^\eta\), giving
\eqref{eq:prime-bound} by \(7\eta\le\varepsilon\).
\leanlinks{\leanref{eventually_primeRemainder_le_rpow}}
\end{proof}

\begin{proof}[Proof of Theorem~\ref{thm:grh}]
\textup{(1)}
Let \(\chi\) be any character modulo \(q\) and suppose
\(\RePart\rho>1/2\). At \(\rho=1\) the principal character is excluded,
and for nonprincipal characters Dirichlet's theorem gives
\(L(1,\chi)\ne0\). Hence assume \(\rho\ne1\).
Set \(q_0=18480q\) and \(d_0=q_0/2-1\).
We have \(d_0\equiv7\pmod8\) and \((q_0,d_0)=1\), since
\(q_0=2(d_0+1)\) and \(d_0\) is odd.
Choose \(1/2<\beta<\RePart\rho\).
Corollary~\ref{cor:prime-bound} with \(\varepsilon=\beta-1/2\)
gives the hypothesis of Proposition~\ref{prop:one-sided-criterion}
for every prime in the residue class \(d_0\) modulo \(q_0\).
Let \(\chi_0\) be the character modulo \(q_0\) induced by \(\chi\).
Then \(L(\rho,\chi_0)\ne0\). The induction formula
\[
 L(s,\chi_0)=L(s,\chi)
  \prod_{\substack{p\mid q_0\\p\nmid q}}\bigl(1-\chi(p)p^{-s}\bigr)
 \qquad(s\ne1)
\]
yields \(L(\rho,\chi)\ne0\).

\textup{(2)}
Suppose \(L(s,\chi)=0\) with \(0<\RePart s<1\).
Let \(\psi\) be the primitive character inducing \(\chi\).
The extra Euler factors in the induction formula are nonzero
in this strip, so \(L(s,\psi)=0\).
The gamma factor is holomorphic and nonzero here, giving a zero
of the completed \(L\)-function.
The functional equation for \(\psi\) gives a zero of the completed
\(L\)-function of \(\psi^{-1}\) at \(1-s\).
Removing the gamma factor gives \(L(1-s,\psi^{-1})=0\).
Since \(\psi^{-1}\) induces \(\chi^{-1}\), it follows that
\(L(1-s,\chi^{-1})=0\).
\leanlinks{\leansrc{Lib/Analytic/ZeroSymmetry}{Reflection lemma for possibly imprimitive characters}}
Applying (1) to the two zeros gives
\(\RePart s\le1/2\) and \(1-\RePart s\le1/2\).

\textup{(3)}
Suppose \(\zeta(s)=0\), excluding the trivial zeros.
The value \(\zeta(0)=-1/2\) excludes \(s=0\), and \(s=1\) is the pole.
The factor \(\pi^{-s/2}\Gamma(s/2)\) is finite and nonzero at \(s\),
so the completed zeta function vanishes there. The completed function
is invariant under \(s\mapsto1-s\). Since \(1-s\ne0,1\), multiplying
its value there by the reciprocal gamma factor gives
\(\zeta(1-s)=0\).
Apply (1) modulo \(1\) to both zeros to conclude that
\(\RePart s=1/2\).
\leanlinks{\leanref{grh}}
\end{proof}


\begin{thebibliography}{9}
\bibitem{alpoege}
Levent Alpöge, \emph{Announcement of a formal proof of the generalized
Riemann hypothesis}, 2026.
\href{https://alpo.ge/remark.pdf}{PDF}.

\bibitem{grh-source}
Levent Alpöge with Claude, \emph{GRH for Dirichlet L-functions, in Lean 4},
readable edition by Luke Sernau and collaborators.
\href{\upstream/grh/README.md}{README},
\href{\upstream/grh/GUIDE.md}{reading guide}, and
\href{\upstream/grh/REFERENCES.md}{references}.
Commit \texttt{72fead102f66}.

\bibitem{mv}
Hugh L. Montgomery and Robert C. Vaughan,
\emph{Multiplicative Number Theory I. Classical Theory},
Cambridge Studies in Advanced Mathematics 97,
Cambridge University Press, 2007.
\href{https://doi.org/10.1017/CBO9780511618314}{DOI}.

\bibitem{slijepcevic}
Siniša Slijepčević,
\emph{Outline of Alpöge's proof of the generalised Riemann hypothesis},
21 November 2026.
\href{https://hexagonmath.org/pdf/2611.00812}{PDF}.

\bibitem{zhang}
Robin Zhang,
\emph{A preliminary note on Alpöge's proposed proof of the generalized Riemann hypothesis},
22 November 2026.
\href{https://www.ihes.fr/~rzhang/files/GRH\%20note.pdf}{PDF}.

\bibitem{rubinstein}
Michael O. Rubinstein,
\emph{Alpöge/Claude proof of the Generalized Riemann Hypothesis},
\href{https://math.uwaterloo.ca/~mrubinst/alpoge_claude_grh/}{project page}.
\href{https://math.uwaterloo.ca/~mrubinst/alpoge_claude_grh/zeta-paper_v2.pdf}{Exposition for the zeta function},
version 2, working draft dated 22 November 2026.

\end{thebibliography}

\paperlanguage{}{\clearpage}
\section*{\paperlanguage{底本の Lean 宣言との対応}{Correspondence with the Lean source}}
\addcontentsline{toc}{section}{\paperlanguage{底本の Lean 宣言との対応}{Correspondence with the Lean source}}
\paperlanguage{本文の参照番号は、底本のコミットにある以下の宣言を指す。名前から該当する行を開ける。}{The numbered references in the text point to the following declarations in the pinned Lean source. Each identifier links to its declaration.}

\subsection*{\paperlanguage{解析}{Analytic reduction}}
\begin{description}
\leanentry{integral_primeRemainder}
  {剰余のラプラス変換と留数}
  {Laplace transform and residues of the remainder}
  {\leandecl{New/Reduction/Zeros/PrimeError}{293}{integral\_\allowbreak primeRemainder}\enspace \leandecl{New/Reduction/Zeros/PrimeError}{414}{isPolar\_\allowbreak primeRemainder}}
\leanentry{eventually_rpow_lt_thetaTailIntegral}
  {零点がある場合のテータ積分の下界}
  {Lower bound on the theta integral in the presence of a zero}
  {\leandecl{Lib/Analytic/QuadraticTwist}{275}{eventually\_\allowbreak rpow\_\allowbreak lt\_\allowbreak thetaTailIntegral}}
\leanentry{sum_thetaTailIntegral_le}
  {テータ積分の総和の上界}
  {Upper bound on the sum of the theta integrals}
  {\leandecl{Lib/Analytic/QuadraticTwist}{615}{sum\_\allowbreak thetaTailIntegral\_\allowbreak le}}
\leanentry{zero_density_quadTwist}
  {二次捻りの零点密度と補助素数の選択}
  {Zero density for quadratic twists and the choice of auxiliary prime}
  {\leandecl{Lib/Analytic/QuadraticTwist}{732}{zero\_\allowbreak density\_\allowbreak quadTwist}\enspace \leandecl{Lib/Analytic/QuadraticTwist}{796}{exists\_\allowbreak entireL\_\allowbreak quadTwist\_\allowbreak ne\_\allowbreak zero}}
\leanentry{IsPolar.landau}
  {片側評価と留数の符号による極の排除}
  {Exclusion of poles by a one-sided bound and the signs of residues}
  {\leandecl{Lib/Analytic/Landau}{217}{IsPolar.landau}}
\leanentry{LFunction_ne_zero_of_primeRemainder_le}
  {剰余の片側評価による零点判定}
  {Zero-free criterion from a one-sided bound on the remainder}
  {\leandecl{New/Reduction/Zeros/Criterion}{37}{LFunction\_\allowbreak ne\_\allowbreak zero\_\allowbreak of\_\allowbreak primeRemainder\_\allowbreak le}}
\leanentry{eventually_norm_logDeriv_entireL_le}
  {ディリクレ L 関数の古典的な零点のない領域}
  {Classical zero-free region for Dirichlet L-functions}
  {\leandecl{Lib/Analytic/PrimeNumberTheorem/ZeroFreeRegion}{196}{eventually\_\allowbreak norm\_\allowbreak logDeriv\_\allowbreak entireL\_\allowbreak le}}
\leanentry{eventually_norm_logDerivMean_div_le}
  {固定法における対数微分の評価}
  {Fixed-modulus bounds for logarithmic derivatives}
  {\leandecl{New/Reduction/Zeros/PrimeNumberTheorem}{178}{eventually\_\allowbreak norm\_\allowbreak logDerivMean\_\allowbreak div\_\allowbreak le}}
\leanentry{eventually_abs_le_inv_pow_of_transform}
  {打切り点を大きく取ったときの剰余}
  {Remainder bound at the larger cutoff}
  {\leandecl{Lib/Analytic/PrimeNumberTheorem/Tauberian}{319}{eventually\_\allowbreak abs\_\allowbreak le\_\allowbreak inv\_\allowbreak pow\_\allowbreak of\_\allowbreak transform}}
\end{description}

\subsection*{\paperlanguage{補助数体と素数冪和}{Auxiliary field and prime-power sums}}
\begin{description}
\leanentry{Parameters.constantBound}
  {固定定数と許容データの定義}
  {Fixed constant and definition of admissible data}
  {\leandecl{New/Reduction/Data/AdmissibleData}{50}{Parameters.constantBound}\enspace \leandecl{New/Reduction/Data/AdmissibleData}{77}{Data}}
\leanentry{Parameters.IsSparse}
  {高い重さを持つ素点の疎性}
  {Sparsity of the places of higher weight}
  {\leandecl{New/Reduction/Data/AdmissibleData}{58}{Parameters.IsSparse}}
\leanentry{exists_three_term_free}
  {桁の構成による疎集合}
  {Sparse sets from the digit construction}
  {\leandecl{New/Reduction/Data/Existence}{80}{exists\_\allowbreak three\_\allowbreak term\_\allowbreak free}}
\leanentry{Parameters.exists_isSparse}
  {位数と大きさを制御した疎集合の存在}
  {Existence of sparse sets with the required bounds}
  {\leandecl{New/Reduction/Data/Existence}{244}{Parameters.exists\_\allowbreak isSparse}\enspace \leandecl{New/Reduction/Data/Existence}{488}{Parameters.exists\_\allowbreak isSparse\_\allowbreak of\_\allowbreak ranks}}
\leanentry{eventually_exists_data}
  {許容データとレベル素数の選択}
  {Choice of admissible data and level primes}
  {\leandecl{New/Reduction/Data/Existence}{898}{eventually\_\allowbreak exists\_\allowbreak data}\enspace \leandecl{New/Reduction/Data/Existence}{820}{Parameters.IsLargeScale.exists\_\allowbreak isAdmissibleData}}
\leanentry{ramifiedSourceForms}
  {分岐条件を課した始域}
  {Source forms subject to the ramification condition}
  {\leandecl{New/UpperBound/Size/ShortColumns}{741}{ramifiedSourceForms}}
\leanentry{primeSum_le}
  {重み付き素数和と共通主項の評価}
  {Bounds for the weighted prime sum and the common main term}
  {\leandecl{New/Reduction/PrimeSums}{206}{primeSum\_\allowbreak le}\enspace \leandecl{New/Reduction/PrimeSums}{253}{jetSize\_\allowbreak sub\_\allowbreak primeSum\_\allowbreak mem\_\allowbreak Icc}}
\leanentry{totalRemainder_eq}
  {二つの打切り点を使った総剰余の恒等式}
  {Identity for the total remainder at the two cutoffs}
  {\leandecl{New/Reduction/PrimeSums}{177}{totalRemainder\_\allowbreak eq}}
\end{description}

\subsection*{\paperlanguage{捻りつき格子と積公式}{Twisted lattices and the product formula}}
\begin{description}
\leanentry{gramForm_mem_embValRing}
  {共役する有限素点での内積の整性}
  {Integrality of the pairing at conjugate finite places}
  {\leandecl{Lib/Adelic/AutomorphicForms/PAdicGauge}{706}{gramForm\_\allowbreak mem\_\allowbreak embValRing}}
\leanentry{Data.targetLattice}
  {目標空間の捻りつき格子}
  {Twisted lattice of target forms}
  {\leandecl{New/Reduction/IntegralTargetForms}{64}{Data.targetLattice}}
\leanentry{IndexLe}
  {対数指数と二乗余体積、および Hadamard の評価}
  {Index, squared covolume, and Hadamard's bound}
  {\leandecl{New/Reduction/Lattices/TwistedLattice}{159}{IndexLe}\enspace \leandecl{New/Reduction/Lattices/TwistedLattice}{275}{CovolumeLe}\enspace \leandecl{New/Reduction/Lattices/TwistedLattice}{317}{covolumeLe\_\allowbreak of\_\allowbreak normSqAt\_\allowbreak le}}
\leanentry{isAlgebraic_of_norm_complexToCp_le}
  {有限素点での一様有界性と代数性}
  {Algebraicity from a uniform finite-place bound}
  {\leandecl{Lib/General/NumberTheory/PadicLimitsValuation}{50}{isAlgebraic\_\allowbreak of\_\allowbreak norm\_\allowbreak complexToCp\_\allowbreak le}}
\leanentry{HeightLe.nonneg_of_one}
  {高さについての積公式}
  {The product formula for heights}
  {\leandecl{New/Reduction/Lattices/Heights}{191}{HeightLe.nonneg\_\allowbreak of\_\allowbreak one}\enspace \leandecl{New/Reduction/Lattices/Heights}{225}{HeightLe.neg\_\allowbreak le\_\allowbreak of\_\allowbreak norm\_\allowbreak map\_\allowbreak le}}
\leanentry{covolume_add_index_nonneg}
  {積公式による行列式と指数の比較}
  {Comparison of determinant and index by the product formula}
  {\leandecl{New/Reduction/Lattices/TwistedLattice}{477}{covolume\_\allowbreak add\_\allowbreak index\_\allowbreak nonneg}}
\end{description}

\subsection*{\paperlanguage{局所整除性と上界}{Local divisibility and the upper bound}}
\begin{description}
\leanentry{exists_isAdmissibleCombination_seesawDatum}
  {合成持ち上げの係数と不変性}
  {Coefficients and invariance in the composite lift}
  {\leandecl{New/UpperBound/Seesaw/SeesawIdentity}{592}{exists\_\allowbreak isAdmissibleCombination\_\allowbreak seesawDatum}\enspace \leandecl{New/UpperBound/Seesaw/SeesawIdentity}{509}{isInvariant\_\allowbreak seesawDatum}}
\leanentry{exists_hasNormalForm}
  {テータ積を整な因子と指数因子に分ける標準形}
  {Normal form for the product of theta series}
  {\leandecl{New/Derivatives/FormalGroup/ThetaProducts}{443}{exists\_\allowbreak hasNormalForm}}
\leanentry{exists_norm_schurJet_le_norm_jet_minorProd}
  {シュール微分から小行列式に沿う微分への帰着}
  {Reduction of Schur derivatives to derivatives along products of minors}
  {\leandecl{New/Derivatives/Divisibility/Count}{800}{exists\_\allowbreak norm\_\allowbreak schurJet\_\allowbreak le\_\allowbreak norm\_\allowbreak jet\_\allowbreak minorProd}}
\leanentry{exists_newDirection_dvd}
  {整除性を高める微分方向の選択}
  {Choice of differentiation directions for divisibility}
  {\leandecl{New/Derivatives/Divisibility/NullFrames}{361}{exists\_\allowbreak newDirection\_\allowbreak dvd}}
\leanentry{norm_jet_minorProd_le_of_charges}
  {方向のスケーリングを含む付値評価}
  {Valuation bound including the scaling of the directions}
  {\leandecl{New/Derivatives/Divisibility/Count}{662}{norm\_\allowbreak jet\_\allowbreak minorProd\_\allowbreak le\_\allowbreak of\_\allowbreak charges}}
\leanentry{le_divisibilityBound}
  {一段の持ち上げにおける付値の利得}
  {Valuation gain in one lift}
  {\leandecl{New/UpperBound/Divisibility/JetExponent}{504}{le\_\allowbreak divisibilityBound}}
\leanentry{exists_isFullMark}
  {例外素数で共役素点の片方を選ぶ補題}
  {Choice of one place from each conjugate pair at an exceptional prime}
  {\leandecl{New/UpperBound/Divisibility/MarkedPrimes}{230}{exists\_\allowbreak isFullMark}}
\leanentry{gaugeLe_scaledHeckeOperator_pow}
  {スケーリングした Hecke 作用素の冪による付値評価}
  {Valuation bound after powers of scaled Hecke operators}
  {\leandecl{New/UpperBound/Divisibility/ValuationAtNonSplitPrimes}{984}{gaugeLe\_\allowbreak scaledHeckeOperator\_\allowbreak pow}}
\leanentry{sum_allowedLoss_mul_log_le}
  {各素数の損失を足し合わせた上界}
  {Bound for the total loss over the primes}
  {\leandecl{New/UpperBound/Size/Estimates}{342}{sum\_\allowbreak allowedLoss\_\allowbreak mul\_\allowbreak log\_\allowbreak le}}
\leanentry{sqrt_weylNormSq_jetMap_le}
  {一段のテータ核の無限素点でのノルム評価}
  {Archimedean norm bound for a single theta kernel}
  {\leandecl{New/UpperBound/Size/ArchimedeanSize}{226}{sqrt\_\allowbreak weylNormSq\_\allowbreak jetMap\_\allowbreak le}}
\leanentry{sum_rowSizeSum_le}
  {行の長さから得る階乗の寄与の総和}
  {Sum of the factorial contributions from the rows}
  {\leandecl{New/UpperBound/Size/ArchimedeanSize}{290}{sum\_\allowbreak rowSizeSum\_\allowbreak le}}
\leanentry{weylNormSq_restrictedChain_le}
  {制限した連鎖のノルムの上界}
  {Norm bound for the restricted chain}
  {\leandecl{New/UpperBound/Size/ArchimedeanSize}{359}{weylNormSq\_\allowbreak restrictedChain\_\allowbreak le}}
\leanentry{span_isShort_eq_top}
  {短い整形式による始域の生成}
  {Spanning the source by short integral forms}
  {\leandecl{New/UpperBound/Size/ShortColumns}{981}{span\_\allowbreak isShort\_\allowbreak eq\_\allowbreak top}}
\end{description}

\subsection*{\paperlanguage{二次モーメントと生成}{Second moments and spanning}}
\begin{description}
\leanentry{targetDim_eq_finrank_goodForms_spaceW}
  {定符号空間との多重度の一致}
  {Equality of multiplicities with the definite space}
  {\leandecl{New/Spanning/Multiplicities/EqualMultiplicities}{406}{targetDim\_\allowbreak eq\_\allowbreak finrank\_\allowbreak goodForms\_\allowbreak spaceW}}
\leanentry{finrank_principalForms_spaceW}
  {定符号空間の主合同レベルでの次元}
  {Dimension at principal level on the definite space}
  {\leandecl{New/Spanning/SecondMoment/SpanningTheorem}{69}{finrank\_\allowbreak principalForms\_\allowbreak spaceW}}
\leanentry{hasAmplitudeBlocks_scaledPoincare}
  {ポアンカレ級数の内積の行列表示}
  {Block expansion of the Poincaré-series pairing}
  {\leandecl{New/Spanning/SecondMoment/SpanningTheorem}{168}{hasAmplitudeBlocks\_\allowbreak scaledPoincare}}
\leanentry{besselTranslate_mul_left}
  {Bessel 部分群の作用と展開式の同変性}
  {Bessel subgroup action and equivariance of the expansion}
  {\leandecl{New/Spanning/Unfolding/Moments}{231}{besselTranslate\_\allowbreak mul\_\allowbreak left}\enspace \leandecl{New/Spanning/Unfolding/Moments}{382}{HasAmplitudeBlocks.adjoint\_\allowbreak comp\_\allowbreak rep\_\allowbreak mul\_\allowbreak left}}
\leanentry{linkSecondMoment}
  {同じフレームの対応に属する項の二次モーメント}
  {Second moment for terms with the same frame correspondence}
  {\leandecl{New/Spanning/RemainingTerms/Size}{80}{linkSecondMoment}}
\leanentry{filling}
  {Bessel 系とレベルに対する許容条件}
  {Admissibility of the Bessel data and the level}
  {\leandecl{New/Spanning/Unfolding/FillingAdmissible}{743}{instance D.filling.IsAdmissible}}
\leanentry{sqrt_secondMoment_compl_besselS_le}
  {Bessel 部分群外での二次モーメントの分解}
  {Decomposition of the second moment away from the Bessel subgroup}
  {\leandecl{New/Spanning/Unfolding/SquareExpansion}{559}{sqrt\_\allowbreak secondMoment\_\allowbreak compl\_\allowbreak besselS\_\allowbreak le}}
\leanentry{secondMoment_remainingTerms_le}
  {主項へ入らない項の二次モーメント}
  {Bound for the terms outside the main contribution}
  {\leandecl{New/Spanning/RemainingTerms/Sum}{682}{secondMoment\_\allowbreak remainingTerms\_\allowbreak le}}
\leanentry{linkSecondMoment_le_of_le_meetNorm}
  {共通部分の格子による二次モーメント評価}
  {Second-moment bound from the lattice of the common intersection}
  {\leandecl{New/Spanning/RemainingTerms/Sum}{759}{linkSecondMoment\_\allowbreak le\_\allowbreak of\_\allowbreak le\_\allowbreak meetNorm}}
\leanentry{linkSecondMoment_isOpenElement_le}
  {開な還元に対する局所密度の評価}
  {Local density bound on the open reductions}
  {\leandecl{New/Spanning/MainTerm/EnergyEstimate}{631}{linkSecondMoment\_\allowbreak isOpenElement\_\allowbreak le}}
\leanentry{secondMoment_compl_besselS_le}
  {Bessel 部分群の近傍での誤差}
  {Error near the Bessel subgroup}
  {\leandecl{New/Spanning/Unfolding/SquareExpansion}{644}{secondMoment\_\allowbreak compl\_\allowbreak besselS\_\allowbreak le}}
\leanentry{bad_types_rare}
  {射影で取り除かれる型の寄与}
  {Bound for the types discarded by the projection}
  {\leandecl{New/Spanning/SecondMoment/SpanningTheorem}{244}{bad\_\allowbreak types\_\allowbreak rare}}
\leanentry{trace_besselTranslate_eq_mul_reducedChar}
  {Bessel の跡と共役平均の評価}
  {Bessel traces and a bound on their conjugacy average}
  {\leandecl{New/Spanning/Unfolding/Moments}{751}{trace\_\allowbreak besselTranslate\_\allowbreak eq\_\allowbreak mul\_\allowbreak reducedChar}\enspace \leandecl{New/Spanning/Unfolding/Moments}{991}{sum\_\allowbreak norm\_\allowbreak reducedCharMean\_\allowbreak sq\_\allowbreak le}}
\leanentry{isNearIsometry_piTensorElem_comp}
  {射影ポアンカレ写像の近似等長性}
  {Near-isometry of the projected Poincaré map}
  {\leandecl{New/Spanning/Unfolding/Moments}{1230}{HasAmplitudeBlocks.isNearIsometry\_\allowbreak piTensorElem\_\allowbreak comp}}
\leanentry{secondMoment_fixer_univ_le}
  {すべての選択素点を固定した場合の二次モーメント}
  {Second moment with fixing conditions at every selected place}
  {\leandecl{New/Spanning/Unfolding/MomentBounds}{202}{secondMoment\_\allowbreak fixer\_\allowbreak univ\_\allowbreak le}}
\leanentry{card_mul_sum_det_eq_zero_le}
  {有限 Fourier 展開による退化小行列式の評価}
  {Bound for singular minors by finite Fourier expansion}
  {\leandecl{New/Spanning/SecondMoment/PositiveType}{153}{card\_\allowbreak mul\_\allowbreak sum\_\allowbreak det\_\allowbreak eq\_\allowbreak zero\_\allowbreak le}}
\leanentry{sum_not_isOpenElement_le}
  {開でない還元の寄与}
  {Contribution from reductions outside the open set}
  {\leandecl{New/Spanning/Unfolding/MomentBounds}{399}{sum\_\allowbreak not\_\allowbreak isOpenElement\_\allowbreak le}}
\leanentry{sum_secondMoment_le}
  {全群にわたる二次モーメントの評価}
  {Second moment over the full group}
  {\leandecl{New/Spanning/Unfolding/MomentBounds}{441}{sum\_\allowbreak secondMoment\_\allowbreak le}}
\leanentry{isSpread_of_second_moment}
  {二次モーメントから像の階数へ}
  {Rank of images from the second moment}
  {\leandecl{New/Spanning/SecondMoment/FewTranslates}{320}{isSpread\_\allowbreak of\_\allowbreak second\_\allowbreak moment}}
\leanentry{IsSpread.exists_few_le_finrank_iSup_range}
  {少数の平行移動で目標空間を生成}
  {Spanning the target with few translates}
  {\leandecl{New/Spanning/SecondMoment/FewTranslates}{479}{IsSpread.exists\_\allowbreak few\_\allowbreak le\_\allowbreak finrank\_\allowbreak iSup\_\allowbreak range}}
\leanentry{finrank_ker_one_sub_scalingGenerator_le}
  {中心作用の固定空間の次元}
  {Dimension of the fixed space of the central action}
  {\leandecl{New/UpperBound/Size/ShortColumns}{857}{finrank\_\allowbreak ker\_\allowbreak one\_\allowbreak sub\_\allowbreak scalingGenerator\_\allowbreak le}}
\leanentry{enough_forms}
  {分岐条件を課した始域からの生成}
  {Spanning with ramification imposed on the source}
  {\leandecl{New/Reduction/ShortAndEnough}{448}{enough\_\allowbreak forms}}
\end{description}

\subsection*{\paperlanguage{行列式の比と結論}{Determinant ratios and conclusion}}
\begin{description}
\leanentry{clearingBound}
  {射影の分母に対する局所補正を含む上界}
  {Bounds for denominators of projections with the local corrections}
  {\leandecl{New/Congruences/FirstOrder/TraceRelations}{128}{clearingBound}}
\leanentry{heckeValueProd}
  {交換に現れる Hecke 値の積}
  {Product of Hecke values in the exchange identities}
  {\leandecl{New/OneLift/ComparisonOperator/ThreeExchanges}{181}{heckeValueProd}}
\leanentry{integral_forms_long}
  {二乗余体積と指数を合わせた下界}
  {Lower bound combining squared covolume and index}
  {\leandecl{New/Determinant/Twist/LongForms}{214}{integral\_\allowbreak forms\_\allowbreak long}}
\leanentry{star_inequality}
  {主項を相殺して得られる比較不等式}
  {Comparison inequality after cancellation of the main term}
  {\leandecl{New/Reduction/MainInequality}{87}{star\_\allowbreak inequality}}
\leanentry{eventually_primeRemainder_le_rpow}
  {素数和の剰余の片側上界}
  {One-sided bound for the prime-sum remainder}
  {\leandecl{New/Reduction/OneSidedBound}{200}{eventually\_\allowbreak primeRemainder\_\allowbreak le\_\allowbreak rpow}}
\leanentry{grh}
  {右半平面の非消滅と臨界線上の零点}
  {Nonvanishing in the right half-plane and zeros on the critical line}
  {\leandecl{New/Reduction/OneSidedBound}{239}{grh}\enspace \leandecl{New/Main}{28}{zeros\_\allowbreak on\_\allowbreak line}\enspace \leandecl{New/Main}{47}{zeta\_\allowbreak zeros\_\allowbreak on\_\allowbreak line}}
\end{description}

\end{document}
